\documentclass[11pt,a4paper]{article}
\usepackage[T1]{fontenc}
\usepackage{lmodern}
\hyphenation{OpenAI}
\usepackage[margin=25mm,headheight=14pt]{geometry}
\usepackage{amsmath,amssymb,amsthm,mathtools,mathrsfs}
\usepackage{booktabs,array,tabularx,longtable}
\usepackage{enumitem}
\usepackage[expansion=false]{microtype}
\usepackage{xcolor}
\definecolor{linkblue}{HTML}{173F65}
\usepackage[colorlinks=true,linkcolor=linkblue,citecolor=linkblue,urlcolor=linkblue,hypertexnames=false]{hyperref}
\usepackage{xurl}
\usepackage{fancyhdr}
\pagestyle{fancy}
\fancyhf{}
\fancyhead[L]{\small Unitary Weil $L$-functions}
\fancyhead[R]{\small A seven-eighths half-plane}
\fancyfoot[C]{\thepage}
\renewcommand{\headrulewidth}{0pt}
\setlength{\emergencystretch}{3em}
\setlength{\parskip}{2pt}
\setlist{topsep=4pt,itemsep=2pt,parsep=0pt}
\allowdisplaybreaks[2]
\numberwithin{equation}{section}
\newtheorem{theorem}{Theorem}[section]
\newtheorem{proposition}[theorem]{Proposition}
\newtheorem{lemma}[theorem]{Lemma}
\newtheorem{corollary}[theorem]{Corollary}
\theoremstyle{definition}
\newtheorem{definition}[theorem]{Definition}
\theoremstyle{remark}
\newtheorem{remark}[theorem]{Remark}
\newcommand{\eps}{\varepsilon}
\newcommand{\RR}{\mathbb R}
\newcommand{\CC}{\mathbb C}
\newcommand{\NN}{\mathbb N}
\newcommand{\ZZ}{\mathbb Z}
\newcommand{\QQ}{\mathbb Q}
\newcommand{\Rea}{\operatorname{Re}}
\newcommand{\Ima}{\operatorname{Im}}
\newcommand{\Norm}{\operatorname{N}}
\newcommand{\rad}{\operatorname{rad}}
\newcommand{\vol}{\operatorname{vol}}
\newcommand{\Res}{\operatorname{Res}}
\newcommand{\cF}{\mathcal F}
\newcommand{\cB}{\mathcal B}
\newcommand{\Qsrc}{\textup{Q}}
\newcommand{\Psrc}{\textup{P}}
\newcommand{\ind}{\mathbf 1}
\hypersetup{pdftitle={A seven-eighths zero-free half-plane for unitary Weil L-functions},pdfsubject={Hecke analytic core, unitary Weil representations, effective Chebotarev, and proof dependencies}}
\title{\vspace{-1em}A seven-eighths zero-free half-plane\\for unitary Weil $L$-functions}
\author{Boyang Hu}
\date{13 October 2026}
\begin{document}
\maketitle
\thispagestyle{plain}
\begin{abstract}
We prove that for every number field $F$ and every finite-dimensional
continuous unitary representation $r$ of the global Weil group $W_F$,
the $L$-function $L_F(s,r)$ has no zeros in $\Rea s>7/8$, and that its
only poles in this half-plane are at $s=1+it$, with order equal to the
multiplicity of $\mathcal N_F^{it}$ in $r$. The analytic core is the
rank-one case: $L_F(s,\psi)\ne0$ for $\Rea s>7/8$ for every unitary
Hecke character $\psi$ of every number field, including characters of
infinite order. This extends the cubic-theta method of the OpenAI
quasi-Riemann preprint, which treats finite-order characters over
$\QQ(\sqrt{-3})$, by transferring its estimates to number fields
containing $\mu_{12}$ and to families of fixed infinity type, followed
by an abelian descent. Brauer induction and Frobenius reciprocity then
give the result in every dimension. Consequences include the same
half-plane for all Artin $L$-functions and their unitary Hecke twists,
an effective Chebotarev theorem, cuspidal weight-one newforms, and CM
motives after the appropriate weight shift.
\end{abstract}

\noindent\textit{2020 Mathematics Subject Classification.} Primary 11M41;
secondary 11R39, 11R42, 11R44, 11F37, 11L40, 11G15.

\section{Introduction}\label{sec:introduction}

Let $F$ be a number field and let $W_F$ be its classical global Weil
group.  We use the norm character $\mathcal N_F$ normalized by
$\mathcal N_F(\operatorname{Frob}_{\mathfrak p})=\Norm\mathfrak p$,
and put $\nu_{F,t}=\mathcal N_F^{it}$ for $t\in\RR$.  Thus
$L_F(s,\nu_{F,t})=\zeta_F(s-it)$.  For a continuous unitary
representation $r:W_F\longrightarrow U(d)$ on $V_r$, write
\[
 L_F(s,r)=\prod_{\mathfrak p}
 \det\!\left(1-r(\operatorname{Frob}_{\mathfrak p})
     (\Norm\mathfrak p)^{-s}\bigm|V_r^{I_{\mathfrak p}}\right)^{-1}
 \qquad(\Rea s>1).
\]
We use the meromorphic continuation of this finite-place Euler product.  At a ramified prime the
Frobenius acts on the inertia invariants, so the definition is independent
of its lift.  We use the classical Weil-group formalism of
\cite[Chapters~1, 2 and~20]{Langlands}; finite-dimensional unitary
representations are automatically semisimple.

\begin{theorem}[Unitary Weil form]\label{thm:weil-main}
For every number field $F$ and every finite-dimensional continuous
unitary representation $r$ of $W_F$, the meromorphic function
$L_F(s,r)$ has no zeros in $\Rea s>7/8$.  Its only poles in this
half-plane are prescribed by the norm-character constituents:
\begin{equation}\label{eq:weil-main-divisor}
 \operatorname{ord}_{s=1+it}L_F(s,r)
   =-m_t(r),\qquad
 m_t(r)=\dim_{\CC}\operatorname{Hom}_{W_F}(\nu_{F,t},r).
\end{equation}
Here positive order denotes a zero and negative order a pole.
Consequently, if $r$ contains no $\nu_{F,t}$, its $L$-function is
holomorphic and nowhere zero in this half-plane.  The statement and
the pole orders are unchanged on deleting any fixed finite set of
finite Euler factors.  The boundary is independent of $d$, the field,
the conductor, and the infinity type.
\end{theorem}

Sections~\ref{sec:field}--\ref{sec:closure} prove the following
one-dimensional result. Section~\ref{sec:weil} then deduces
Theorem~\ref{thm:weil-main} by classical induction theory.

\begin{theorem}[Hecke analytic core]\label{thm:main}
Let $F$ be a number field and let
$\psi:F^\times\backslash\mathbb A_F^\times\longrightarrow S^1$
be a continuous unitary Hecke character. The function $L_F(s,\psi)$
is holomorphic and nonzero in $\Rea s>7/8$ unless
$\psi=\mathcal N_F^{it}$ for some $t\in\RR$; in that case it is
nonzero there and its only singularity is a simple pole at $s=1+it$.
The same statement holds after deletion of any fixed finite set of
Euler factors.
\end{theorem}

In rank one the finite Euler product is
$\prod_{\mathfrak p}(1-\psi(\mathfrak p)(\Norm\mathfrak p)^{-s})^{-1}$,
with the usual convention at ramified primes.  Its meromorphic
continuation and pole classification are the standard Hecke theory
\cite{Tate}.

The principal analytic sources are the 30 September 2026 preprint
\cite{Q}, denoted $\Qsrc$, and the 4 October 2026 preprint
\cite{P}, denoted $\Psrc$. The theorem in \cite[Theorem~1.1]{Q} treats finite-order Hecke
characters over $\mathbb Q(\sqrt{-3})$ and Dirichlet characters, with
boundary $7/8$; \cite[Theorem~1.2]{P} treats finite-order Hecke characters
over cyclotomic fields containing $\mu_{12}$, with boundary $1-10^{-6}$.

Independently, Cureton and Thorner \cite{CT} followed the strategy of
$\Psrc$ over an arbitrary number field $K$ and proved that $\zeta_K(s)$,
and hence every finite-order Hecke $L$-function over $K$, has no zeros in
$\Rea s>1-1/6000$; their proofs were generated and verified in Lean.
Theorem~\ref{thm:main} implies their nonvanishing statement, but the
present argument has not been formally verified.

References to $\Qsrc$ and $\Psrc$ use the dated editions specified in the
bibliography; all their numerical locators refer to those editions.

\paragraph{Earlier work.}
Patterson's cubic theta series \cite{Patterson} and the metaplectic
theta construction of Kazhdan--Patterson \cite{KP,KPcorr} underlie
the reflection.  The unconditional Eisenstein-field cusp formulas used in
\cite[Section~5]{Q} are recorded in
\cite[Section~5.1 and Appendix~A]{DR}.  For the general-field
automorphic normalization we use \cite[Sections~3, 5--7 and Appendix~A]{P},
with the corrected Whittaker normalization explained below.
The marked reflection and physical probe estimates are taken from
\cite[Sections~5--16]{Q}; the canonical inverse and paired fourth-moment
arguments are those of \cite[Lemmas~17.1--17.2 and~18.1]{Q}.
The sixth-power amplification is \cite[Lemma~17.6]{Q}, and the witness
comparison follows \cite[Proposition~19.2]{Q}.
We retain their finite-seminorm, polynomial-height, strict-width,
and common-mask assertions, and use the signal-to-zero-free principle
of \cite[Proposition~2.1]{Q}.

\paragraph{Field and infinity type.}
Sections~\ref{sec:transfer} and~\ref{sec:canonical} establish the
ideal-norm and fixed-infinity-type versions of the source constructions.
The unit quotient, complete marked expansion, and recursive norm
identities are the three interfaces summarized in
Appendix~\ref{app:transfer-map}. Sections~\ref{sec:rows}
and~\ref{sec:closure} then make the exponent comparison and the
common-family closure.

\paragraph{Weil representations and geometric applications.}
Brauer's induction method \cite{Brauer}, in the Weil-group form of
\cite[Lemma~2.2]{Langlands}, supplies a virtual decomposition into
characters induced from finite extensions.  Lemma~\ref{lem:unitary-brauer}
verifies that the inducing characters can be taken unitary.  Frobenius
reciprocity then accounts exactly for the apparent zeros that negative
Brauer coefficients could produce at Hecke poles.  This gives
\eqref{eq:weil-main-divisor}; it is the same classical pole-counting
mechanism used on $\Rea s\geq1$ in
\cite[Proposition~11]{Johansson}, now applied with the Hecke half-plane
of Theorem~\ref{thm:main}.  

The applications include Artin representations and their unitary Hecke
twists, the effective Chebotarev theorem of
Corollary~\ref{cor:chebotarev}, weight-one newforms through
Deligne--Serre \cite[Theorems~4.1 and~4.6]{DS}, and potentially abelian
geometric realizations through \cite[Corollary~10 and Lemma~14]{Johansson}.
For pure weight $w$, the change from unitary to arithmetic normalization
moves the boundary to $w/2+7/8$; the CM realization used here is
\cite[Corollaries~7.6, 7.9 and Proposition~7.12]{Fargues}.

\subsection{Proof outline}\label{subsec:proof-outline}

The argument compares two estimates for one auxiliary function. Its
principal term contains the reciprocal of the target $L$-function;
controlling the difference extends that reciprocal to a larger
half-plane.

Fix $K=F(\zeta_{12})$ and a unitary character $\omega$ of $K$.
The family
$\cF_\omega=\{\theta\omega^e\Norm_K^{it}:e\in\{-1,0,1\},\
\theta\text{ finite order},\ t\in\RR\}$ contains every character
that appears in the estimates. Let $B$ be the supremum of the real
parts of their nontrivial zeros, with $B\ge1/2$, as in
\eqref{eq:common-boundary}. The Euler product gives $B\le1$.
Assume $B>7/8$ and fix a target $\psi$ in this family.

The cubic theta series is summed against an additive test and finitely
many smooth prime sums. This defines the auxiliary function
$I_{\psi,\mathrm{mod}}$ in \eqref{eq:full-compensation}. The
compensation at each prime is chosen so that its main prime factor
contains only the finite-order row character; see
\eqref{eq:can-pairing} and \eqref{eq:one-compensation}. Thus the
prime estimates used for an infinite-order target involve
finite-order Hecke characters in the same family. Poisson summation
and Mellin inversion express the auxiliary function as a sum over
sixth-power-free rows. The principal row has poles at $w=1$ and
$z=1/6$. After taking these residues and normalizing, its contribution is
\[
 f_\psi(Z)=\frac{1}{2\pi i}\int_{(2)}
 Z^{C(s)}e^{(s-5/6)^2}\frac{H_\psi^S(s)}{L_K^S(s,\psi)}\,ds,
 \qquad C(s)=s-\frac{11}{16}.
\]
The explicit Euler factor in Lemma~\ref{lem:principal-nonzero}
makes $H_\psi^S$ holomorphic and nonzero on $\Rea s>2/3$.

There are two ways to estimate the normalized auxiliary function
$J_\psi$. On the theta side, the complete marked reflection and the
additive Gram estimate give
$|J_\psi(Z)|\ll_{\psi,\eps}Z^{3/16+\eps}$.
On the Euler-product side, the zero detector associates exceptional
rows with large values of inverse or ordinary character sums.
The two canonical moments bound the number of these rows. With the
actual fourth-moment parameter $\kappa=2B-1$, the comparison in
Section~\ref{sec:rows} and its positive certificate leave a fixed
saving. The resulting estimates, for the same $J_\psi$, are
\[
 \begin{aligned}
 |J_\psi(Z)|&\ll_{\psi,\eps}Z^{C(7/8)+\eps},\\
 |J_\psi(Z)-f_\psi(Z)|&\ll_\psi Z^{C(B)-\eta_0},
 \end{aligned}
\]
where $\eta_0>0$ is common to the later finite-order twists of the
fixed infinity type. Consequently the Mellin transform of $f_\psi$
is holomorphic on $\Rea s>\max\{7/8,B-\eta_0\}$. It equals
$e^{(s-5/6)^2}H_\psi^S(s)/L_K^S(s,\psi)$ there. The nonvanishing
numerator excludes all zeros to the right of a line strictly smaller
than $B$, a contradiction. Abelian descent proves the Hecke theorem
over $F$; unitary Brauer induction and Frobenius reciprocity give the
Weil theorem and its pole formula.

The main geometric change occurs before these estimates. Over $K$,
associates are organized on the compact quotient
$T_0/\Gamma$, where $T_0=K_S^1$ and
$\Gamma=(\mathcal O_{K,S}^{\times})^6$. Weighted integration on this
quotient retains the continuous unit phases. Balancing the generators
keeps the additive shapes in a fixed compact set, so the two Poisson
transforms have the original norm scales; the additional kernel
logarithms are absorbed through finite-depth propagation.

Relative to the source framework, the proof has five components:
\begin{enumerate}[label=\textup{(\alph*)},leftmargin=*,itemsep=2pt]
\item Ideal and unit-quotient formulas over every $K\supset\mu_{12}$,
with complete local coefficients and both Poisson zero orbits.
\item Fixed-infinity-type reflection and paired prime coefficients,
which supply the moments for infinite-order targets.
\item One prime-slot system for the entire moderate-row range and the
alternative positive certificate \eqref{row:square-certificate}.
\item A common positive saving and the single Mellin closure of
Lemma~\ref{lem:mellin-closure}.
\item The Weil deduction, its exact pole orders, and the arithmetic
consequences of Section~\ref{sec:weil}.
\end{enumerate}
Appendix~\ref{app:transfer-map} matches each source step with the
field-sensitive input, its replacement here, and the norm calculation
that remains unchanged. The parameter choices are listed in order in
Section~\ref{subsec:parameter-order}.

\subsection{Terminology}\label{subsec:terminology}

All lengths below are logarithmic: a variable of length $a$ has norm
of order $Z^a$. The following terms describe the sums and the finite
recursions used in Sections~\ref{sec:canonical}--\ref{sec:rows}.
\begin{description}[style=nextline,leftmargin=0pt,labelwidth=0pt,labelsep=0pt,font=\normalfont\itshape,
                   itemsep=3pt]
\item[Auxiliary function (physical probe).]
The theta--additive sum $I_{\psi;D}(X,Y,Z)$ in
\eqref{eq:physical-I}, with its compensated and normalized form
$J_\psi$ in \eqref{eq:signal}. Its principal Mellin contribution is
$f_\psi$.
\item[Prime slot and mark.]
A slot is one factor $Q_{u,i}$: a smoothly weighted sum over primes
in a prescribed annulus. In a completed theta sum its corresponding
mark is the divisibility condition $\ind_{p_i\mid A}$ on the whole
index $A=cn^3$.
\item[Ordinary sum (plain) and inverse sum.]
The ordinary sum has coefficients $\lambda_u(\mathfrak a)$; the
inverse sum has coefficients $\mu_K(\mathfrak a)\lambda_u(\mathfrak a)$.
Both are normalized by the square root of their length scale.
\item[Capacity and strict width.]
A capacity is an allowed total prime-slot length, specified by the
moment inequalities. A strict width is the positive difference between
the two sides of such an inequality. A width decrement is the fixed
amount reserved when passing to a recursive child.
\item[Terminal and canonical energy.]
A terminal is a range estimated directly, without another recursive
call. The canonical energy \eqref{eq:can-energy} is the row-and-label
second moment whose coefficient form is recovered after two Poisson
transforms.
\item[Moving puncture and common mask.]
A puncture is a radical $\mathfrak r$ imposing
$\ind_{(\mathfrak a,\mathfrak r)=1}$. It is fixed in the current inner
sum but may vary with outer data; its logarithmic norm is charged to
the capacity. A common mask is the same coprimality restriction in all
factors of the current expression.
\item[Frozen outer ideals.]
These are ideals held fixed during an inner estimate and then summed
with their displayed weights. Uniformity in them means that their
allowed norm ranges, rather than their individual values, determine
the estimate.
\item[Coprime residual support.]
This is the support left after the displayed greatest-common-divisor
extractions and the original character zeros have been imposed.
The full M\"obius expansion is applied to this support before any
unrestricted coefficient extension.
\end{description}

\section{Field reduction}\label{sec:field}

\subsection{Abelian descent}

\begin{lemma}\label{lem:descent}
Let $K/F$ be a finite abelian extension and $\sigma_0>0$.  If
$L_K(s,\psi\circ\Norm_{K/F})$ is zero-free in $\Rea s>\sigma_0$, then
so is $L_F(s,\psi)$.
\end{lemma}
\begin{proof}
Put $G=\operatorname{Gal}(K/F)$ and view $\widehat G$ as finite-order
Hecke characters of $F$.  After omitting the ramified primes and the
conductors from one fixed finite set $S$, class field theory gives
\begin{equation}\label{eq:descent-product}
 L_K^S(s,\psi\circ\Norm_{K/F})
   =\prod_{\chi\in\widehat G}L_F^S(s,\psi\chi).
\end{equation}
For completeness, at an unramified prime whose Frobenius has order $f$,
the values $\chi(\operatorname{Frob}_{\mathfrak p})$ run through the
$f$th roots of unity, each $|G|/f$ times.  This proves the Euler identity
on the absolute half-plane, and continuation proves
\eqref{eq:descent-product} everywhere; see also\ \cite{Milne}.

If every factor is holomorphic at $s_0$, a zero of one factor would give a
zero on the left.  Otherwise the Hecke pole classification gives a unique
$\chi_0$ with $\psi\chi_0=\Norm_F^{it}$ and $s_0=1+it$.  Its pullback
is $\Norm_K^{it}$.  Both the left side and the $\chi_0$ factor have a
simple pole, so meromorphic orders give
\[
 -1=-1+\sum_{\chi\ne\chi_0}
          \operatorname{ord}_{s_0}L_F^S(s,\psi\chi).
\]
Every summand is nonnegative and is therefore zero.  Finally, a deleted
Euler polynomial $1-\xi(\mathfrak p)(\Norm\mathfrak p)^{-s}$ is nonzero
for $\Rea s>0$ when $\xi$ is unitary.  This proves the lemma, including
possible norm-character poles.
\end{proof}

The extension $K=F(\zeta_{12})$ is abelian over $F$, has degree at most
four, is totally imaginary, and contains $\mu_{12}$.    We henceforth work over such a fixed $K$.

\subsection{The fixed-type family}

Fix a unitary Hecke character $\omega$ of $K$ and set
\begin{equation}\label{eq:family}
 \cF_\omega=
 \{\theta\omega^e\Norm_K^{it}: e\in\{-1,0,1\},\
       \theta\text{ finite order},\ t\in\RR\}.
\end{equation}
Ideal values are normalized by
$\Norm_K^{it}(\mathfrak a)=(\Norm\mathfrak a)^{it}$.
The family contains all finite-order characters and is closed under finite
twists, conjugation, and vertical translation.  Define its common boundary by
\begin{equation}\label{eq:common-boundary}
 B=\sup\bigl(\{1/2\}\cup
  \{\Rea\rho:\eta\in\cF_\omega,\ L_K(\rho,\eta)=0,
                         \ 1/2\leq\Rea\rho\leq1\}\bigr).
\end{equation}
The Euler product gives $B\leq1$.  The parameter $t$ does not affect
this supremum.  Under the eventual contradiction $B>7/8$, put
\begin{equation}\label{eq:dynamic}
 \Delta=B-\frac78\in(0,1/8],\qquad
 \kappa=2B-1=\frac34+2\Delta\in(3/4,1],\qquad
 \alpha=\frac56.
\end{equation}

All arithmetic sums use ideal norms $q_{\mathfrak a}=\Norm\mathfrak a$.
When a roman letter denotes a generator, its finite class and unit sector
are specified by the allowed-generator convention below.  
Every power-residue character retains its zero extension on nonunits,
including powers whose exponent is zero modulo six.

The real losses, strict-width decrements, and finite analytic slot system
will be fixed for the three types $0,\omega_\infty,\bar\omega_\infty$
before a later finite-order twist is chosen.  Constants, seminorm orders,
height degrees, and lower thresholds may depend on that later character.
Section~\ref{sec:closure} records the complete order of choices.

\section{Arithmetic and reflection}
\label{sec:transfer}

Throughout this section $K\supset\mu_{12}$ is fixed and
$[K:\mathbb Q]=2r$.  At a finite place we use
$|\varpi_v|_v=(\operatorname N v)^{-1}$, and at a complex place we use
$|z|_v=|z|_{\rm ord}^2$.  Fix the standard trace additive character and
self-dual additive measures.  Choose $S$ to contain the complex places,
the primes over $6$, the conductors of the fixed characters, the
ramification of the additive character, and the fixed exceptional places
of the theta construction.  Enlarge $S$ by primes generating the ideal
class group.  Thus $R=\mathcal O_{K,S}$ is principal.  Write
\[
 K_S=\prod_{v\in S}K_v,\qquad
 |y|_S=\prod_{v\in S}|y_v|_v,\qquad
 T_0=K_S^1=\{y\in K_S^\times:|y|_S=1\},\qquad
 \Gamma=(R^\times)^6.
\]
An exterior ideal is an integral ideal of $R$.  Its norm is
$q_{\mathfrak a}=|R/\mathfrak a|$, extended multiplicatively to fractional
ideals.  The positive norm section $\ell_x$ has
finite components $1$ and complex components $x^{1/(2r)}$, so
$|\ell_x|_S=x$ and $\ell_x\ell_y=\ell_{xy}$.

\subsection{Generators, reciprocity and the Gauss corrections}

\begin{proposition}[Arithmetic calibration]\label{prop:arithmetic-calibration}
The following choices and identities hold over $K$.
\begin{enumerate}[label=\textup{(\roman*)}]
\item The diagonal lattice $R\subset K_S$ is self-dual and has covolume
$1$.  For $a\in K^\times$, $q_a=|a|_S=q_{aR}$ and
$(aR)^\perp=a^{-1}R$.
\item Put $V=K_S^\times/K_S^{\times6}$ and
$H_S(x,y)=\prod_{v\in S}(x_v,y_v)_{6,v}$, with the reciprocity convention
of \cite[Section~3]{P}.  The image $E$ of $R^\times/(R^\times)^6$ is
Lagrangian, and $V=E\oplus B$ for an isotropic subgroup $B$.
Each exterior fractional ideal has an allowed generator $a$ with
$[a]\in B$, unique modulo $\Gamma$.  Such generators can be balanced so
that $a/\ell_{q_a}$ lies in a fixed compact subset of $T_0$.
\item For an exterior prime $\mathfrak p$, define the sextic residue
character $\chi_{\mathfrak p}$ by the power $(q_{\mathfrak p}-1)/6$.
Every character power is extended by zero on nonunits, including powers
divisible by $6$.  On allowed coprime generators,
\begin{equation}\label{eq:allowed-reciprocity}
 \chi_{\mathfrak a}(b)=\chi_{\mathfrak b}(a).
\end{equation}
For $u\in R\setminus\{0\}$, $\mathfrak a\mapsto\chi_{\mathfrak a}(u)$ is
induced away from its displayed zeros by a finite-order Hecke character,
whose primitive conductor satisfies
\begin{equation}\label{eq:kummer-conductor}
 \operatorname N\mathfrak f_u\ll_{K,S}
       \operatorname N\operatorname{rad}(uR)\le q_u.
\end{equation}
Two such inducing characters agree exactly when their parameters have
quotient in $K^{\times6}$.  Characters of $B$ extend to finite-order
Hecke characters unramified outside $S$.
\item If $\gamma_j(\mathfrak c)$ denotes the normalized Gauss sum of
$\chi_{\mathfrak c}^j$ for squarefree $\mathfrak c$, there are a unitary
Hecke character $\Lambda$ and a finite-order Hecke character $G$,
unramified outside $S$, for which
\begin{equation}\label{eq:gauss-calibration}
 \gamma_2(\mathfrak p)^3=-\Lambda(\mathfrak p),\qquad
 G(\mathfrak p)=\overline{\chi_{\mathfrak p}(4)}\gamma_3(\mathfrak p),
 \qquad \gamma_1(\mathfrak p)\gamma_2(\mathfrak p)
       =-\Lambda(\mathfrak p)G(\mathfrak p).
\end{equation}
The complex angular exponents of $\Lambda$ are $\pm1$, and it has no
radial infinity parameter.
\end{enumerate}
\end{proposition}

\begin{proof}
Apply adelic Poisson summation to
$\phi\otimes\prod_{v\notin S}\mathbf1_{\mathcal O_v}$.
The exterior factors are self-dual by the choice of $S$, giving the
lattice $R$ with scalar factor one on both sides.  Strong approximation and
additive adelic duality identify $K_S/R$ with the corresponding compact
adelic quotient and identify the annihilator with $R$.  Scaling and the
product formula prove (i); see also \cite[Lemma~3.1]{P} and \cite{Tate}.

For $\ell=2,3$, local Kummer theory and the $S$-unit theorem give
\[
 \dim_{\mathbb F_\ell}K_S^\times/K_S^{\times\ell}=2|S_f|+2r,
 \qquad
 \dim_{\mathbb F_\ell}R^\times/(R^\times)^\ell=|S_f|+r.
\]
Indeed a finite place contributes
$2+\mathbf1_{v\mid\ell}[K_v:\mathbb Q_\ell]$, and complex places
contribute zero.  An $S$-unit locally an $\ell$th power at $S$ defines a
Kummer extension unramified outside $S$ and split at $S$.  The trivial
$S$-class group makes this extension trivial, proving injection of the
unit quotient.  Hilbert reciprocity makes its image isotropic.
Since $-1=\zeta_{12}^6$, the diagonal Hilbert symbols are trivial, so the
pairings are alternating.  The dimension calculation gives a Lagrangian
unit image; complementary isotropic subspaces for $\ell=2,3$ combine to
give $B$.  Multiplying any generator by an $S$-unit puts its class in $B$;
uniqueness follows from $E\cap B=1$.  The norm-one idele class group is
compact, whence $T_0/R^\times$, and therefore $T_0/\Gamma$, is compact.
This proves the balancing assertion.

The exterior Hilbert symbols of two allowed generators and the trivial
$S$-pairing give \eqref{eq:allowed-reciprocity}.  If
$m=\epsilon b$ with $b$ allowed and $\epsilon\in R^\times$, the same
calculation gives, with the original zeros retained,
\begin{equation}\label{eq:unit-reciprocity}
 \chi_{\mathfrak a}(m)=\chi_{\mathfrak a}(\epsilon)
                         \chi_{\mathfrak b}(a),\qquad
 \chi_{\mathfrak a}(\epsilon)=H_S(\epsilon,a)^{-1}.
\end{equation}
Kummer theory identifies the inducing character with the action on
$u^{1/6}$.  At an exterior prime it is tame, with conductor exponent at
most one, and is unramified if $u$ is a unit.  At $S$ only finitely many
local power classes occur.  This proves \eqref{eq:kummer-conductor} and
the equality criterion.  A character of $B$ extends trivially across
$E$; the resulting character of $K_S^\times$, trivial on $R^\times$,
defines the asserted Hecke character on exterior ideals.  In particular,
reducing exterior valuations modulo six gives distinct sixth-power-free
row representatives modulo $\Gamma$; only the representative $1$
induces the principal character.

For completeness, the Gauss calibration is independent of any Galois
hypothesis on $K$.  Put $K_0=\mathbb Q(\sqrt{-3})$.  If
$\pi_+(\mathfrak q)$ is the generator congruent to $1$ modulo $3$, the
cubic Jacobi identity gives
$J_{q}(\rho,\rho)=-\pi_+(\mathfrak q)$ and
$\gamma_2^3=q^{-1/2}J_q(\rho,\rho)$.
For residue degree $f$, the classical Davenport--Hasse lifting formula,
in the normalization recorded in \cite[Theorem~1.1]{HK}, gives
\[
 J_{q^f}(\rho\circ N,\rho\circ N)
   =(-1)^{f-1}J_q(\rho,\rho)^f=-\pi_+(\mathfrak q)^f.
\]
Consequently $\Lambda=\Lambda_{K_0}\circ N_{K/K_0}$, where
$\Lambda_{K_0}(\mathfrak a)=\pi_+(\mathfrak a)/|\pi_+(\mathfrak a)|$,
has the first property in \eqref{eq:gauss-calibration}.  The local norm
at infinity is a product of coordinates or their conjugates, giving
angular weights $\pm1$.  For the quadratic sum we use Weil's local
index and product formula \cite[Chapter~II]{Weil}, in the explicit
normalization of \cite[Lemma~3.4, equation~(3.13)]{P}:
the normalized indices satisfy
$w_v(a)w_v(b)=(a,b)_{2,v}w_v(ab)$ and $\prod_vw_v(a)=1$.
Thus $\gamma_3(\mathfrak p)=w_S(p^{-1})^{-1}$ is a square-class
character on $B$, since its multiplicative defect is $H_S^3=1$ there.
Extend it as in (iii) and multiply by $\overline{\chi_\bullet(4)}$ to
obtain $G$.  Finally the finite-field duplication identity
$g_1g_4=\overline{\chi(4)}g_2g_3$, together with $g_2g_4=q$, gives the
last identity.  Chinese remaindering gives, for coprime squarefree ideals,
\begin{equation}\label{eq:gauss-crt}
 \gamma_j(\mathfrak c_1\mathfrak c_2)
 =\chi_{\mathfrak c_1}(c_2)^j\chi_{\mathfrak c_2}(c_1)^j
      \gamma_j(\mathfrak c_1)\gamma_j(\mathfrak c_2).
\end{equation}
These are the calculations of \cite[Lemmas~3.2--3.4]{P}; all their
hypotheses have now been checked over $K$.
\end{proof}

\subsection{The exact theta period with fixed infinity type}

We use the determinant-modified Kazhdan--Patterson cover with parameters
$(d,n,c)=(2,3,2)$, divided by its split scalar subgroup.  Its exceptional
datum has normalized exponents $1/6,-1/6$.  The construction of \cite{KP},
with the normalization in \cite[Proposition~5.1 and Appendix~A]{P},
applies over $K$.  Indeed the scalar cocycle has exponent
$1+4c=9$, so the scalar section is central and split; on a coroot cube
the selected datum has value $|t|_v$, the exceptional condition.
The rational splittings are compatible by reciprocity and the product
formula.  A rational ratio that is a cube at every completion is a cube
in $K$, since otherwise its Kummer extension has a nonsplit unramified
prime by Chebotarev.  This verifies the global compatibility of the
datum.  The residual representation is the one in
\cite[Theorem~II.2.1]{KP}; its Whittaker coefficients are given by
\cite[Theorem~II.2.2]{KP}, read with the correction \cite{KPcorr}.
Nonvanishing of the global functional and uniqueness of the local
models use \cite[Theorem~II.2.5(a),(b)]{KP}.
The correction inserts the finite-$S$ normalizing factor $T(S)^{-1}$,
with
$T(S)=\prod_{v\in S_f}(1+q_v^{-1})$ in rank two.
We use the normalized-functional realization in
\cite[Appendix~A.2]{P}, which incorporates this correction.
At a complex place the representation is the spherical principal series
of $\operatorname{PGL}_2(\mathbb C)$ with these exponents.  Its constant
term has radial complex inducing law, for every fixed right argument.

Use the rational Weyl representative
$w=\left(\begin{smallmatrix}0&1\\1&0\end{smallmatrix}\right)$ of
\cite[Section~6.1]{P}.  The normalized exterior spherical Whittaker values,
in the convention of \cite[Proposition~5.1(ii)]{P}, are
\begin{equation}\label{eq:theta-spherical-values}
 W_{\mathfrak p}(a(\varpi^e u))=\chi_{\mathfrak p}(u)^{2e}
 \begin{cases}
 q_{\mathfrak p}^{-j},&e=3j\ge0,\\
 q_{\mathfrak p}^{-j-1/2}\gamma_2(\mathfrak p),&e=3j+1\ge1,\\
 0,&e<0\text{ or }e\equiv2\pmod3.
 \end{cases}
\end{equation}
The local identities and \eqref{eq:gauss-crt} show that the exterior
coefficient at an allowed $A$ is zero unless
$\mathfrak A=\mathfrak c\mathfrak n^3$, with $\mathfrak c$ squarefree,
and is then $\gamma_2(\mathfrak c)/(q_{\mathfrak c}^{1/2}q_{\mathfrak n})$.
The ideals $\mathfrak c$ and $\mathfrak n$ may share primes.

Fix a unitary target $\psi$ and put
$\vartheta=\overline\Lambda\,\overline G\psi$.
Let $\mathcal L$ be its actual $S$-period character, so
$\mathcal L(A)=\vartheta(\mathfrak A)$ on allowed generators and
$\mathcal L(R^\times)=1$.  At a complex place write
\begin{equation}\label{eq:period-infinity-type}
 \mathcal L_v(z)=(z/|z|_{\rm ord})^{m_v}|z|_v^{i\tau_v},
 \qquad m_v\in\mathbb Z,\quad\tau_v\in\mathbb R.
\end{equation}
All relative radial parameters are retained.  In ideal-character
convention $m_v=k_v-b_v$, where $k_v$ is the target angular weight and
$b_v\in\{\pm1\}$ is the weight of $\Lambda$.  Replacing $\psi$ by its
conjugate if necessary ensures $m_v\ne0$ at some place, since $k-b$ and
$-k-b$ cannot both vanish.  This orientation is fixed throughout its
finite-order twists.

\begin{proposition}[Fixed-type theta reflection]\label{prop:fixed-type-theta}
For the preceding target, the unmarked theta periods, including their
whole-index coprimality masks and finite sectors, admit exact reflection for every
smooth test supported in a fixed compact subinterval of $(0,\infty)$.
For every fixed finite $S$-branch the reflected transform satisfies
\begin{equation}\label{eq:theta-kernel-bound}
 |(x\partial_x)^h\Phi_{\kappa,V}(x)|
 \ll_{K,S,\mathcal L,A,h}p_J(V)\min\{x^{1/4},x^{-A}\},
 \qquad A\ge1,\quad J>4rA+h+1.
\end{equation}
The same identities give exact Gaussian smoothing.  A moving norm phase
$V(x)x^{iT}$ has a fixed polynomial cost in $1+|T|$.
\end{proposition}

\begin{proof}
At a complex place normalize the holomorphic and antiholomorphic upper
root operators $E,\bar E$ so that they multiply the Whittaker function by
$z,\bar z$, respectively; this follows from Whittaker equivariance
\cite[(5.10)]{P}.  Starting from the spherical vector, take
$\phi_m=E^m\phi^\circ$ for $m\ge0$ and
$\phi_m=\bar E^{-m}\phi^\circ$ for $m<0$.  Then
\begin{equation}\label{eq:higher-ktype}
 W_{\phi_m}(a(z))=|z|_v^{a_v}(z/|z|_{\rm ord})^m
                    K_{1/3}(c_v|z|_v^{1/2}),
 \qquad a_v=(|m|+1)/2,\quad c_v>0.
\end{equation}
Spherical invariance gives $(F-\bar E)\phi^\circ=0$, and $F$ commutes
with $\bar E$.  Our Weyl representative satisfies
$\operatorname{Ad}(w)E=F$ and fixes $\phi^\circ$, so
$\pi(w)\phi_m=\phi_{-m}$.  Thus the reflected angular factor is exactly
the opposite one.  The other customary Weyl representative introduces
only the corresponding fixed sign.

At the finite places choose a common compact open product
$U\subset K_{S_f}^\times$ inside the sixth powers and the kernels of the
finitely many unit restrictions of $\mathcal L$ that occur.  A vector
with Whittaker function $1_U$ exists by \cite[Lemma~5.3]{P}: convolve a
nonzero Whittaker vector by an upper-unipotent Schwartz function whose
additive Fourier transform is $1_U$, and normalize at $1$.
This construction uses the open kernels on units, irrespective of the
orders of uniformizer values.

Let $X_l=\{y:|y|_S=l\}$ and choose measures so that
$d^\times y=d_l^\times y\,dl/l$.  The unreflected sheet profile is
\[
 V_0(l)=\int_{X_l}1_U(y_f)
       \prod_{v\mid\infty}t_v^{a_v-i\tau_v}
                         K_{1/3}(c_v\sqrt{t_v})\,d_l^\times y,
 \qquad t_v=|y_v|_v.
\]
We use $\widehat V(s)=\int_0^\infty V(x)x^s\,dx/x$.
The Bessel integral and absolute Fubini give
\begin{equation}\label{eq:theta-direct-mellin}
 \widehat V_0(s)=C_0A_0^s
       \prod_{v\mid\infty}\prod_{\epsilon=\pm1}
        \Gamma(a_v+s-i\tau_v+\epsilon/6),
 \qquad C_0\ne0,\quad A_0=\prod_v4/c_v^2>0.
\end{equation}
Here the finite Mellin factor is the constant $\operatorname{vol}(U)$.

We justify the global identity before dividing these factors.
Fix a nonzero row $m\in R$.  If
$m=\epsilon_m m_0$ with $m_0$ allowed, put
$\mathcal L_m(y)=\mathcal L(y)H_S(\epsilon_m,y)^{-1}$.
This belongs to a fixed finite family with the same infinity type.
For $j_{\mathfrak p}=v_{\mathfrak p}(m)\bmod6$, retain a digit even when
$j_{\mathfrak p}=0$.  An additional whole-index coprimality mask is
obtained by inserting the zero digit at each of its new primes.
The coefficients
\[
 \alpha_{\mathfrak p,j}(v)=
 \begin{cases}
 \tau_{\mathfrak p,j}^{-1}\chi_{\mathfrak p}(v)^{-j},
       &1\le j\le5,\ v\ne0,\\
 0,&1\le j\le5,\ v=0,\\
 1-q_{\mathfrak p}^{-1},&j=0,\ v=0,\\
 -q_{\mathfrak p}^{-1},&j=0,\ v\ne0,
 \end{cases}
 \quad
 \tau_{\mathfrak p,j}=\sum_{v\ne0}
        \chi_{\mathfrak p}(v)^{-j}\psi_{\mathfrak p}(v/p)
\]
satisfy $\sum_v\alpha_{\mathfrak p,j}(v)
\psi_{\mathfrak p}(tv/p)=\chi_{\mathfrak p}(t)^j$ with the
specified zero extension; here $p$ is the chosen allowed generator of
$\mathfrak p$.  Form the combined digit sums $D_m^+(y)$ and
$D_m^-(y)$ by evaluating theta at $g(y,v)$ and $wg(y,v)$, respectively,
where $g$ has component $a(y)$ at $S$ and $n(v/p)$ at row
primes.  Rational automorphy gives $D_m^+=D_m^-$.
They descend under $\Gamma$: replacing $v$ by $\gamma v$ gives the
rational diagonal $a(\gamma)$ before $w$, $a(\gamma^{-1})$ after $w$,
and the same exterior compact factor on the far right.  These compact
factors fix the spherical vectors, and $\alpha(\gamma v)=\alpha(v)$.
The constant terms descend by the same root relation and the product
formula; this is the exact argument in \cite[(6.6)--(6.7)]{P}.

Consequently the period
\[
 \mathcal P_m(l)=\int_{X_l/\Gamma}
       1_B([y])\mathcal L_m(y)^{-1}D_m^\pm(y)\,d_l^\times y
\]
is well defined on a compact quotient.  Rotation at a complex place with
$m_v\ne0$ kills both combined constant terms: their inducing law is
radial, including at the finite right arguments after $w$, whereas the
weight has nontrivial angular character.
The signed orbit substitution is
\begin{equation}\label{eq:theta-orbit-unfolding}
 \sum_{\gamma\in\Gamma}\int_{\Omega_l}
 1_B([y])\mathcal L_m(y)^{-1}f(\gamma Ay)\,d_l^\times y
 =\mathcal L_m(A)\int_{X_{q_A l}}
 1_B([A]^{-1}[v])\mathcal L_m(v)^{-1}f(v)\,d^\times v.
\end{equation}
The sheet measures are invariant under this substitution.  On the direct cutoff $[v]=1$;
the torus cocycle is $H_S(\nu,y)^2$, which is $1$ there by alternation.
The condition $[y]\in B$ then forces $\nu=\gamma A$ with $A$ allowed.
Thus the direct expansion is exactly
\begin{equation}\label{eq:theta-direct-period}
 \mathcal P_m(l)=
 \sum_{\mathfrak c\ {\mathrm{squarefree}},\,\mathfrak n}
 \frac{\gamma_2(\mathfrak c)\vartheta(\mathfrak c\mathfrak n^3)
              \chi_{\mathfrak c\mathfrak n^3}(m)}
      {q_{\mathfrak c}^{1/2}q_{\mathfrak n}}
       V_0(q_{\mathfrak c\mathfrak n^3}l).
\end{equation}

All these interchanges are absolute.  At a finite place every fixed
smooth Whittaker function vanishes for sufficiently large $|y|_v$.
For small $|y|_v$, the ordinary Jacquet module gives a recurrence on
triple valuation steps with eigenvalue moduli $q_v^{-1},q_v^{-2}$;
its polynomial-times-geometric solutions are $O(|y|_v^\delta)$ for
every $\delta<1/3$ \cite[(5.13)]{P}.  The complex Bessel functions and
their Euler derivatives have the corresponding positive small-end
power and arbitrary large-end decay.  Integrating on norm sheets gives
$O(\min(l^\alpha,l^{-A}))$ for every $\alpha<\delta$, also with any
fixed unitary character and finite power-class indicator.  One can see
this directly by taking one complex norm as the dependent coordinate;
the remaining logarithmic integrals converge after a smaller positive
power is reserved.  This is \cite[Lemma~5.4]{P}.
Together with the convergent exterior coefficient majorant for
$\Re s>1/2$, it proves absolute orbit unfolding.

Expand the reflected digits using
$wn(H)=n(H^{-1})a(-H^{-2})\bar n(H^{-1})$ and
\eqref{eq:theta-spherical-values}.  The finite coefficients are the
exact finite-field identities of \cite[Lemmas~6.2--6.3]{P}; their
hypotheses are $q_{\mathfrak p}\equiv1\pmod{12}$ and the retained zero
extensions, both valid here.  After the finite power-class separation,
write the resulting identity as
\begin{equation}\label{eq:theta-reflected-period}
 \mathcal P_m(l)=\sum_\kappa c_{\kappa,m}
       \sum_{\mathfrak a}b_{\kappa,m}(\mathfrak a)
       H_\kappa(\lambda_{\kappa,m}q_{\mathfrak a}/l),
 \qquad \lambda_{\kappa,m}>0.
\end{equation}
The exterior branch sum is finite for each row.  The $S$-functions
$H_\kappa$ belong to a fixed finite list; all $S$-valuations are integrated in the sheets.
The exterior coefficients and scales in this formula are given by
that finite Weyl operation.  The raw valuation table and the full
expansion with additional divisibility marks are proved separately in
Theorem~\ref{thm:marked-expansion}; the present argument uses only
the unmarked expansion of \Psrc.  

Let $\mathcal D_m(s)$ and $\mathcal D_{\kappa,m}(s)$ be the Dirichlet
series with the respective exterior coefficients in
\eqref{eq:theta-direct-period} and \eqref{eq:theta-reflected-period}.
The two expansions give rapid decay of $\mathcal P_m$ at opposite ends,
with all Euler derivatives.  Its entire Mellin transform therefore gives
\begin{equation}\label{eq:theta-completed-equation}
 \widehat V_0(s)\mathcal D_m(s)
 =\sum_\kappa c_{\kappa,m}\lambda_{\kappa,m}^{s}
       \widehat H_\kappa(-s)\mathcal D_{\kappa,m}(-s).
\end{equation}
Reciprocal gamma functions continue $\mathcal D_m$ to an entire
function.  On a right line its series is bounded; on a left line
\eqref{eq:theta-completed-equation} and Stirling give polynomial growth.
Inside the strip its preliminary growth is at most polynomial times
$e^{\pi r|t|}$.  Dividing by a sufficiently large power of a linear
factor with no zero in the strip and applying the maximum principle
after multiplication by $e^{\varepsilon s^2}$ proves polynomial strip
growth, as in \cite[proof of Proposition~5.1, p.~33]{Q}.

In the reflected sheet the character is $\mathcal L_m$, and the complex
vector is $\phi_{-m_v}$.  Consequently
\begin{equation}\label{eq:theta-multiplier}
 R_\kappa(s)=\frac{\widehat H_\kappa(s)}{\widehat V_0(-s)}
 =C_\kappa A_0^{2s}B_\kappa(s)
   \prod_{v\mid\infty}\prod_{\epsilon=\pm1}
   \frac{\Gamma(a_v+s+i\tau_v+\epsilon/6)}
        {\Gamma(a_v-s-i\tau_v+\epsilon/6)}.
\end{equation}
The finite factor $B_\kappa$ is holomorphic for $\Re s>-1/3$ and
bounded on every closed vertical substrip, by the preceding finite
Whittaker bounds.  The only finite factor divided out was
$\operatorname{vol}(U)\ne0$.
For an annular $V$, contour shifting in
\eqref{eq:theta-completed-equation} now gives the exact transform
\begin{align}\label{eq:theta-annular-transform}
 \Phi_{\kappa,V}(x)&=\frac1{2\pi i}\int_{(0)}
       \widehat V(-s)R_\kappa(s)x^{-s}\,ds,\\
 \sum_{\mathfrak a}a_m(\mathfrak a)V(q_{\mathfrak a}/X)
   &=\sum_\kappa c_{\kappa,m}\sum_{\mathfrak a}b_{\kappa,m}(\mathfrak a)
       \Phi_{\kappa,V}(\lambda_{\kappa,m}q_{\mathfrak a}X),\notag
\end{align}
where $a_m$ denotes the direct coefficient in
\eqref{eq:theta-direct-period}.  Lemma~\ref{lem:gamma-reflection} below
proves \eqref{eq:theta-kernel-bound}.  Finally define
\[
 K(T)=\frac1{2\pi i}\int_{(\sigma)}
       \frac{e^{s^2}}{\widehat V_0(s)}T^s\,ds,
 \qquad \widehat K(-s)=\frac{e^{s^2}}{\widehat V_0(s)}.
\]
The quotient is entire and has Gaussian vertical decay times a fixed
exponential and polynomial.  Moving the line in either direction gives
every absolute power moment of $K$ and its Euler derivatives.
Absolute Mellin Fubini therefore proves the exact identity
\[
 \int_0^\infty K(T)V_0(Tx)\,\frac{dT}{T}
       =\frac1{2\sqrt\pi}\exp\!\left(-\frac{\log^2x}{4}\right).
\]
The reflected profile is the equally convergent
$\int_0^\infty K(T)H_\kappa(x/T)\,dT/T$.
\end{proof}

\begin{lemma}[A common gamma-quotient estimate]\label{lem:gamma-reflection}
Suppose
\[
 R(s)=C A_*^s B(s)\prod_{j=1}^d
       \frac{\Gamma(a_j+s+it_j)}{\Gamma(a_j-s-it_j)},
 \qquad A_*>0,\quad a_j\ge1/3,
\]
with fixed parameters, where $B$ is holomorphic for $\Re s>-1/3$ and
bounded on closed vertical substrips.  For $A\ge1$, an integer $h\ge0$
and $J>2dA+h+1$, the transform
$\mathcal H_RV(x)=(2\pi i)^{-1}\int_{(0)}\widehat V(-s)R(s)x^{-s}\,ds$
satisfies
\begin{equation}\label{eq:unified-reflection-bound}
 |(x\partial_x)^h\mathcal H_R(Vx^{iT})(x)|
 \ll p_J(V)(1+|T|)^J\min\{x^{1/4},x^{-A}\}.
\end{equation}
\end{lemma}
\begin{proof}
There is no numerator pole in $\Re s\ge-1/4$.  On a fixed vertical line,
the opposite imaginary arguments cancel the exponential factors in
Stirling's formula, separately in each quotient, giving
$|R(\sigma+it)|\ll(1+|t|)^{2d\sigma}$.
Shift to $\sigma=-1/4$ or $A$ and integrate by parts $J$ times in the
Mellin transform.  Euler derivatives insert only $(-s)^h$, and
$p_J(Vx^{iT})\ll(1+|T|)^Jp_J(V)$.  The same bounds justify the
horizontal joins.
\end{proof}

For theta take $d=2r$ and the parameters $a_v\pm1/6\ge1/3$.
For ordinary Hecke reflection, if the completed local factors are
$\Gamma_{\mathbb C}(s+|k_v|/2+i\tau_v)$, its Mellin quotient is
\[
 \widehat W^\sharp(s)=\widehat W(1-s)
       \prod_{v\mid\infty}
       \frac{\Gamma(s+|k_v|/2-i\tau_v)}
            {\Gamma(1-s+|k_v|/2+i\tau_v)},
\]
up to fixed positive exponential scales.  Set
$U(x)=x^{1/2}W^\sharp(x)$ and $V(x)=x^{1/2}W(x)$.
Then $\widehat U(s)=\widehat V(-s)R(s)$ with $d=r$,
$a_v=1/2+|k_v|/2$ and $t_v=-\tau_v$.  The conductor factor gives dual
norm length $Q/X$, with $Q=|d_K|\operatorname N\mathfrak f$.
Apply the common-mask identities before replacing a presentation by its
primitive inducing character \cite[(18.4)--(18.8)]{Q}.
If the inducing character is not a pure norm character, the reflection
has no residue.  A pure norm phase is absorbed into the moving height,
and the resulting principal residue is retained in the exceptional
calculation.
The central small-annulus mass $O(2^{-j/4})$ is summable, which is the
only small-end property used in \cite[Section~18.1]{Q}.
Theta retains the precise $x^{1/4}$ envelope used in
\cite[Lemma~5.3]{Q}.

\subsection{Whole-index marked reflection}\label{sec:full-marked-reflection}

The following exterior expansion supplies the coefficient independence
needed for the quadratic--cubic norm estimate.  The period character in this subsection
may be any unitary $\vartheta$ unramified outside $S$ whose $S$-character
has a nonzero angular weight.  The physical choice is
$\bar\Lambda\bar G\psi$; the canonical choice will be
$\bar\Lambda\nu$ times a fixed finite-order character.

Write $U_6=R^\times/\Gamma$, and fix representatives of this finite group.
For a row $m=\epsilon_m m_0$, with $m_0$ allowed, set
$\mathcal L_\upsilon(y)=\mathcal L(y)H_S(\epsilon_m,y)^{-1}$,
where $\upsilon=\epsilon_m\Gamma$.  Isotropy of the unit image makes
this character trivial on $R^\times$.  Its angular and relative radial
parameters are those of $\mathcal L$.

\begin{lemma}[Local mark digits]\label{lem:mark-digits}
For an exterior prime $p$ put $Q=q_p$ and
$\mu_p(v)=Q^{-1}$ for every $v\in R/p$.  Then
\[
 \sum_{v\bmod p}\mu_p(v)\psi_p(tv/p)=\ind_{p\mid t}.
\]
The zero digit, after reflection, contributes $Q^{-1}W_p(a(\nu))$.
For a unit $t$ the nonzero digits contribute
\begin{equation}\label{eq:mark-active-local}
 \sum_{v\ne0}\mu_p(v)W_p(a(p^{e-2}t)wn(v/p))
 =\begin{cases}
 Q^{-1/2}\gamma_2(p)\chi_p(t)^2,&e=0,\\
 0,&e\ne0.
 \end{cases}
\end{equation}
For comparison, the reflected nonzero-digit part for an ordinary exponent
$j\in\{0,\ldots,5\}$ is $s_{p,j,e}\chi_p(t)^{b_{j,e}}$, where
\begin{equation}\label{eq:can-weyl-table}
 \begin{aligned}
 b_{j,0}&\equiv8-j\pmod6,&
 s_{p,j,0}&=\tau_{p,j}^{-1}
       \sum_{x\ne0}\chi_p(x)^{j-4}\psi_p(x/p) &&(j\ne0),\\
 b_{0,0}&\equiv2\pmod6,&
 s_{p,0,0}&=-Q^{-1/2}\gamma_2(p),\\
 b_{j,e}&\equiv4+2e\pmod6,&
 s_{p,j,e}&=\ind_{j=4}(Q-1)\tau_{p,4}^{-1}W_p(a(p^e)) &&(e>0).
 \end{aligned}
\end{equation}
All coefficients with $e<0$ vanish.  For $e>0$ they also vanish unless
$e\equiv0,1\pmod3$.  At a row prime with $j=1$ only $e=0$ occurs,
and $s_{p,1,0}=\omega_p$ has modulus one.  In the $j=0$ ordinary
branch the zero digit contributes $(1-Q^{-1})W_p(a(\nu))$.
\end{lemma}
\begin{proof}
The first identity is additive orthogonality.  The zero-digit assertions
follow from compact invariance of the spherical vector.  The local Weyl
identity of \cite[Lemma 6.2]{P}, before choosing digit weights, is
\[
 \sum_{v\ne0}\beta(v)W_p(a(p^{e-2}t)wn(v/p))
 =W_p(a(p^e))\chi_p(t)^{4+2e}
       \sum_{v\ne0}\beta(v)\chi_p(v)^4\psi_p(p^{e-1}t/v)
\]
for $e\ge0$, with zero left-hand side for $e<0$.  This identity is
linear in the weights $\beta$.  For $\beta=\mu_p$ and $e>0$,
$\chi_p^4$ is nonprincipal and its sum is zero.  At $e=0$ substitute
$x=t/v$.  The Gauss sum of $\chi_p^{-4}=\chi_p^2$ contributes
$Q^{1/2}\gamma_2(p)$, and the remaining unit power is $\chi_p(t)^8$.
This proves \eqref{eq:mark-active-local}.  The same substitution with
$\beta=\alpha_{p,j}$ gives \eqref{eq:can-weyl-table}.  In particular,
for $j=4$ its positive-valuation sum is $(Q-1)/\tau_{p,4}$; for all
other $j$ it vanishes.  These are raw local factors in the convention
of $\Psrc$.  Their globally paired cubic exponent is $-2$, as computed below.
\end{proof}

Fix a factorization $(m)=h d$, where $d$ is squarefree,
$(h,d)=1$, and $v_p(m)=1$ for $p\mid d$.
The primes of $h$ are frozen, and may also have valuation one.
Let $r$ be a squarefree puncture with $(r,d)=1$, put
$\mathcal P_0=\operatorname{supp}(h)\cup\operatorname{supp}(r)$,
and take a finite set of mark primes $\mathcal Q$ disjoint from
$\operatorname{supp}(d)\cup\mathcal P_0$.
At $p\in\mathcal P_0$ put $j_p=v_p(m)\bmod6$; a new puncture has
$j_p=0$.  For a test $V$ define
\begin{equation}\label{eq:marked-primal-exact}
 T_{\mathcal Q}(X;m,r,V)=
 \sum_{\substack{c\ \mathrm{sf}\\n}}
 \frac{\gamma_2(c)}{q_c^{1/2}q_n}\,
 \vartheta(cn^3)\chi_{cn^3}(m)\ind_{(cn^3,r)=1}
 \prod_{q\in\mathcal Q}\ind_{q\mid cn^3}\,
 V(q_cq_n^3/X).
\end{equation}
Here $c$ ranges over squarefree exterior ideals and $n$ over all
exterior ideals, allowing common factors.
If a mark meets $m$ or $r$, the primal expression is identically zero;
such a tuple is removed before making the expansion.

An active branch is a partition $\mathcal Q=\mathcal A\sqcup\mathcal I$;
write $P=\prod_{q\in\mathcal A}q$.
Let $f$ run over squarefree products of primes of $\mathcal P_0$
containing every prime with $j_p\ne0$, and put
$ I_f=\prod_{p\in\mathcal P_0,\,p\nmid f}(1-q_p^{-1})$.
Let $a'$ be supported on $f$, with every exponent $0$ or $1$ modulo $3$.
We retain all actual zero coefficients in this convenient larger support.
For $p\mid f$ abbreviate
$e_p=v_p(a')$, $s_p=s_{p,j_p,e_p}$ and $b_p=b_{j_p,e_p}$.
Choose a frequency-unit representative $\epsilon\in U_6$.
The variables $c',n'$ satisfy $c'$ squarefree,
$(c',f)=1$ and $(n',fP)=1$.
Use allowed generators for all displayed ideal products and put
\begin{equation}\label{eq:marked-frequency}
 \nu_0=\frac{\epsilon a'c'(n')^3}{(dfP)^2},\qquad
 \kappa=[\epsilon][a'][c'][n']^3[d]^{-2}[f]^{-2}[P]^{-2}\in V.
\end{equation}
The symbol zeros encode the remaining restrictions at $d$ and $P$.

For $L=\mathcal L_\upsilon$ and $t=(f,a',n',\epsilon)$ define
\begin{align}
 \Omega_t={}&I_f L(a')^{-1}L(n')^{-3}L(f)^2
   \prod_{p\mid f}s_p
    \chi_p\bigl(\epsilon(a'/p^{e_p})(n')^3(f/p)^{-2}\bigr)^{b_p},
       \label{eq:marked-Omega}\\
 \rho_t(d)={}&L(d)^2
   \prod_{p\mid d}\omega_p
       \chi_p\bigl(\epsilon a'(d/p)^{-2}f^{-2}\bigr)
   \prod_{p\mid f}\chi_p(d)^{-2b_p},\label{eq:marked-rho}\\
 \pi_t(P)={}&L(P)^2
   \prod_{q\mid P}\gamma_2(q)
       \chi_q\bigl(\epsilon a'(P/q)^{-2}f^{-2}\bigr)^2
   \prod_{p\mid f}\chi_p(P)^{-2b_p},\label{eq:marked-pi}\\
 \lambda_t(c')={}&L(c')^{-1}
   \prod_{p\mid f}\chi_p(c')^{b_p}
   \prod_{l\mid c'}\gamma_2(l)
       \chi_l\bigl(\epsilon a'(c'/l)f^{-2}\bigr)^2.
       \label{eq:marked-lambda}
\end{align}
Extend $\rho$ by zero off squarefree ideals prime to $\mathcal P_0$,
$\pi$ by zero off squarefree ideals prime to $\mathcal P_0$,
and $\lambda$ by zero off squarefree ideals prime to $f$.
In particular, $\rho$ is independent of $P,c',n'$,
$\pi$ is independent of $d,c',n'$, and $\lambda$ is independent of
$d,P,n'$.  Their absolute values are at most one.
The only dependence of $\Omega$ on moving dual variables is allowed
inside the dual coefficient.

For absolute majorants set
\begin{equation}\label{eq:frozen-beta}
 \beta_p(e)=
 \begin{cases}
 1,&e=0,\\ q_p^{-k},&e=1+3k,	k\ge0,\\
 q_p^{1/2-k},&e=3k,	k\ge1,\\0,&\text{otherwise},
 \end{cases}
 \qquad \beta_f(a')=\prod_{p\mid f}\beta_p(e_p).
\end{equation}
Then $|\Omega_t|\le\beta_f(a')$ by the local table.

\begin{theorem}[Separated marked expansion]\label{thm:marked-expansion}
With the preceding data, fix $\upsilon=\epsilon_m\Gamma$.
The exact reflection of \eqref{eq:marked-primal-exact} is
\begin{align}\label{eq:marked-reflection-exact}
 T_{\mathcal Q}(X;m,r,V)
 ={}&\sum_{\mathcal Q=\mathcal A\sqcup\mathcal I}
       \left(\prod_{q\in\mathcal I}q_q^{-1}\right)q_P^{-1/2}
       \sum_{f,a',n',\epsilon}\ \sum_{s_d,s_c\in B}
       \frac{\Omega_t}{q_{n'}}\pi_t(P)\rho_t(d)
       \ind_{[d]=s_d}\ind_{(P,n')=1}\ind_{(P,d)=1}
       \notag\\[-2pt]
 &\hspace{4mm}\cdot\sum_{c'\ \mathrm{sf}}
       \frac{\lambda_t(c')\ind_{[c']=s_c}}{q_{c'}^{1/2}}
       \chi_d(c'n')^3\chi_P(c')^{-2}
       \Phi_{\upsilon,\kappa,V}
           \left(\frac{Xq_{a'}q_{c'}q_{n'}^3}{(q_dq_fq_P)^2}\right).
\end{align}
Here the $f,a',n',\epsilon,c'$ ranges are those above, and $\kappa$
is determined by \eqref{eq:marked-frequency}.  The mask $(P,d)=1$ is redundant under the stated
disjointness assumption, but is displayed for use in row norms.
The functions $\Phi_{\upsilon,\kappa,V}$ belong to the fixed finite
family of \eqref{eq:theta-annular-transform} and satisfy
\eqref{eq:theta-kernel-bound}.  The expansion converges absolutely.
Besides the common kernel, the only factors coupling different moving
kinds of variables are
\begin{equation}\label{eq:marked-only-couplings}
 \chi_d(c'n')^3\chi_P(c')^{-2}
       \ind_{(P,n')=1}\ind_{(P,d)=1}.
\end{equation}
The same assertion holds for a test whose Mellin transform is entire
and rapidly decreasing on every closed vertical strip, with the
required weighted Euler seminorms finite at both ends; this includes the fixed
Gaussian used in the physical construction.
\end{theorem}
\begin{proof}
\emph{The automorphic identity.}
In the digit average defining $D_m^\pm$, use $\alpha_{p,1}$ at $d$,
$\alpha_{p,j_p}$ at $\mathcal P_0$, and $\mu_q$ at the marks.
The digit weights are invariant under multiplication by $\Gamma$.
Consequently the same quotient descent as in
Proposition~\ref{prop:fixed-type-theta} applies.  Both constant terms
are radial at every complex place.  Integration against
$L^{-1}$ kills them by rotation at a place with nonzero angular weight;
the relative radial parameters do not affect this cancellation.
Unfolding the unreflected side and using Lemma~\ref{lem:mark-digits}
gives
\[
 F(l)=\sum_A a^+(A)V_0(q_A l),
\]
where $a^+(A)$ is precisely the arithmetic coefficient of
\eqref{eq:marked-primal-exact}.  The finite expansion acts on the digit functions.

\emph{Frequencies and phases.}
At a row prime the local valuation is $-2$.
At an active ordinary prime it is $e_p-2$; at an active mark it is
$-2$.  Elsewhere the spherical support consists of valuations $3k$
and $3k+1$.  These conditions give exactly the frequencies
$\gamma\nu_0$, $\gamma\in\Gamma$, with the indices in
\eqref{eq:marked-frequency}.  Their parametrization is unique:
reduce valuations modulo $3$ outside $f$, and use allowed-generator
uniqueness for the remaining unit.  Frequencies with a locally zero
$s_p$ still carry coefficient zero.

For $p\mid d$, $p\mid f$, and $q\mid P$, respectively, the unit parts are
\[
 \epsilon a'c'(n')^3(d/p)^{-2}f^{-2}P^{-2},\quad
 \epsilon(a'/p^{e_p})c'(n')^3d^{-2}(f/p)^{-2}P^{-2},\quad
 \epsilon a'c'(n')^3d^{-2}f^{-2}(P/q)^{-2}.
\]
At $l\mid c'$ the spherical unit factor is
$\chi_l(\epsilon a'(c'/l)(dfP)^{-2})^2$; its cube contribution has
exponent divisible by six.  The $S$-character separates as
\[
 L(\nu_0)^{-1}=
 L(a')^{-1}L(c')^{-1}L(n')^{-3}L(d)^2L(f)^2L(P)^2,
\]
since $L(\epsilon)=1$.
Using reciprocity, the mixed factors are exactly
\[
\begin{array}{ccl}
 d,c'&:&\chi_d(c')\chi_{c'}(d)^{-4}=\chi_d(c')^3,\\
 d,P&:&\chi_d(P)^{-2}\chi_P(d)^{-4}=\ind_{(d,P)=1},\\
 P,c'&:&\chi_P(c')^{2}\chi_{c'}(P)^{-4}=\chi_P(c')^{-2},\\
 d,n'&:&\chi_d(n')^3,\\
 P,n'&:&\chi_P(n')^6=\ind_{(P,n')=1}.
\end{array}
\]
These identities retain their zero extensions.  The self-interactions
and interactions with frozen primes are exactly
\eqref{eq:marked-Omega}--\eqref{eq:marked-lambda}.
This proves coefficient separation, including its asserted independence.

The $S$-unfolding gives the sheet function
\begin{equation}\label{eq:marked-H-explicit}
 H_{\upsilon,\kappa}(x)=
  \int_{X_x}\ind_B(\kappa[v]^{-1})H_S(\kappa,[v])^2
                 L(v) f_{w,\upsilon}(v)\,d_x^\times v.
\end{equation}
All $S$-valuations have been integrated here.  The exterior coefficient
is the product in \eqref{eq:marked-reflection-exact} with
$\Phi$ replaced by $H(|\nu_0|_S/l)$.
This gives an exact fixed-scale reflected expansion of $F(l)$.
The class multiplier is trivial at complex places, so its Mellin
multiplier is precisely \eqref{eq:theta-multiplier}.

\emph{Convergence and arbitrary tests.}
For $\sigma>1/2$, the absolute reflected Dirichlet majorant, apart
from the finitely many row and mark factors, is at most
\begin{equation}\label{eq:marked-dirichlet-majorant}
 \zeta_K^S(\sigma+1/2)\zeta_K^S(3\sigma+1)
 \prod_{p\mid f}\left(1+
       \frac{q_p^{-\sigma}+q_p^{-1/2-3\sigma}}
            {1-q_p^{-1-3\sigma}}\right).
\end{equation}
This follows by summing \eqref{eq:frozen-beta}, and by positive
extension of the squarefree and cube ranges.  It proves absolute convergence of the orbit sums.
For each fixed primal row and mark tuple, the direct and reflected
expansions make $F(l)$ rapidly decreasing at infinity and at zero,
respectively.  Thus $\widetilde F(s)=\int_0^\infty F(l)l^s\,dl/l$
is entire and bounded on fixed strips.  On a right half-plane it is
$\widehat V_0(s)\mathcal D_+(s)$.  On a left half-plane it is the
sum of $a^-(j)|\nu_0(j)|_S^s\widehat H_{\upsilon,\kappa(j)}(-s)$.
Since $1/\widehat V_0$ is entire, this continues $\mathcal D_+$.
Its preliminary strip growth is at most $e^{C|\Im s|}$ times a
polynomial.  On the two boundary lines the preceding majorant and
\eqref{eq:theta-multiplier} give polynomial growth.  The strip
Phragm\'en--Lindel\"of argument therefore gives polynomial growth
throughout the intervening strip.

Mellin inversion of $V(q_A/X)$ may now be shifted from a line
$\Re s>1/2$ to a line $\Re s<-1/2$, without residues.
Inserting the reflected series, interchanging absolutely, and putting
$w=-s$ gives
\[
 \frac1{2\pi i}\int_{(\sigma)}
   \widehat V(-w)\frac{\widehat H_{\upsilon,\kappa}(w)}
                           {\widehat V_0(-w)}
      (X|\nu_0|_S)^{-w}\,dw,
 \qquad \sigma>-1/3.
\]
This is $\Phi_{\upsilon,\kappa,V}(X|\nu_0|_S)$, and proves
\eqref{eq:marked-reflection-exact}.  The shifts to $-1/4$ and $A$
and Lemma~\ref{lem:gamma-reflection} give the claimed seminorm bounds.
The argument applies directly to the indicated two-ended test class;
for the Gaussian it agrees with the smoothing already constructed.
\end{proof}

\subsection{Weighted orbital Poisson summation}

For the rest of the section $\Gamma$ may be replaced by one fixed
finite-index subgroup of $(R^\times)^6$ to accommodate fixed sectors.
It remains fixed throughout all moving moduli and rows.
The following formula keeps finite arithmetic coefficients and fixed
unitary backgrounds in their respective variables.

\begin{lemma}[Weighted masked orbital Poisson]\label{lem:weighted-orbital-poisson}
Let $\chi:K_S^\times\to\mathbb C^\times$ be fixed, smooth and unitary,
with $\chi(\Gamma)=1$.  For $\phi\in\mathcal S(K_S)$ set
\[
 W_{\phi,\chi}(x)=\int_{X_x}\chi(y)^{-1}\phi(y)\,d^\times y,
 \qquad
 \widetilde W_{\phi,\chi}(x)=\int_{X_x}\chi(y)\widehat\phi(y)\,d^\times y.
\]
Let $f:R/cR\to\mathbb C$ be periodic and $\Gamma$-invariant.  With
continuous Fourier sign $e(-xy)$, define
$\widehat f_d(h)=\sum_{a\bmod c}f(da)e(ha/c)$.
For every exterior radical mask $\mathfrak R$ and $L>0$,
\begin{align}\label{eq:weighted-orbital-poisson}
 &\sum_{[k]\ne0}\chi(k)f(k)\mathbf1_{(k,\mathfrak R)=1}
       W_{\phi,\chi}(q_k/L)\notag\\
 &\quad=\frac{L\chi(c)}{q_c}
       \sum_{d\mid\operatorname{rad}\mathfrak R}\frac{\mu(d)\chi(d)}{q_d}
       \sum_{[h]\ne0}\bar\chi(h)\widehat f_d(h)
       \widetilde W_{\phi,\chi}(Lq_h/(q_dq_c))+Z_{\chi,f},\\
 Z_{\chi,f}
 &=\chi(\ell_L)\mathbf1_{\chi|_{T_0}=1}\operatorname{vol}(T_0/\Gamma)
 \left\{\frac{L\widehat\phi(0)}{q_c}
       \sum_d\frac{\mu(d)}{q_d}\widehat f_d(0)
       -\phi(0)f(0)\sum_d\mu(d)\right\}.\notag
\end{align}
Here brackets mean $\Gamma$-orbits, $q_k=|k|_S$, and $c,d$ are exterior
generators.  All zero extensions are retained in $f$ and the mask.
\end{lemma}
\begin{proof}
Integrate the full lattice sum, including its actual zero, over
$T_0/\Gamma$ with weight $\chi(\ell_L)\chi(u)^{-1}$.
Expand the mask before the finite transform.  Additive Poisson on the
self-dual lattice gives exactly
\[
 \sum_{n\in R}f(dn)\phi(dnu/\ell_L)
  =\frac{L}{q_dq_c}\sum_{h\in R}\widehat f_d(h)
                    \widehat\phi(\ell_Lh/(dcu)).
\]
The finite transform is $\Gamma$-invariant by the substitution
$a\mapsto\gamma a$.  Unfolding the nonzero original orbits gives the
left side of \eqref{eq:weighted-orbital-poisson}.  On the dual side,
inversion and Haar translation give the factor $\chi(dc)\bar\chi(h)$
and the sheet of norm $Lq_h/(q_dq_c)$.  The dual kernel is thus common to all modulus shapes.  The integral of
$\chi(u)^{-1}$ over the compact quotient is its volume if
$\chi|_{T_0}=1$ and is zero otherwise.  Retaining the dual zero and
subtracting the included primal zero proves the displayed $Z_{\chi,f}$.

\end{proof}

Every smooth annular norm test has such a lift: multiply a compact shape
function of integral one by the test and by $\chi(y)$.  Thus $\chi=1$
gives the additive interface, while fixed nontrivial $\chi$ gives the
character-equivariant interface for canonical descendants.  The factor
$\chi(dc)$ has modulus one; the finite Fourier table remains the table
of $f$.  If $f=\xi$ is primitive, its specialization is
$\widehat f_d(h)=\xi(d)\sqrt{q_c}\gamma(\xi)\bar\xi(h)$.
This includes the primitive principal case with its actual zero term.

We record two quantitative details.  First, a fixed compact unit
partition represents orbit sums by lattice sums with compact shape
support and finite multiplicity.
Second, for a fixed finite local type, the sheet integrals of a
Schwartz function and their Euler derivatives have arbitrary power
decay at infinity and at most $(1+|\log x|)^b$ growth at zero, for a
fixed $b=b(K,S)$.  This follows by bounding by product Schwartz functions:
the logarithmic complex norms and the finite valuation shells form a
simplex of logarithmic volume near zero.  Euler derivatives do not
increase its dimension.  On polynomial scale ranges these logarithms
consume any prescribed positive loss, with only the threshold depending
on the fixed data.  These conclusions hold for the weighted profiles
by unitarity and the fixed infinity parameters.

\begin{lemma}[Balanced lattice decay]\label{lem:balanced-lattice-decay}
Let $C$ be a fixed compact subset of $K_S^1$, and let $\Phi$ range over a
bounded Schwartz family on $K_S^n$ with fixed finite local type.  For
$a\in C$, $\rho\ge1$ and $A>0$,
\begin{equation}\label{eq:balanced-lattice-decay}
 \sum_{\mathbf h\in R^n\setminus\{0\}}
       |\widehat\Phi(\ell_\rho a\mathbf h)|\ll_A\rho^{-A}.
\end{equation}
The bound uses finitely many Schwartz seminorms; their orders may depend
on $A[K:\mathbb Q]$.
\end{lemma}
\begin{proof}
The finite support and the compact range of $a_f$ put all contributing
$\mathbf h$ in a fixed fractional $\mathcal O_K$-lattice.  At infinity
$\ell_\rho$ dilates by $\rho^{1/[K:\mathbb Q]}$, and $a$ ranges over a
compact family of invertible diagonal maps.  The lattice has a fixed
positive minimum distance.  A Schwartz order greater than
$\max\{A[K:\mathbb Q],n[K:\mathbb Q]\}$ makes its nonzero sum
$O(\rho^{-A})$.
\end{proof}

\subsection{The full additive Gram estimate}

We prove the Gram bound by two successive correlation calculations.  For an allowed generator $s$
put
\[
 g_{\mathfrak s}(k)=\sum_{a\bmod s}\chi_{\mathfrak s}(a)e(ka/s).
\]
Let $c(\mathfrak s)$ be a fixed finite linear combination of the finite
ray coefficients of the additive probe.  In particular it may contain
a fixed class indicator, a fixed-numerator residue character and the
inverse of a fixed character on its unit support.  Extend it by zero
before evaluating any such inverse.  The fixed set $S$ includes all
finite coefficient conductors and fixed numerator primes; thus the
remaining finite ray restrictions occur only in the fixed local
profiles at $S$.  For $W\in C_c^\infty(0,\infty)$,
define
\begin{equation}\label{eq:additive-polynomial}
 A_m(Y;\nu)=Y^{-3/2}\sum_{\mathfrak s}
 c(\mathfrak s)W(q_{\mathfrak s}/Y)
       (q_{\mathfrak s}/Y)^{-1+i\nu}g_{\mathfrak s}(-m).
\end{equation}
This is the exact central normalization in \cite[Section~6 and
Proposition~15.2]{Q}, now on exterior ideals.
The coefficient $c$ is independent of the row and of $\nu$.

\begin{proposition}[Additive Gram bound]\label{prop:additive-gram}
Let $Q,Y\ge1$ range in fixed polynomial ranges in $Z$, and suppose
$P_a=Y^2/Q\ge1$.  For every fixed row-ball constant and every $\eps>0$,
\begin{equation}\label{eq:additive-gram-bound}
 \sum_{\substack{[m]\ne0\\q_m\ll Q}}|A_m(Y;\nu)|^2
 \ll_{K,S,c,W,\eps}(1+|\nu|)^J\frac QY
             \left(1+P_a^{1/6}+\frac{P_a^2}{Y}\right)Z^\eps.
\end{equation}
The same bound holds for rows in a fixed compact additive shape.  The
seminorm order $J$ is fixed independently of the moving scales and
height.  The arithmetic coefficient is fixed independently of the row.
\end{proposition}

\begin{proof}
Choose a nonnegative Schwartz/compact-open majorant whose orbital
profile is at least one on the row ball.  This is possible by the
balanced representatives of Proposition~\ref{prop:arithmetic-calibration}.
The full lattice average includes a nonnegative term at $m=0$, which
may be retained for an upper bound.  Expand the square in
\eqref{eq:additive-polynomial}, keeping its one joint annular profile,
and write $s_i=Cn_i$ with $(n_1,n_2)=1$.

The complete finite correlation, with every zero extension retained, is
\begin{equation}\label{eq:full-additive-correlation}
 F(u,v;j)=\sum_{\substack{x\bmod u,\ y\bmod v\\
                         vx-uy\equiv j\ (uv)}}
                         \chi_u(x)\overline{\chi_v(y)}.
\end{equation}
Expanding the two Gauss sums and using additive orthogonality shows
that their complete row transform is $q_uq_vF(u,v;j)$.
The Poisson factor is $Q/(q_uq_v)$, so the exact common prefactor after
the central normalization is $Q/Y^3$.
The correlation is zero unless $C\mid j$.  For $j=Ck$, reciprocity on
the allowed generators gives
\begin{equation}\label{eq:correlation-common-part}
 F(Cn_1,Cn_2;Ck)
 =\chi_{n_1}(k)\overline{\chi_{n_2}(-k)}L_C(n_1,n_2;k),
\end{equation}
where $L_C$ is the corresponding congruence sum modulo $C$.
At $\mathfrak p^a\Vert C$, when both residual columns are units, its
local value is
\[
 q_{\mathfrak p}^{a-1}\chi_{\mathfrak p}(n_1/n_2)^a
 \begin{cases}
 q_{\mathfrak p}-1,&\mathfrak p\mid k,\\
 -1,&\mathfrak p\nmid k,\ 6\nmid a,\\
 q_{\mathfrak p}-2,&\mathfrak p\nmid k,\ 6\mid a.
 \end{cases}
\]
If exactly one column is divisible by $\mathfrak p$, the value is
$q_{\mathfrak p}^{a-1}(q_{\mathfrak p}-1)
\mathbf1_{6\mid a}\mathbf1_{\mathfrak p\nmid k}$.
These formulas follow by reduction to the residue field: in the
nonzero-unit case substitute the ratio $(1+t)/t$ and sum the character
over the units except $1$.  Define the local extension to be zero if
both residual columns meet $\mathfrak p$.  The resulting extension
satisfies $|L_C|\le q_C$ and is periodic modulo $\operatorname{rad}C$
in each column.  On the coprime residual locus it equals the genuine
correlation.  This proves over every residue field the complete
correlation identities used in \cite[Lemma~13.3]{Q}.

Insert the full identity
$\mathbf1_{(n_1,n_2)=1}=\sum_{d\mid n_1,\ d\mid n_2}\mu(d)$ before
using this extension, and write $n_i=dm_i$.  The residual variables $m_1,m_2$ may share factors.  Nonzero terms have $(d,Ck)=1$,
by the extension and the original character zeros.
The fixed sectors are included in fixed local profiles.  By
\eqref{eq:unit-reciprocity}, the remaining exterior coefficient is
periodic in the two columns modulo
\[
 \mathfrak r=\operatorname{lcm}(\operatorname{rad}C,
                              \operatorname{rad}(kR)),
 \qquad q_{\mathfrak r}\ll_{\rm fixed}q_Cq_k.
\]
Its normalized two-variable Fourier coefficients have magnitude at
most $O(q_C)$.  Call $k$ nonexceptional when an exterior prime outside
$C$ has valuation in $kR$ not divisible by $6$.  At such a prime one
column contains a nonprincipal zero-extended residue character; all
other local factors and fixed phases are independent of that residue.
Its complete mean is zero.  Thus the joint Fourier coefficient at
$(0,0)$ vanishes for every nonexceptional $k$.

The frequency $k=0$ uses the actual compact-quotient scalar
$\operatorname{vol}(T_0/\Gamma)\widehat\phi(0)$ from
Lemma~\ref{lem:weighted-orbital-poisson}.
It is exactly diagonal: nonzero residues in
\eqref{eq:full-additive-correlation} force $u=v$, and then
$F(u,u;0)=\varphi(u)$.  There are $O(Y)$ possible $C$ with $q_C\asymp Y$,
each contributing at most $O(Y)$.  Including $Q/Y^3$ gives $O(Q/Y)$.

For nonzero $k$, the first orbital Fourier kernel has argument
$Qq_Cq_k/(q_{s_1}q_{s_2})\asymp q_Cq_k/P_a$.
Use dyadic shells $q_Cq_k/P_a\asymp H$, $H=1,2,4,\ldots$, with the
first shell containing arguments at most two.  Any small-end
logarithms of the orbital kernel cost a preliminary $Z^\eps$ loss.
Put $q_C\asymp C_0$, $q_d\asymp D_0$ and
$N=Y/(C_0D_0)$.  For fixed $C,d,k$, keep a single joint two-column
profile $W_H$.  Its required homogeneous seminorms are
$O((1+|\nu|)^JH^{-B})$: derivatives of the first Fourier kernel are
absorbed by its rapid large-end decay, and norm-power derivatives cost
only fixed powers of $1+|\nu|$.  The factor $H^{-B}$ occurs once in the
linear two-column Fourier estimate.

Use a compact unit partition in both columns and balance an allowed
generator of $\mathfrak r$.  For a nonexceptional $k$, two-column
Poisson and Lemma~\ref{lem:balanced-lattice-decay} give
\begin{equation}\label{eq:gram-two-column-poisson}
 \sum_{\mathbf m\in R^2}a_{C,d,k}(\mathbf m)
                 W_H(\mathbf m/\ell_N)
 \ll(1+|\nu|)^JH^{-B}q_CN^2
                 \min\{1,(q_{\mathfrak r}/N)^A\}.
\end{equation}
Here $A>2$ is arbitrary.  For $N\ge q_{\mathfrak r}$ the zero
frequency is absent and the lemma gives the displayed decay; for
$N<q_{\mathfrak r}$ use the compact-shape point count $O(N^2)$.
A nonempty subunit annulus has $N$ bounded below by a fixed positive
constant.  The field degree enters through the fixed derivative order.

Set $K_H=HP_a/C_0$.  The ideal count, including the fixed unit sectors,
gives $O(K_H)$ nonzero frequency orbits in this shell; the range is
empty below a fixed positive lower bound for $K_H$.
There are $O(C_0D_0K_H)$ choices of $C,d,k$.  Multiplying
\eqref{eq:gram-two-column-poisson} by this count and by $Q/Y^3$ gives
\[
 \ll(1+|\nu|)^JH^{-B}\frac QY
       \frac{HP_a}{C_0D_0}
       \min\{1,(HP_aC_0D_0/Y)^A\}.
\]
For $\Lambda_0=Y/(HP_a)$, the positive dyadic sum is
\[
 \sum_{i,j\ge0}2^{-i-j}
       \min\{1,(2^{i+j}/\Lambda_0)^A\}
 \ll_A\frac{1+\log(2+\Lambda_0)}{\Lambda_0}.
\]
Split at $2^{i+j}\asymp\Lambda_0$ when $\Lambda_0\ge1$; when it is
smaller than one the assertion follows from the convergent geometric
sum.  Summing $H$ with $B>2$ gives the term
$(Q/Y)(P_a^2/Y)Z^\eps$.

An exceptional nonzero frequency has exterior ideal
$kR=\mathfrak a^6\mathfrak e$, where $\mathfrak e$ is sixth-power-free
and supported on $C$.  There are
$6^{\omega(C)}\ll_\eps q_C^\eps$ possibilities for $\mathfrak e$ and
$O(K_H^{1/6})$ for $\mathfrak a$ on a nonempty shell.  The fixed unit
sectors cost a constant.  Use $|L_C|\le q_C$ and the trivial
two-column count.  Before multiplication by $Q/Y^3$, the contribution is
\[
 \ll(1+|\nu|)^JH^{-B}Z^\eps
       Y^2(HP_a)^{1/6}C_0^{-1/6}D_0^{-1}.
\]
The $C_0,D_0$ sums converge after taking a smaller preliminary loss,
and the $H$ sum converges for the same fixed $B>2$.  This gives
$(Q/Y)P_a^{1/6}Z^\eps$.  Adding the three frequency types proves
\eqref{eq:additive-gram-bound}.  A fixed compact additive row shape is
covered by the same positive majorant.
\end{proof}

In the physical probe the fixed target acts on the completed theta
index $\mathfrak c\mathfrak n^3$, while
\eqref{eq:additive-polynomial} is independent of that target.
Its continuous phase is therefore handled by
Proposition~\ref{prop:fixed-type-theta}; the additive Gram calculation
uses $\chi=1$.  The weighted Poisson lemma handles the fixed live
backgrounds in Section~\ref{sec:canonical}.
All complete masks, collision zeros and common joint profiles have been
retained.  For every prescribed real loss the norm powers in these
lemmas are unchanged by $K$, $S$ or the fixed infinity type.  The
constants, finite derivative orders and eventual thresholds may depend
on those fixed data; increasing the derivative orders supplies the
higher-dimensional Poisson decay, and fixed unit logarithms consume
only the prescribed loss.

\section{Canonical moments}\label{sec:canonical}

The two estimates needed for row counting are stated first in
Theorems~\ref{thm:can-inverse} and~\ref{thm:can-fourth}.
Their proofs share the coefficient identities, the marked theta
terminal, and the profile-propagation lemma. The inverse proof is a
two-transform induction on row width
(Section~\ref{subsec:inverse-induction}); the ordinary fourth moment
adds a same-band comparison and an exceptional-row cancellation
(Section~\ref{subsec:fourth-induction}). This order allows the moment
statements to be read independently of the recursive bookkeeping.


We use the arithmetic normalization of
Proposition~\ref{prop:arithmetic-calibration}.  Thus $K\supset\mu_{12}$
is fixed, $R=\mathcal O_{K,S}$ is principal, and an exterior ideal is an
integral ideal supported outside $S$.  Its norm is denoted by $q_{\mathfrak a}$.
Ideal variables in character symbols are represented by allowed generators.
Element rows are taken modulo a fixed finite-index subgroup of
$(R^\times)^6$.  There are only a fixed number of such rows above each
exterior ideal.  A further finite unit-sector subdivision is allowed in a
positive moment estimate, and is recombined afterwards.

For an exterior ideal $\mathfrak n$, write
$\chi_{\mathfrak n}(u)$ for the sextic presentation, and put
$\chi_u(\mathfrak n)=\chi_{\mathfrak n}(u)$.
Every character power is extended by zero at a nonunit, including a power
divisible by six.  The modulus-one character is identically one.
All additional coprimality conditions below are imposed before an inverse
character is evaluated.  An inducing character means the primitive Hecke character, used
together with all displayed zero factors.

Fix a finite group $\Theta$ of finite-order Hecke characters, with
conductors supported in $S$, containing the finite Gauss and reciprocity
characters and the finite characters used in the local sector tables.
Membership in $\Theta$ refers to primitive inducing characters.
A finite-order background is absorbed into $\Theta$.
Every global norm twist is put in the smooth tests.  The remaining fixed
background $\omega_0$ is either $1$ or a genuine character modulo global
norm twists.  Its restriction to a compact local unit group has finite
conductor; its possibly infinite-order values on uniformizers and at
infinity remain part of the fixed character.

All tests have support in prescribed compact subintervals of
$(0,\infty)$.  We use Euler seminorms
\[
 p_J(W)=\max_{0\le j\le J}\sup_{x>0}|(x\partial_x)^jW(x)|.
\]
For a tuple of tests we use $1+\sum p_J(W)$.
Pure norm powers are normalized as $(q_{\mathfrak a}/X)^{it}$ at the
stated scale $X$.  Bounds below have the form of a fixed power of a finite
test seminorm times a fixed polynomial in the separated heights.
The permitted rowwise smooth choices are scale and parameter suprema
obtained by logarithmic Fourier inversion and Sobolev's inequality on a
fixed finite-dimensional parameter set.  Arithmetic column coefficients remain fixed independently of the row.

\subsection{Moment statements}

Let $\psi$ be a fixed representative oriented as in
Proposition~\ref{prop:fixed-type-theta}, and put
\begin{equation}\label{eq:can-presentations}
 \mathcal X_u=
 \{\theta\chi_u,\theta\bar\chi_u:\theta\in\Theta\}
 \ \cup\
 \{\psi\theta\bar\chi_u,\bar\psi\theta\chi_u:\theta\in\Theta\}.
\end{equation}
For $\lambda_u=\nu\chi_u^{\epsilon_\chi}$ in this family, define
\[
 M_u(Z^r;W)=Z^{-r/2}\sum_{\mathfrak n}
    \mu_K(\mathfrak n)\lambda_u(\mathfrak n)
    W(q_{\mathfrak n}/Z^r).
\]
Let $I$ be a fixed finite set of slots.  Their underlying prime lists
$\mathcal P_i$ are disjoint, lie outside $S$, and satisfy
$q_{\mathfrak p}\asymp Z^{z_i}$, where $z_i\ge0$.
For coefficients $|a_i(\mathfrak p)|\le1$, independent of $u$, set
\[
 Q_{u,i}=Z^{-z_i/2}\sum_{\mathfrak p\in\mathcal P_i}
    a_i(\mathfrak p)\lambda_u(\mathfrak p)
    W_i(q_{\mathfrak p}/Z^{z_i}),
 \qquad z=\sum_{i\in I}z_i.
\]
Empty products equal one.  The lists may contain all good primes in the annuli.

\begin{theorem}[Marked inverse moment]\label{thm:can-inverse}
Fix bounded nonnegative ranges for $m,r,z_i$, a bound for $|I|$, and
$c_1,c_2>0$.  If
\[
 r+2z\le m-c_1,\qquad 2r+8z\le3m-c_2,
\]
then, for every $\eps>0$,
\begin{equation}\label{eq:can-inverse}
 \sum_{0<q_u\ll Z^m}
 \left|M_u(Z^r;W)\prod_{i\in I}Q_{u,i}\right|^2
 \ll Z^{m+\eps}.
\end{equation}
The bound is uniform in the prime lists and their bounded coefficients,
and holds for every subcollection of the slots and either permitted
orientation.  No slot mesh is required.  The initial masks consist precisely of the fixed exclusions and
displayed row-character zeros.
The implied constant has the finite-seminorm and polynomial-height
dependence specified above.
\end{theorem}

For the ordinary moment a larger internal class is needed.  Write
\[
 \xi_u(\mathfrak n)=\vartheta(\mathfrak n)
      \prod_j\chi_{\mathfrak n}(b_j)^{e_j}\chi_u(\mathfrak n),
 \qquad \lambda_u=\omega_0\xi_u,
 \qquad M=m+q,
\]
where $\vartheta\in\Theta$ and the $b_j,e_j$ are fixed within the
current row sum.  The union of the supports of their displayed moving
factors, before cancellations and reduction modulo six, has norm at most
$Z^q$.  Each such prime is a common zero of every column.  A redundant
factor may be moved to a common radical mask only through an exact
factorization preserving its remaining phase.  An extra mask
$\mathbf1_{(\mathfrak n,\mathfrak R)=1}$ of polynomial norm is allowed
if it is fixed within the row sum and is used in both plain columns and
every slot.  Cancellations with the varying $u$ remain part of the natural row zeros.

For $X=Z^n$ let
\[
 S_{\lambda_u,\mathfrak R}(n;W)=X^{-1/2}
   \sum_{\mathfrak a}\lambda_u(\mathfrak a)
      \mathbf1_{(\mathfrak a,\mathfrak R)=1}W(q_{\mathfrak a}/X).
\]
The paired slot coefficient is
\begin{equation}\label{eq:can-pairing}
 a_i(\mathfrak p)=\omega_0(\mathfrak p)^{-1}
     \sum_j c_{ij}\vartheta_{ij}(\mathfrak p),
 \qquad \vartheta_{ij}\in\Theta.
\end{equation}
The finite lists and coefficients are fixed independently of $Z,u$.
Thus the actual character in each prime sum is finite order:
$\lambda_u a_i=\xi_u\sum_jc_{ij}\vartheta_{ij}$.
In this ordinary moment the slots are the complete smooth sums
\[
 Q_{u,i,\mathfrak R}
 =Z^{-z_i/2}\sum_{\mathfrak p\notin S}
   a_i(\mathfrak p)\lambda_u(\mathfrak p)
   \mathbf1_{(\mathfrak p,\mathfrak R)=1}
   W_i(q_{\mathfrak p}/Z^{z_i}),
\]
with disjoint underlying supports before the masks.
Fixed finite-ray restrictions are encoded in \eqref{eq:can-pairing}.
Let $\mathcal R_0$ consist of the rows for which $\lambda_u^*$ is
nonprincipal; for $z>0$ let $\mathcal R_z$ consist of those for which
$\xi_u^*\notin\Theta$.  We always also restrict $0<q_u\ll Z^m$.

\begin{theorem}[Paired fourth moment]\label{thm:can-fourth}
Let $3/4\le\kappa\le1$.  If $z>0$ and $\kappa<1$, suppose that every finite-order
Hecke $L$-function over $K$ has no zero with real part greater than
$(1+\kappa)/2$.  No zero-free premise is imposed when $\kappa=1$.
For every $\eps>0$ and prescribed bounded ranges of all real lengths,
there is $\eta>0$, uniform in $\kappa\in[3/4,1]$, such that a fixed slot
system with $z_i\le\eta$ satisfies
\begin{equation}\label{eq:can-fourth}
 \sum_{u\in\mathcal R_z}
 \left|S_{\lambda_u,\mathfrak R}(n_1;W_1)
       S_{\lambda_u,\mathfrak R}(n_2;W_2)
       \prod_iQ_{u,i,\mathfrak R}\right|^2
 \ll Z^{M+\eps}
\end{equation}
in the following ranges:
\begin{equation}\label{eq:can-fourth-capacity}
 \begin{cases}
 n_1,n_2\ge0 \text{ arbitrary in the prescribed bounded ranges},&z=0,\\
 n_1+n_2+6\kappa z\le M,&z>0.
 \end{cases}
\end{equation}
The mesh precedes the fixed number of slots.  Constants, finite
seminorm orders, and polynomial height degrees may depend on that
number and on the fixed arithmetic data, but are uniform in the moving
support, common masks, rows, and frozen outer labels.
\end{theorem}

Theorems~\ref{thm:can-inverse} and \ref{thm:can-fourth} are the
field-and-type versions of \cite[Lemmas~17.1 and~18.1]{Q}, respectively;
the internal inverse induction is \cite[Lemma~17.2]{Q}.
Their source statements concern the original finite-character setting.
Below we retain those canonical reductions and give the coefficient
identities, field-dependent estimates, and termination and loss
calculations needed to transfer them.  

\subsection{Finite coefficients and backgrounds}

At an exterior prime of norm $P$, define the normalized complete sum
$\mathcal G(\mathfrak p^a,h)$ using the original zero-extended
character of exponent $a$ and the primitive additive character.
Separating the higher residue digits gives, for $a\ge1$,
\begin{equation}\label{eq:can-local-G}
 \begin{aligned}
 |\mathcal G(\mathfrak p^a,h)|
   &=P^{(a-1)/2}\mathbf1_{v_{\mathfrak p}(h)=a-1},&&6\nmid a,\\
 \mathcal G(\mathfrak p^a,h)
   &=P^{a/2}\mathbf1_{\mathfrak p^a\mid h}
     -P^{a/2-1}\mathbf1_{\mathfrak p^{a-1}\mid h},&&6\mid a.
 \end{aligned}
\end{equation}
Indeed the sum over the $P^{a-1}$ lifts vanishes unless
$\mathfrak p^{a-1}\mid h$.  The remaining field sum is a nonprincipal
Gauss sum, of absolute value $\sqrt P$, or the Ramanujan sum
$P\mathbf1_{\mathfrak p\mid h/\mathfrak p^{a-1}}-1$.
This proves the formulas over every residue field, with all zeros intact.

The Chinese remainder theorem gives the complete two-column
correlation of \cite[Lemmas 13.3--13.4]{Q} with $P=q_{\mathfrak p}$.
Its diagonal is zero unless the two full moduli agree, and then equals
their Euler function.  At a common prime of multiplicities $a>b$ the
correlation vanishes unless $6\mid b$ and the divided frequency is a
unit; its remaining size is at most $P^b$.  At equal multiplicities
$a$ its size is at most $P^a$, with the stronger $P^{a-1}$ when
$6\nmid a$ and the divided frequency is a unit.  These are finite-field
identities: solving the two-column congruence leaves either a complete
nonprincipal sum or a sum over the units with one value omitted.
The finite factors at $S$ are kept in fixed sector tables.

To use the complete-support correlation, first write the full products
as $Da,Eb$, where $(a,b)=1$ and $(ab,DE)=1$; $D,E$ contain every
common prime with its full multiplicity.  After the genuine identity is
established, extend its right-hand side to individually allowed
residual pairs and insert
\begin{equation}\label{eq:can-full-mobius}
 \mathbf1_{(a,b)=1}=\sum_{t\mid(a,b)}\mu_K(t).
\end{equation}
  At every stage an extracted prime
retains its zero on all affected residual columns.

The calibrated squarefree identities from
Proposition~\ref{prop:arithmetic-calibration} are
\begin{equation}\label{eq:can-calibration}
 \gamma_2(\mathfrak n)^3=\mu_K(\mathfrak n)\Lambda(\mathfrak n),
 \qquad
 \bar\Lambda(\mathfrak n)\gamma_2(\mathfrak n)\gamma_1(\mathfrak n)
     =\mu_K(\mathfrak n)G(\mathfrak n),
\end{equation}
where $\Lambda$ is unitary and $G$ is finite order.
The canonical squarefree coefficient is therefore
$a_0(\mathfrak n)=\bar\Lambda(\mathfrak n)\gamma_2(\mathfrak n)$.
The first primitive transform changes $a_0\nu$ into
$\mu_K\nu\vartheta$, and the second changes it back into
$a_0\nu\vartheta\vartheta'$, with $\vartheta,\vartheta'\in\Theta$.
For example, $\mu_K\bar\gamma_1=a_0\bar G$; the remaining factors are
the fixed finite-sector coefficients just described.

The identification with the actual theta coefficient, including the
cube-completion masks, is made explicit in
Lemma~\ref{lem:canonical-theta-bridge} below.

In the first transform write $n_i=A_iCu_i$.  The two column backgrounds
satisfy
\[
 \nu(n_1)\overline{\nu(n_2)}
   =\nu(A_1)\overline{\nu(A_2)}\,
       \nu(u_1)\overline{\nu(u_2)},
\]
because $|\nu(C)|=1$.  The finite table is expanded while
$\nu(u_i)$ remains on its column.  The common factor in the second
transform cancels in the same way.  Cube completion contributes
$\nu(b)^3$ on the cube variable; it is a bounded coefficient on that variable, while the live residual
character remains $\nu$.
At initialization a presentation $\nu\chi_u^{\epsilon_\chi}$ is
conjugated in its entirety when $\epsilon_\chi=1$.  Its canonical
background is $\nu^{-\epsilon_\chi}$.  Thus
\eqref{eq:can-presentations} gives exactly the closed class
\begin{equation}\label{eq:can-background-class}
 \nu_{\rm can}\in\Theta\cup\psi\Theta.
\end{equation}
Finite twists preserve the chosen archimedean type and the nonzero
angular condition of Proposition~\ref{prop:fixed-type-theta}.

For the ordinary reduction each side keeps one background on its
whole residual product.  Complete extraction contributes only
frozen unit scalars.  Both finite transforms leave the original
individual slot coefficients unchanged, so
\eqref{eq:can-pairing} continues to cancel the live background in every
actual prime polynomial.  A change of row orientation conjugates both
plain columns and all slots, coefficients, and tests together.
After reflecting one plain column, it alone may be conjugated with
its test inside the absolute value of the product.  This restores its
original character without changing the squared norm.

\begin{lemma}[The canonical coefficient and cube completion]\label{lem:canonical-theta-bridge}
On squarefree exterior ideals the coefficient
$a_0=\bar\Lambda\gamma_2$ has the exact identities
\eqref{eq:can-calibration}.  Let $\tau$ be a completely multiplicative
unit-modulus character times any of the original completely
multiplicative coprimality indicators.  For a whole-index weight $d$
put
\[
 \mathscr T_{d,\tau}(X)=
  \sum_{\substack{c\ \mathrm{sf}\\b}}
     \frac{\gamma_2(c)\tau(cb^3)}{q_c^{1/2}q_b}
             d(cb^3)V(q_cq_b^3/X),\qquad
 d^{[h]}(A)=d(h^3A).
\]
Then, for a compactly supported test,
\begin{equation}\label{eq:canonical-cube-inversion}
 \sum_{c\ \mathrm{sf}}\frac{\gamma_2(c)\tau(c)}{q_c^{1/2}}
                       d(c)V(q_c/X)
 =\sum_h\frac{\mu_K(h)\tau(h)^3}{q_h}
          \mathscr T_{d^{[h]},\tau}(X/q_h^3).
\end{equation}
Here and below, a puncture is the indicator
$\rho(a)=\ind_{(a,r_\rho)=1}$, fixed in the current row--label sum.
For a fixed squarefree label $f$ and
$\tau(a)=\bar\Lambda(a)\nu(a)\rho(a)\chi_a(k)\chi_a(f)^4$,
the completed sum for one mark tuple is exactly
\eqref{eq:marked-primal-exact}, with period
$\vartheta=\bar\Lambda\nu$, row $m=kf^4$, and puncture
$\rad(r_\rho)$.  Product-form sums of tuples follow by linearity.
Fixed finite factors can be absorbed into $\nu$.
\end{lemma}
\begin{proof}
In the cube of the CRT law \eqref{eq:gauss-crt} with $j=2$, both
cross exponents are divisible by six.  The cross exponents in the
product of its $j=1$ and $j=2$ laws are also divisible by six.
The two squarefree expressions are therefore multiplicative on coprime
ideals.  Their values at primes, including the minus signs, are
\eqref{eq:gauss-calibration}; this proves \eqref{eq:can-calibration}.
The finite order of $G$ follows from the Weil-index calculation on $B$.

Expand the right-hand side of \eqref{eq:canonical-cube-inversion}
and put $t=hb$.  Every term has common factor
$\gamma_2(c)\tau(ct^3)d(ct^3)/(q_c^{1/2}q_t)$ and common test
$V(q_cq_t^3/X)$.  Its remaining coefficient is
$\sum_{h\mid t}\mu_K(h)$, which is zero unless $t=1$.
All sums are finite for compactly supported $V$.  For the bounded
product-form mark weights used here, the Gaussian case also follows
by absolute convergence with the weighted divisor majorant.  Complete multiplicativity preserves the puncture indicators.

After cube extraction a product-form mark has argument $h^3A$,
with the assignment rule
\begin{equation}\label{eq:mark-cube-priority}
 \ind_{p\mid h^3A}=\ind_{p\mid h}
                     +\ind_{p\nmid h}\ind_{p\mid A}.
\end{equation}
Expanding this identity for the individual slots leaves either a mark
already assigned to $h$, or one live whole-index mark with its original
coefficient and its original exclusion at $h$.
Finally $\tau(cb^3)=\bar\Lambda(c)\nu(c)\rho(c)\chi_c(k)\chi_c(f)^4$
times the cube of the same multiplicative factors at $b$.
This is the coefficient in \eqref{eq:can-completed-marked}.
To match a centrally normalized squarefree polynomial, take
$V_*(y)=y^{1/2}W(y)$: then
$q_c^{-1/2}V_*(q_c/X)=X^{-1/2}W(q_c/X)$.
This identifies the test and the normalization as well as the
arithmetic coefficient in the canonical terminal.

After $h$ is frozen, its extracted phase factors as a row factor
$\chi_h(k)^3$, a label factor $\chi_h(f)^{12}$, and a scalar
independent of both.  Each has absolute value at most one;
$\chi_h(f)^{12}$ is still its coprimality indicator.
Multiplication by their product is a contraction in the joint
row--label Hilbert space.  Hence the source's Minkowski step may freeze $h$ before applying
the completed estimate.
The factor $\nu(h)^3$ in the inversion is on this extracted cube
variable; the live theta period remains $\bar\Lambda\nu$.
\end{proof}

\begin{lemma}[Preservation of paired whole-product phases]\label{lem:paired-phase-invariant}
In the ordinary canonical reduction, suppose the live plain character
is $\omega_0\xi$, with $\xi$ finite, and each live slot has relative
coefficient $\omega_0^{-1}\vartheta_i$, $\vartheta_i\in\Theta$.
Retain the source's whole-product allocation and common masks.
The two finite transforms, extraction, and allowed conjugations
preserve this property, up to replacing $\xi$ and $\vartheta_i$ by
finite twists.  Thus only a finite-order character is evaluated in
the prime cancellation premise.
\end{lemma}
\begin{proof}
Write a separated whole index as $A=DA'$.  The genuine factor is
exactly $\omega_0(D)\omega_0(A')$.  Complete extraction freezes the
first factor; a common factor in a coefficient and its conjugate
cancels by $|\omega_0(C)|^2=1$.  The complete-support finite Fourier
tables of \cite[Lemma 18.2]{Q} involve only $\xi$ and the residue
symbols.  They therefore multiply the surviving $\omega_0(A')$ by a finite character.  

On a live slot the two genuine factors are
$\omega_0(p)$ from the whole product and $\omega_0(p)^{-1}$ from
its specified coefficient.  They cancel before any prime sum is
estimated.  If that slot is assigned to an extracted factor, both
factors belong to the same frozen allocation; the remaining slots
retain their individual coefficients.  This remains true when an
index occurs in both rectangles of a centered expression: the
allocation is made on the common whole product, before their
smooth profiles are separated.

Conjugating a row orientation conjugates both plain factors, every
slot, and all of their coefficients and tests.  It replaces the
pair $\omega_0,\omega_0^{-1}$ by
$\bar\omega_0,\bar\omega_0^{-1}$.
Ordinary reflection changes one plain character to its conjugate;
conjugating that one reflected sum and its test inside the product
absolute square restores the original orientation.  Use this operation for an uncentered product or a separately
estimated comparison product. A centered difference retains its
common profiles, whole-product character, and residue cancellation
until the two rectangles are split.  The conductor
and common-mask identities are still applied to the original
presentation before primitive reduction.  These verifications prove
the invariant used in the source's canonical reductions.
\end{proof}

\subsection{The two sieve inputs}\label{sec:sieve-identification}

Heath-Brown's quadratic and cubic power-residue large sieves
\cite{HBQ,HBC} are the classical predecessors of the two bounds used
here.  The number-field inputs are the quadratic-family theorem of
Goldmakher--Louvel \cite[Theorem~1.1]{GL} and the $n$-th-order family
theorem of Blomer--Goldmakher--Louvel \cite[Theorem~1.3]{BGL}, specialized
to $n=3$.  
The quadratic bound for exterior allowed symbols is the identification
in \cite[Lemma~6.4]{P}; we check its family hypotheses over $K$ below.
In the present notation it says
\begin{equation}\label{eq:exterior-quadratic-sieve}
 \sum_{\substack{k\ \mathrm{sf}\\q_k\le K_0}}
       \left|\sum_{\substack{j\ \mathrm{sf}\\q_j\le N_0}}
                         z_j\chi_k(j)^3\right|^2
 \ll_{K,S,\eps}(K_0N_0)^\eps(K_0+N_0)\sum_j|z_j|^2.
\end{equation}
Here the family hypotheses can be checked directly.  Let $\eta_j$
be the primitive quadratic character of $K(\sqrt j)/K$, with $j$
represented by an allowed generator.  Reciprocity identifies its
exterior values with $\chi_{\bullet}(j)^3$.  For squarefree $j$,
every prime dividing $j$ is tamely ramified and every other exterior
prime is unramified, so its conductor is $j\mathfrak s_j$ with
$\mathfrak s_j$ supported on $S$.  The finitely many local square
classes at $S$ give a fixed common conductor bound there.
The class map from exterior ideals to the square-class image of
$B$ is a homomorphism, and the reciprocity defect on allowed
generators is one.  If coprime squarefree $j_1,j_2$ have the same
class, their quotient is a square at $S$.  Thus
$\eta_{j_1}\bar\eta_{j_2}$ is unramified at $S$ and has exact
conductor $j_1j_2$ outside $S$.  These are the quadratic-family
hypotheses of \cite[Definition 1 and Theorem 1.1]{GL}, and prove
\eqref{eq:exterior-quadratic-sieve}.  This is also the verification in \cite[Lemma 6.4]{P}.
All characters here have trivial infinity type.  Fixed row
restrictions and unimodular row or column factors are harmless.
Moving pair-dependent restrictions are treated by their exact
M\"obius identities below.

\begin{lemma}[Identification of the cubic sieve family]\label{lem:exterior-cubic-sieve}
For coefficients $c_P$ on squarefree exterior ideals $q_P\le V_0$,
\begin{equation}\label{eq:exterior-cubic-sieve}
 \sum_{\substack{j\ \mathrm{sf}\\q_j\le U_0}}
       \left|\sum_{P\ \mathrm{sf}}c_P\chi_P(j)^{\pm2}\right|^2
 \ll_{K,S,\eps}(U_0V_0)^\eps
       \{U_0+V_0+(U_0V_0)^{2/3}\}\sum_P|c_P|^2.
\end{equation}
Both variables are prime to $S$, and the symbols retain their zeros
at a common prime.
\end{lemma}
\begin{proof}
We apply \cite[Theorem 1.3]{BGL} with $n=3$ to the Fisher--Friedberg
family described in \cite[Section~2, Lemma~2.1 and Appendix~II]{BGL}.
Write $\chi_j^{\rm BGL}$ for its primitive character.  For a fixed modulus $\mathfrak c$ supported on $S$, its
reciprocity defect depends only on the two classes in
$H_{\mathfrak c}\otimes\ZZ/3$.  It has
$\chi_{(b)}^{\rm BGL}=(b/\cdot)_3$ when $b\equiv1\pmod{\mathfrak c}$,
and $\chi_{(1)}^{\rm BGL}=1$.  For squarefree $j$ its conductor lies
between $j$ and $\mathfrak c j$.  

Choose a modulus $\mathfrak m$ supported on $S$, divisible by
$\mathfrak c$, so deep that
$1+\mathfrak m\mathcal O_v\subset K_v^{\times6}$ at every finite
$v\in S$.  Let $H_{\mathfrak m}$ be its ray class group.  For a fixed
squarefree $j$ compare, on ideals $b$ prime to $\mathfrak m j$,
\[
 \Xi_j(b)=\chi_b(j)^2,
       \qquad\Xi_j^{\rm BGL}(b)=\chi_j^{\rm BGL}(b).
\]
Both functions are multiplicative.  If $(b)$ is ray-principal modulo
$\mathfrak m$, then $[b]=1$ in $K_S^\times/K_S^{\times6}$, so $b$
itself is an allowed generator.  Reciprocity gives
$\Xi_j((b))=\chi_j(b)^2=(b/j)_3$.
For the BGL family the reciprocity defect with a trivial ray class is
one: use its dependence on classes and the normalization at $(1)$.
Hence
\[
 \Xi_j^{\rm BGL}((b))
       =\chi_{(b)}^{\rm BGL}(j)=(b/j)_3.
\]
It follows that $\xi_j=\Xi_j/\Xi_j^{\rm BGL}$ descends to a
character of $H_{\mathfrak m}$.  More explicitly, every ray class has
a representative prime to $j$ by approximation; the kernel of this
restricted ideal-to-class map consists exactly of its ray-principal
ideals.  

Split the columns $P$ into their finitely many classes
$h\in H_{\mathfrak m}$.  On one such class,
\[
 \chi_P(j)^2=\xi_j(h)\chi_j^{\rm BGL}(P).
\]
The factor $\xi_j(h)$ is a unimodular row factor and disappears inside
the absolute square.  Cauchy over $h$, followed by the BGL theorem
with coefficients $c_P\ind_{[P]=h}$, proves the plus sign in
\eqref{eq:exterior-cubic-sieve}.  At a common prime both sides are
zero, because the BGL conductor contains squarefree $j$.
Complex conjugation proves the minus sign.  All losses from the
class splitting are constants for the fixed $K,S$.
\end{proof}

\begin{lemma}[Completed quadratic reduction]\label{lem:completed-quadratic-reduction}
Let $k,n$ be squarefree exterior ideals with $q_k\ll K_0$,
$q_n\asymp N_0$, and let $q_b\asymp B_0$, with $b$ unrestricted.
For $\beta(n,b)$ independent of $k$, the positive reduction in
\cite[Lemma 5.5]{Q} holds over $K$ with the original zeros.
In particular, if $|\beta|\le1$, then
\begin{equation}\label{eq:completed-quadratic-simple}
 \sum_k\left|\sum_{n,b}\beta(n,b)\chi_k(nb)^3\right|^2
 \ll_\eps(K_0N_0B_0)^\eps(K_0+N_0B_0)N_0B_0.
\end{equation}
\end{lemma}
\begin{proof}
Here is the positive form that will also be used with an arbitrary
coefficient.  Write uniquely $b=gt^2$, with $g$ squarefree, and put
$c=(n,g)$, $n=cm$, $g=ch$.  Then $c,m,h$ are pairwise coprime and squarefree, while $g,t$
may share factors.  Exactly,
\[
 \chi_k(nb)^3=\chi_k(mh)^3\ind_{(k,ct)=1}.
\]
Keep the $k$--$t$ restriction and expand the full $k$--$c$ mask by
$\sum_{r\mid(k,c)}\mu_K(r)$.  Rowwise divisor Cauchy costs a
small power.  Put $k=rk'$, $c=rc'$, and separate dyads
$q_c\asymp C$, $q_g\asymp G$, $q_t\asymp T$, so $GT^2\asymp B_0$.
Cauchy over $t$ costs $O(T)$, and Cauchy over $c'$ costs $O(C/q_r)$.
Group the squarefree product $mh$, whose norm is $O(N_0G/C^2)$;
its factorizations have divisor-bounded multiplicity.  Applying
\eqref{eq:exterior-quadratic-sieve} gives, up to a small power,
\begin{equation}\label{eq:completed-quadratic-positive}
 T\sum_{q_t\asymp T}\ \sum_{\substack{r\ \mathrm{sf}\\q_r\ll\min(K_0,C)}}
 \left(\frac{K_0}{q_r}+\frac{N_0G}{C^2}\right)\frac{C}{q_r}
 \sum_{c',m,h}|\beta(rc'm,rc'ht^2)|^2.
\end{equation}
The last sum keeps its inherited support. The unit factor at $r$
is bounded after taking the square; the row--$c'$ mask has been
resolved by the divisor sum.  This proves the positive
reduction for arbitrary $\beta$.  Counting its remaining ideal ranges
when $|\beta|\le1$ and summing dyads gives
\eqref{eq:completed-quadratic-simple}.  All factorizations take place among exterior ideals.
\end{proof}

\subsection{The marked quadratic--cubic terminal}

\begin{lemma}[Quadratic--cubic norm]\label{lem:can-hybrid}
Let $K_0,N_0,B_0,L_0\ge1$.  Let $k,n,P$ be squarefree exterior ideals,
$q_k\ll K_0$, $q_n\asymp N_0$, $q_P\asymp L_0$, and let $b$ be an
exterior ideal with $q_b\asymp B_0$.  If $|a(P)|,|\beta(n,b)|\le1$,
with $a$ independent of $k,n,b$ and $\beta$ independent of $k,P$, then
\begin{align}\label{eq:can-hybrid}
 &\sum_k\left|
   \sum_P\frac{a(P)}{\sqrt{q_P}}\mathbf1_{(P,k)=1}
   \sum_{n,b}\beta(n,b)\chi_k(nb)^3\chi_P(n)^{-2}
                       \mathbf1_{(P,b)=1}\right|^2\notag\\
 &\hspace{12mm}\ll_{\eps}(K_0N_0B_0L_0)^\eps
       (K_0+N_0B_0)B_0
       \{N_0+L_0+(N_0L_0)^{2/3}\}.
\end{align}
Fixed row restrictions and fixed coefficient-support restrictions are
allowed if these independence conditions are preserved.
\end{lemma}
\begin{proof}
Use \eqref{eq:exterior-quadratic-sieve} and
Lemma~\ref{lem:exterior-cubic-sieve}.  The positive completed
quadratic reduction for arbitrary coefficients is
\eqref{eq:completed-quadratic-positive}.

We spell out the two moving masks in applying these sieves.
Expand the first by $\sum_{D\mid(P,k)}\mu_K(D)$, then use divisor
Cauchy and put $P=DP'$, $k=Dk'$.  The factors at $D$ join $\beta$ or the fixed $P'$ coefficient,
with the $(P,k)$ mask resolved by the divisor sum.
For fixed $D$ apply the completed quadratic reduction, writing
$b=gt^2$, $c=(n,g)$, $n=cm$, $g=ch$, and resolving the full row--$c$
mask by its divisor $r$.  Let $G,C,T$ be the dyadic norm scales of
$g,c,t$, respectively.  Squarefreeness gives pairwise coprime $c,m,h$; the factors $g,t$
may overlap.
The remaining cubic column is $j=c'm=n/r$, and the exact identity
\[
 \chi_{P'}(rc'm)^{-2}\mathbf1_{(P',rc'ht^2)=1}
 =\chi_{P'}(r)^{-2}\chi_{P'}(c'm)^{-2}
       \mathbf1_{(P',ht)=1}
\]
preserves all missing zeros.  Its $P'$ coefficient is independent of
$j$ and has squared norm $O_\eps((K_0N_0B_0L_0)^\eps)$ after division
by $q_{P'}$.  On a block $q_D\asymp D_0$, $q_r\asymp r_0$, the cubic
sieve therefore contributes
\[
 \frac GC\left\{\frac{N_0}{r_0}+\frac{L_0}{D_0}
       +\left(\frac{N_0L_0}{r_0D_0}\right)^{2/3}\right\}
\]
up to the chosen small power.  Here $G/C$ includes the $h$ count.
The remaining quadratic and $t$ factors are $T^2C/r_0$; their
product with $G/C$ is $B_0/r_0$, since $T^2G=B_0$.
There are $O(r_0)$ possible $r$
and $\sum_{D\asymp D_0}q_D^{-1}\ll1$ up to that power.
The quadratic factor is bounded by $K_0+N_0B_0$.
Summing the divisor and dyadic ranges proves \eqref{eq:can-hybrid}.
This is the factorization proof of \cite[Lemma 14.2]{Q}; the argument
just given verifies its two sieve inputs and its masks over $K$.
\end{proof}

For a product-form mark write
\begin{equation}\label{eq:can-mark}
 d_I(A)=\sum_{(\mathfrak p_i)\in\prod_{i\in I}\mathcal P_i}
           \prod_{i\in I}a_i(\mathfrak p_i)\mathbf1_{\mathfrak p_i\mid A}.
\end{equation}
The whole-index marked expansion is Theorem~\ref{thm:marked-expansion}.
Its raw local table, mixed phases, and coefficient independence will all
be retained in the following norm calculation.

\begin{lemma}[Marked reflected energy]\label{lem:can-marked-energy}
Put $\Psi_{k,f}(a)=\nu(a)\rho(a)\chi_a(k)\chi_a(f)^4$.
The completed marked sum is
\begin{equation}\label{eq:can-completed-marked}
 \mathcal T_d(X;k,f)=
 \sum_{\substack{n\ \mathrm{sf}\\ b}}
 \frac{a_0(n)\Psi_{k,f}(n)\,
       \bar\Lambda(b)^3\Psi_{k,f}(b)^3}{q_n^{1/2}q_b}
       d(nb^3)V(q_nq_b^3/X),
\end{equation}
where both ideals avoid $S$ but may share primes, $V$ is a fixed
annular profile, and $d$ has the product form \eqref{eq:can-mark}.
Let the row ball have norm scale $Z^M$ and let the squarefree label
$f$ have scale $Z^V$.  Write $Q=\log_Zq_{\rad(r_\rho)}$ for the
fixed puncture length and $Z^{N_*}$ for the completed scale.
Freeze the powerful row part of length $O$, and retain the actual
residual squarefree row annulus of length $H$.
Let $A_0$ be the active nonresidual radical length; $S_0$ the length
of its small Ramanujan and active zero-mask terms; and $N_0,B_0$ the
lengths assigned, respectively, to the dual squarefree variable and
to the cube variable alone.  Let $z_a$ be the active slot length.
On a dual block of squarefree and cube lengths $v,\ell_b$, put
\begin{align}\label{eq:can-energy-data}
 T_d&=2H+2A_0+2z_a-N_*-N_0-3B_0,\notag\\
 u_0&=\min\{v,z_a,(v+z_a)/3\},\\
 E_{\rm ref}&=\frac O2+\max(H,v+\ell_b)-S_0-B_0+z_a-u_0-\ell_b
                  -\frac12(T_d-v-3\ell_b)_+ .\label{eq:can-Eref}
\end{align}
After summing the frozen powerful parts, the squared row norm of this
block is $\ll Z^{E_{\rm ref}+\eps}$, with finite seminorm and fixed
polynomial height factors.  For every fixed $\tau_{\rm ref}>0$,
blocks with $v+3\ell_b>T_d+\tau_{\rm ref}$ have arbitrarily small
total size after a fixed further kernel shift.  The conductor
comparisons on the actual annuli are
\begin{equation}\label{eq:can-conductor-comparison}
 H\le M-O+O(1/\log Z),\qquad
 T_d-S_0-B_0\le2M-O+Q+V-N_*+2z_a+O(1/\log Z).
\end{equation}
The constants in these $O$-terms depend only on the fixed support
intervals and arithmetic data.
\end{lemma}
\begin{proof}
Lemma~\ref{lem:canonical-theta-bridge} identifies the coefficient,
including its unrestricted cube and all of its masks, with
Theorem~\ref{thm:marked-expansion}.  We give the resulting norm
calculation, keeping track of its coefficient independence.

\emph{Rows, collisions, and inactive slots.}
Write the exterior ideal of a row as
$k=k_{\rm pow}k_{\rm sup}d$.  Here $k_{\rm pow}$ contains the full
prime powers of exponent at least two, $k_{\rm sup}$ contains the
valuation-one primes supported on $f\rad(r_\rho)$, and $d$ contains
the remaining valuation-one primes.  Fix the first two parts and the
unit class of the row, and keep the actual annulus $q_d\asymp Z^H$.
To avoid conflict with the canonical label $f$, denote the active
ordinary-prime product in Theorem~\ref{thm:marked-expansion} by $f_0$
in this proof.

A tuple meeting the row, label, or puncture is zero on the primal
side.  Remove it before reflecting.  After fixing all inactive
primes, their exclusions from $d$ are a fixed row restriction, and
Minkowski in the row Hilbert space removes their sums at cost
$\prod_i\sum_{p\in\mathcal P_i}q_p^{-1}\ll Z^{\eps_1}$.
An inactive prime contributes its scalar weight and remains allowed
in the dual index.  Active collisions
are the explicit mask $\ind_{(P,d)=1}$ in
\eqref{eq:marked-reflection-exact}.

\emph{Frozen valuations and the full cube.}
For $p\mid f_0$ write uniquely
\[
 e_p=e_p^{\min}+3x_p,\qquad
 e_p^{\min}=\begin{cases}
 0,&e_p=0,\\1,&e_p\equiv1\pmod3,\\3,&e_p\ge3,	e_p\equiv0\pmod3.
 \end{cases}
\]
Thus $a'=a_{\min}x^3$ with $x$ supported on $f_0$.
The local table gives
\[
 |s_{p,4,0}|=q_p^{-1/2},\quad
 |s_{p,4,1}|\le1,\quad |s_{p,4,3}|\le q_p^{-1/2},\quad
 |s_{p,0,0}|=q_p^{-1/2}.
\]
These are precisely the factors counted by $S_0,N_0,B_0$.
For the remaining valuations the coefficient is bounded by the
corresponding minimum-valuation coefficient times $q_x^{-1}$;
the exponents $b_{j,e}$ are unchanged on the allowed positive
progressions modulo three.  Keep the displayed powers before using
any positive majorant.

Put $a_{\min}$ outside the norm by Minkowski, but keep $x$ inside,
with the dual cube variable
\begin{equation}\label{eq:energy-full-cube}
 b=xn',\qquad (n',f_0)=1.
\end{equation}
This decomposition is unique.  In the row factor
$\chi_d(a')=\chi_d(a_{\min})\chi_d(x)^3$, giving the joint symbol
$\chi_d(c'b)^3$.  In the mark and squarefree column factors the
extra exponent is six, so their $x$-dependence is one, because
$P,c'$ are prime to $f_0$.  The remaining $L(x)^{-3}$ and local
coefficients are functions of the dual variables alone.
The kernel argument is $Xq_{a_{\min}}q_{c'}q_b^3/(q_dq_{f_0}q_P)^2$.
The normalized product of the finitely many local choices has
mass at most $C^{\omega(k_{\rm pow}k_{\rm sup}f r_\rho)}$ after
the assigned powers are extracted; a subsidiary divisor bound
absorbs it into $Z^{\eps_1}$.  Equivalently one may sum the frozen
majorants using \eqref{eq:marked-dirichlet-majorant}.
The per-tuple coefficient exponent on a dual block is consequently
\begin{equation}\label{eq:energy-coefficient-size}
 -v/2-\ell_b-(S_0+B_0+z_a)/2.
\end{equation}

\emph{A common profile on the actual annuli.}
Fix the finitely many classes of $d,c',P,b$ and the frequency unit.
Then the selected $\kappa$ in \eqref{eq:marked-frequency} is fixed.
Writing each varying norm as its block scale times a normalized
coordinate, the argument of the kernel is
\[
 Yy_{c'}y_b^3y_d^{-2}\prod_{i\ \mathrm{active}}y_i^{-2},
       \qquad Y\asymp Z^{v+3\ell_b-T_d}.
\]
All coordinates lie in fixed compact intervals.  Apply the common
profile separation of \cite[Lemma 5.3]{Q}, using
\eqref{eq:theta-kernel-bound}.  Every weighted logarithmic Fourier
norm is bounded by finitely many test seminorms times
\begin{equation}\label{eq:energy-common-kernel-scalar}
 m_{\rm ker}=Z^{-(T_d-v-3\ell_b)_+/4}
\end{equation}
(up to fixed support factors).  Extract this scalar from the whole
linear row vector before the norm inequality or any row extension.
The Fourier measure is common to $d,P,c',b$.

\emph{The hybrid norm.}
At a separated Fourier mode, after the frozen coefficients are
extracted, the remaining row vector is a unimodular row factor times
\[
 \sum_P\frac{a(P)}{q_P^{1/2}}\ind_{(P,d)=1}
       \sum_{c',b}\beta(c',b)\chi_d(c'b)^3\chi_P(c')^{-2}
                                          \ind_{(P,b)=1}.
\]
The factor $a(P)$ consists of \eqref{eq:marked-pi}, the original
individual slot coefficients, and slot norm phases.
The factor $\beta(c',b)$ consists of \eqref{eq:marked-lambda}, the
remaining cube coefficients, dual norm phases, and the class slices.
After \eqref{eq:energy-full-cube}, these have exactly the two
independences in Lemma~\ref{lem:can-hybrid}: $a$ is independent of
$d,c',b$ and $\beta$ is independent of $d,P$.  The mask at $P,b$
equals that at $P,n'$, since $P$ is prime to $f_0$.
Tuple aggregation costs at most a divisor factor.

Apply that lemma with $(K_0,N_0,B_0,L_0)=(Z^H,Z^v,Z^{\ell_b},Z^{z_a})$.
The factor $q_P^{-1/2}$ is already inside the hybrid lemma.
The squared coefficient outside its norm is therefore
$Z^{-v-2\ell_b-S_0-B_0}$.  The hybrid estimate contributes exponent
\[
 \max(H,v+\ell_b)+\ell_b+
              \max\{v,z_a,2(v+z_a)/3\}.
\]
Since the last maximum is $v+z_a-u_0$, the total, before the
kernel, is
$\max(H,v+\ell_b)-S_0-B_0+z_a-u_0-\ell_b$.
Squaring \eqref{eq:energy-common-kernel-scalar} gives the last term
of \eqref{eq:can-Eref}.

\emph{Tails.}
For every fixed $\sigma>1/2$, use
\eqref{eq:marked-dirichlet-majorant} on the genuine reflected series
before Fourier absolutization.  Outside the retained range the
kernel argument exceeds $Z^{\tau_{\rm ref}/2}$ once fixed support
constants have been absorbed in the threshold.  At large-argument
order $A$ its total absolute contribution is at most
\[
 C_A Z^{B_{\rm tail}} Z^{-(A-\sigma)\tau_{\rm ref}/2}.
\]
The exponent $B_{\rm tail}$ accounts for the bounded row, tuple,
and local length ranges and is independent of $A$; the constant and
seminorm order may depend on $A$.  Choosing $A$ large proves the
arbitrary power tail assertion.  A norm twist changes only a fixed
polynomial height factor in these finitely many seminorms.

\emph{Frozen row count and conductor comparison.}
A powerful exterior ideal is uniquely $x_0^2y_0^3$, with $y_0$
squarefree.  Ideal counting and convergence of
$\sum q_{y_0}^{-3/2}$ give $O(Z^{O/2})$ choices in its $O$-dyad.
The supported part divides $\rad(f r_\rho)$ and has $Z^{\eps_1}$
choices.  The unit classes are finite.  Their positive sum adds
$O/2+\eps_1$ once and proves the stated row energy.

For completeness, the contribution to
$2A_0-N_0-S_0-4B_0$ is $2$ for ordinary digits $j=1,2,3,5$;
$1$ for active $j=0$, for the $e=0$ part of $j=4$, or for its
$e=1$ part; $-2$ for its $e=3$ part; and $0$ for an inactive $j=0$.
A residual row prime contributes $2$ through $2H$.
Prime by prime these contributions are bounded by the valuations of
$(k^2/k_{\rm pow})f\rad(r_\rho)$.
Taking norms proves \eqref{eq:can-conductor-comparison}, since
$q_k\ll Z^M$ and $q_{k_{\rm pow}}\asymp Z^O$.
The $S$-variables are contained in the sheet integral
\eqref{eq:marked-H-explicit}.
\end{proof}

\begin{corollary}[Unmarked completed moment]\label{cor:unmarked-completed-moment}
For an empty slot collection and fixed powerful-part dyad $O$,
\begin{equation}\label{eq:unmarked-completed-moment}
 \sum_{0<q_k\ll Z^M}^{\prime}
        |\mathcal T_1(Z^{N_*};k,f)|^2
 \ll Z^{O/2+\max\{M-O,\,2M-O+Q+V-N_*\}+\eps}.
\end{equation}
The prime restricts the powerful part to that dyad.  With slots present
one retains $z_a,u_0$ in \eqref{eq:can-Eref} instead.
\end{corollary}
\begin{proof}
Here $z_a=u_0=0$.  On a retained block,
$v+3\ell_b\le T_d+\tau_{\rm ref}$, so directly
\[
 E_{\rm ref}\le O/2+\max\{H,T_d-S_0-B_0\}+\tau_{\rm ref}.
\]
Indeed the $H$ branch only subtracts nonnegative terms, and the
$v+\ell_b$ branch is at most $T_d-S_0-B_0+\tau_{\rm ref}$.
Insert \eqref{eq:can-conductor-comparison} and choose
$\tau_{\rm ref}$ and the subsidiary dyadic losses within $\eps$.
Sum the logarithmically many retained dyads; the tails were proved
above.  This gives \eqref{eq:unmarked-completed-moment}.
\end{proof}

\subsection{Canonical inverse induction}\label{subsec:inverse-induction}

The inverse moment is proved through a slightly larger second-moment
family. Besides the row $k$, it keeps a squarefree label $f$ and a
puncture $r_\rho$. This is the form that is recovered after two Poisson
transforms: the first produces inverse coefficients and the second
returns the theta-calibrated coefficient $a_0$. The quantity $F=N+V$
is the combined column-and-label length. The two margins below control,
respectively, the puncture budget and the marked reflected terminal.

\begin{proposition}[Canonical energy]\label{prop:can-canonical}
Fix bounded nonnegative ranges for $M,N,V,z_0$ and a positive margin
$c_*$.  Let $f$ run over squarefree exterior ideals of norm
$\asymp Z^V$, with a nonnegative weight $w(f)\le C\tau_K(f)^C$
depending only on $f$.  Let $\nu$ belong to
\eqref{eq:can-background-class}; let
$\rho(n)=\mathbf1_{(n,r_\rho)=1}$ be fixed within the row and label
sum; and let $d$ have the product form \eqref{eq:can-mark}, with
total length at most $z_0$ and coefficients independent of both rows
and labels.  There is no additional residual coefficient on $n$.
Put $F=N+V$.  If
\begin{equation}\label{eq:can-invariants}
 F-M-Q-z_0\ge c_*,\qquad 4F-3M-6z_0\ge c_*,
 \qquad Q=\log_Zq_{\rad(r_\rho)},
\end{equation}
then
\begin{align}\label{eq:can-energy}
 \mathcal E={}&Z^{-V}\sum_f w(f)\sum_{0<q_k\ll Z^M}
 \left|Z^{-N/2}\sum_{n\ \mathrm{sf}}a_0(n)\nu(n)\rho(n)
      \chi_n(k)\chi_n(f)^4d(n)W(q_n/Z^N)\right|^2\notag\\
 &\ll Z^{F+\eps}.
\end{align}
The statement is simultaneous for all subcollections of the original
slots with the same cap $z_0$, and is uniform in the frozen outer
ideals within the stated ranges.
\end{proposition}
\begin{proof}
We induct on the logarithmic row length $M$. A node consists of the
energy \eqref{eq:can-energy}, its lengths $(M,N,V)$, the fixed puncture,
and the surviving subset of the original marks. Its inductive data are
the two margins in \eqref{eq:can-invariants}; its coefficient class is
\eqref{eq:can-background-class}. The two-transform reduction of
\cite[Section~17.5]{Q} has the following form over $K$:
\[
 \mathcal E
 \ \xrightarrow{\text{first Poisson}}\
 \text{inverse-column norms}
 \ \xrightarrow{\text{second Poisson}}\
 \mathcal E_{\rm child}\quad(M'<M).
\]
The identities below specify the coefficients, masks, and norm counts
at each arrow. Choose a fixed step $0<d_0\le c_*/200$.

\smallskip\noindent\emph{Terminal blocks.}
Cube completion expresses a short block through the marked reflected
sum at scale $N_*=N-3b_h$, where the extracted cube variable has
length $b_h<d_0+O(1/\log Z)$.
If $V<d_0$, Lemma~\ref{lem:can-marked-energy} and
\eqref{eq:can-invariants} give the three possible terminal exponents
\begin{equation}\label{eq:can-terminal-ledger}
 F-c_{\rm node},\qquad
 F-2c_{\rm node}+5d_0,\qquad
 F-M/4-3c_{\rm node}/4+5d_0/2,
\end{equation}
apart from prescribed localization and small-power losses.
They correspond to the row branch, the column branch $u_0=z_a$,
and the other column branch.  For example the last uses
$u_0\ge v/2$, the retained dual inequality, and the second invariant.
All are strictly smaller than $F$ when $c_{\rm node}\ge c_*/2$.
The label sum has weighted mass $\ll Z^{V+\eps}$ and cancels the
outside $Z^{-V}$ exactly once.

\smallskip\noindent\emph{First transform.}
For the remaining blocks, reopen the cube completion if necessary.
With cube length $\ell$ and squarefree column length $r=N-3\ell$,
one has $V+\ell\ge d_0$.  Define
\[
 \mathcal T_{k,f}
 =\sum_{b\asymp Z^\ell}\beta(b,f)\chi_b(k)^3
   \sum_{n\ \mathrm{sf}}a_0(n)\nu(n)\rho(n)
       \chi_n(k)\chi_n(f)^4d(nb^3)W(q_n/Z^r).
\]
The block is
\[
 Z^{-r-2\ell-V}\sum_f w(f)
       \sum_k\Phi(q_k/Z^M)|\mathcal T_{k,f}|^2,
\]
where $\beta$ is divisor-bounded and independent of $k,n$.
The original cube masks, including $(b,f)=1$, remain in it.
The cube and squarefree column may share primes.

Apply Lemma~\ref{lem:weighted-orbital-poisson} in the row after
expanding the square.  Write
$n_i=A_iCu_i$, with $A_i=(n_i,\rad(b_1b_2))$ and $C$ the remaining
common factor.  On the genuine coprime residual locus the primitive
row conductor is $u_1u_2R_1$, and the remaining zero mask is
$C\rad(b_1b_2)/R_1$.  The local digit $t_p$ defining $R_1$ is
\[
 t_p=\mathbf1_{p\mid A_1}-\mathbf1_{p\mid A_2}
                  +3v_p(b_1b_2)\pmod6.
\]
The primitive Poisson factor is the row scale divided by the mask
divisor and the square root of this conductor norm.  Formula
\eqref{eq:can-calibration} gives the inverse-type residual
coefficient $\mu_K\nu\vartheta$.
The cube table is particularly useful: write
$b_1b_2=q_0^2J_2^2s$, with $s,J_2$ coprime and squarefree, and put
$J=sJ_2$.  The exact leftover column factor is
\begin{equation}\label{eq:can-cube-factor}
 \chi_n(J)^4\mathbf1_{(n,\rad q_0)=1}.
\end{equation}
Thus $q_0$ is frozen as a moving puncture and $J$ remains an averaged
squarefree label.  

\smallskip\noindent\emph{Mark allocation and the common kernel.}
For every extracted list $D_1,\ldots,D_h$, use the exact priority
identity
\begin{equation}\label{eq:can-priority}
 \mathbf1_{p\mid D_1\cdots D_hn}
 =\sum_{j=1}^h\mathbf1_{p\mid D_j}\prod_{a<j}\mathbf1_{p\nmid D_a}
   +\mathbf1_{p\mid n}\prod_{a\le h}\mathbf1_{p\nmid D_a}.
\end{equation}
This holds with shared primes.  Assigned marks give divisor-bounded
outer coefficients; surviving marks keep their original product
coefficients.  The two copies of a slot in a square have independent
prime choices.

Truncate only a sector-fixed outer frequency ball on the genuine
coprime expression.  Its complement is contained in a large-argument
kernel tail for every supported column pair.  Inside the ball retain
the full smooth kernel, then apply \eqref{eq:can-full-mobius}.
Logarithmic Fourier inversion uses one joint profile in the bare
frequency and all normalized extracted and column norms.  Weighted
Cauchy on the whole outer sum now gives two positive inverse-column
norms.  For the first transform's new row $y=hf^2E$, the old-label
fibre $f^2E\mid y$ has divisor-bounded mass, so the old $f,h$ disappear
at this positive step.  

\smallskip\noindent\emph{Second transform and the child class.}
For one such positive norm, write $x_i=g'z_i$ in its expanded
square.  The second primitive conductor is $z_1z_2$ and its mask
is $g'$.  Apply the same orbital Poisson formula, with divisor
$d_2\mid g'$, and retain the analogous outer ball before the full
Möbius extension.  Formula~\eqref{eq:can-calibration} returns the
coefficient $a_0\nu\vartheta\vartheta'$.
The new row and squarefree label are, in the notation of
\cite[(17.42)--(17.55)]{Q},
\[
 k_{\rm new}=d_kd_2k'',\qquad f_{\rm new}=CJd_2v',\qquad
 \rho_{\rm new}(n)=\rho(n)
            \mathbf1_{(n,\rad(q_0t'g'/d_2))=1}.
\]
The zeros retained in the signed outer scalar force
$f_{\rm new}$ to be squarefree.  Each child test, background,
puncture, and surviving mark is independent of its new row and label.
Only after the second weighted Cauchy inequality are these outer
zeros bounded by one and the positive row and label ranges enlarged.

At this point the child has all fields of the parent state: a
squarefree averaged label, an independent fixed-type background, a
common puncture, and marks with their original individual coefficients.
The old-label fibres are counted by divisors of the new ideals; fixing
a unit sector leaves only a finite multiplicity. This identifies the
recursive input before any numerical estimate is applied.

\smallskip\noindent\emph{Lengths and normalization.}
Here are the identities controlling the induction.
Use the idealized logarithmic centers of \cite[(17.45)--(17.48)]{Q},
write $j$ for the length of $J$, and put
$D_c=\ell+A_i+t+g-\theta+V$.
Then
\begin{align}\label{eq:can-child-identities}
 F_c&=F-D_c-2\ell+j,& M_c&=M-2D_c-2\ell+j,\notag\\
 F_c-M_c&=F-M+D_c,&
 4F_c-3M_c&=4F-3M+2(A_i+t+g-\theta+V)+j.
\end{align}
The structural source fibres imply
$Q_{\mathrm{new}}\le Q+D_c$, $g\ge\theta$, and $j\le2\ell$,
up to their explicitly bounded annular errors.
If each individual support error is at most $\eta_1$, a nonempty
clipped child satisfies
\begin{equation}\label{eq:can-child-margins}
 \begin{gathered}
 F'-M'-Q_{\mathrm{new}}-z_0\ge c_{\rm node}-5\eta_1,\\
 4F'-3M'-6z_0\ge c_{\rm node}-7\eta_1,\\
 M-M'\ge2(V+\ell)-8\eta_1,\qquad F'\le F+11\eta_1.
 \end{gathered}
\end{equation}
These estimates use the source witnesses for the frozen triple
$(q_0,t',g'/d_2)$.

The reconstruction at fixed triple, new row, and new label has
weight at most $D(f_{\rm new})C(q_0,t',g'/d_2)$, with fixed divisor
powers.  The first factor belongs only to the child label measure.
Normalize the positive measure on distinct frozen triples by its
already extracted norm count.  Its remaining mass is a prescribed
small power, used once after the two Cauchy square roots.
The exact prefactor identity is
\begin{equation}\label{eq:can-energy-identity}
 \kappa_i+\lambda_c+C_c+2F_c=F
 \qquad\text{\cite[(17.64)]{Q}}.
\end{equation}
Here the four terms are, respectively, the first prefactor, second
prefactor, frozen-triple count, and restored child normalization
with its asserted exponent.  Hence a child loss $\eps_{\mathrm{child}}$
gives parent loss at most
$40\eta_1+\tau_1+\pi_1+\eps_{\mathrm{child}}$.
The two primitive principal contributions have exponents at most
$M$ and $F+26\eta_1+\tau_1+\pi_1$.
These calculations are norm identities and the finite digit table;
the field change introduces only the fixed sector multiplicity.

\smallskip\noindent\emph{Termination and loss allocation.}
If $D>2+M_{max}/d_0$, choose
$\eta_1\le\min(d_0/16,c_*/(14D))$ and
$D(40\eta_1+\tau_1+\pi_1)<\eps$.
Both margins remain at least $c_*/2$ through depth $D$, and every
nonterminal child decreases the row length by at least $d_0$.
Every branch therefore reaches a terminal before depth $D$.  The terminal
estimates \eqref{eq:can-terminal-ledger} and backward induction
prove \eqref{eq:can-energy}.  The finite-seminorm assertion is
proved in Proposition~\ref{prop:can-uniformity} below.
\end{proof}

\begin{proof}[Proof of Theorem~\ref{thm:can-inverse}]
In a term of $M_u\prod_iQ_{u,i}$ let $P$ be the squarefree product
of its slot primes and put $j=(n,P)$, $n=jn_0$, $P=jP_0$.
Then $c=n_0P_0$ is squarefree and
\[
 \mu_K(n_0)=\mu_K(c)\mu_K(P_0),\qquad
 \lambda_u(n)\lambda_u(P)=\lambda_u(j)^2\lambda_u(c).
\]
All zeros are retained.  If $G$ is the sum of the assigned slot
lengths, the residual cap is $z_0=z-G$ and its column length is
$D'=r+z-2G$.  The factor $Z^{-G}$ cancels the $O(Z^G)$ choices of
the overlap tuple in the row-space triangle inequality.
The scalar $\lambda_u(j)^2$ is a contraction, and the residual
column has the fixed puncture $(c,j)=1$.
The signs $\mu_K(P_0)$ are individual prime coefficients.
One common profile separates the identity
$q_n/Z^r=(q_j/Z^G)(q_c/Z^{D'})/\prod_{i\in I_0}(q_{p_i}/Z^{z_i})$.

Orient the residual polynomial as in
\eqref{eq:can-background-class}.  One masked primitive Poisson
transform, followed by the full mutual-gcd expansion, produces
the canonical variables
\[
 k=d'h',\qquad f=d's,\qquad t=C/d',\qquad
 \rho_{t,j}(n)=\mathbf1_{(n,\rad(tj))=1}.
\]
The common zero $(s,h')=1$ is retained through Cauchy, making the
new label squarefree.  At fixed $t,f,k$ the reconstruction has only
divisor multiplicity, and its $f$-dependent part is the child label
weight.  Before enclosure losses the two canonical margins are
\begin{equation}\label{eq:can-initial-margins}
 m-r-2z+2G,\qquad 3m-2r-8z+10G+2P_1,
 \qquad P_1=\log_Zq_{C/d'}\ge0.
\end{equation}
Their perturbations are at most
$9\eta_1+\tau_1$ and $22\eta_1+3\tau_1$, respectively;
see \cite[(17.85)]{Q}.
Choose them within fixed fractions of $c_1,c_2$ and apply
Proposition~\ref{prop:can-canonical}.
The normalization identity
$\kappa_c^{\rm init}+P_1+2F_c^{\rm init}=m$
then gives $Z^{m+11\eta_1+\pi_1+\eps_{\mathrm{can}}}$.
Principal terms have the direct exponent $m$ and the genuine
frequency tails have arbitrary saving by a fixed kernel order.
Taking the three losses to sum to less than $\eps$ proves
\eqref{eq:can-inverse}, including its arbitrary bounded slot
coefficients and the original masks.
\end{proof}

\subsection{Ordinary reflection and prime factors}\label{subsec:ordinary-reflection}

We first remove a fixed extra mask by an exact divisor identity.  If $T_u(X)$ denotes
the natural plain sum, then
\begin{equation}\label{eq:can-mask-delete}
 S_{u,\mathfrak R}(X)
   =\sum_{d\mid\rad\mathfrak R}
       \mu_K(d)\lambda_u(d)q_d^{-1/2}T_u(X/q_d).
\end{equation}
For a prime slot subtract its terms at primes dividing
$\mathfrak R$.  Apply these identities simultaneously to both
plains and every slot.  Their coefficient masses are bounded by
products of $(1-P^{-1/2})^{-1}$ over a polynomial-size radical,
hence by $Z^\eps$ with any prescribed loss.  Every resulting
length decreases; a nonempty subunit annulus may be replaced by
scale one with a bounded dilation.  If all slots disappear the
unrestricted zero-slot assertion applies.

For a natural presentation let $C_u$ be its primitive conductor
scale, including the fixed discriminant factor, and $R_{0,u}$ its
redundant natural radical.  Tameness and the full declared support
give
\begin{equation}\label{eq:can-conductor-natural}
 C_uq_{R_{0,u}}\ll Z^{m+q}=Z^M.
\end{equation}
At each exterior prime the primitive conductor or the redundant
radical uses that prime once.  The components at $S$
and the infinity type have fixed conductor factors.
Apply the ordinary Hecke functional equation, and only then restore
the deleted redundant factors geometrically.  The dual scales have
the form
\[
 Y=\frac{C_uq_{d_0}}{Xq_{h_0}},
 \qquad d_0\mid R_{0,u},\quad \rad(h_0)\mid R_{0,u},
\]
with coefficient mass bounded by
$\prod_{p\mid R_{0,u}}(1-P^{-1/2})^{-1}$.
Lemma~\ref{lem:gamma-reflection} gives a summable small-annulus
decomposition and arbitrary decay at the large end.
Thus, with any fixed length tolerance $\xi>0$, every retained
reflected scale has nonnegative declared length
\begin{equation}\label{eq:can-reflected-length}
 n_{\rm ref}\le M-n+\xi.
\end{equation}
If the right side is negative there is no retained nonempty annulus.
Conjugating that reflected factor and its profile restores the
original character inside the absolute product.  This is the uncentered operation of Lemma~\ref{lem:paired-phase-invariant}.
The coefficient mass and the permitted rowwise scale supremum cost
only a prescribed small power and finite seminorm orders.

For $\kappa<1$, the zero-free premise in
Theorem~\ref{thm:can-fourth} gives, a fixed distance $e>0$ to the
right of $(1+\kappa)/2$, the logarithmic-derivative bound used in
\cite[(18.11)]{Q}.  The proof by the completed Hecke $L$-function
and its zeros applies over $K$: the degree contributes only a
fixed multiple of the height logarithm.  Mellin inversion of a
smooth von Mangoldt sum, shifted to that line, and removal of
prime powers give the finite-character prime estimate.  On the
permitted rows every $\xi_u^*\vartheta_{ij}$ is nonprincipal.
Consequently a whole fixed product of slots satisfies
\begin{equation}\label{eq:can-prime-bound}
 \left|\prod_iQ_{u,i}\right|^2
 \ll Z^{\kappa z+\eps_1}
       p_J((W_i))^b(1+|\boldsymbol t|)^A.
\end{equation}
The total contribution $2ez$ is included in $\eps_1$.  At $\kappa=1$ the same assertion follows
by absolute prime counting.  The estimates and their norm-power
derivatives are uniform for $3/4\le\kappa\le1$.
The natural redundant radical deletes only $O(\log Z)$ possible
prime labels; its normalized contribution is within the same loss.

\subsection{The exceptional terminal}

\begin{lemma}[Equal-product cancellation]\label{lem:can-exceptional}
Let $\lambda=\omega_0\vartheta$, $\vartheta\in\Theta$, with global
norm twists removed as above.  For a common radical
$\mathfrak R_*$ of norm at most $Z^{B_*}$ put
\[
 L_i(X,t)=\sum_{\mathfrak a}\lambda(\mathfrak a)
       \mathbf1_{(\mathfrak a,\mathfrak R_*)=1}
       q_{\mathfrak a}^{it}W_i(q_{\mathfrak a}/X).
\]
Uniformly for $X>0$,
\begin{equation}\label{eq:can-exceptional-main}
 L_i(X,t)=c_{\lambda,\mathfrak R_*}X^{1+it}I_i(t)
           +O\bigl(Z^{\eps_1}p_J(W_i)(1+|t|)^A\bigr),
 \qquad |L_i(X,t)|\ll p_0(W_i)X,
\end{equation}
where $I_i(t)=\int_0^\infty W_i(y)y^{it}\,dy$ and the main
coefficient is independent of $X,t$.
It vanishes unless $\lambda$ is principal.
If $U_1U_2=V_1V_2=T$ and all four scales are at least $Z^a$, $a\ge0$,
then
\begin{equation}\label{eq:can-exceptional-centered}
 T^{-1/2}|L_1(U_1,t)L_2(U_2,t)-L_1(V_1,t)L_2(V_2,t)|
 \ll Z^{\eps_1}T^{1/2}Z^{-a}p_J((W_1,W_2))^2(1+|t|)^{2A}.
\end{equation}
Without a nonnegative lower length the same expression is
$O(p_0((W_1,W_2))^2T^{1/2})$.
\end{lemma}
\begin{proof}
First omit the mask.  In the Mellin convention used here the exact
integral on an absolute line $c>1$ is
\[
 L_i^{\rm unmasked}(X,t)
   =\frac{X^{it}}{2\pi i}\int_{(c)}
       L_K^S(s,\lambda)\widehat W_i(s+it)X^s\,ds.
\]
Hecke continuation \cite{Tate} therefore gives the residue at
$s=1$ with its stated factor $X^{1+it}I_i(t)$.
For $X\ge1$, shifting to $\Rea s=-1$ gives an
$O(X^{-1}p_J(W_i)(1+|t|)^A)$ remainder.  A norm-reduced genuine background has no pole;
the only remaining pole is the principal one.
For $X<1$ bounded ideal counting gives an $O(p_J(W_i))$ remainder.
Now insert the common mask by inclusion--exclusion.  In a divisor
term the factor $q_d^{it}$ cancels the factor $q_d^{-it}$ from
the residue at scale $X/q_d$.  Thus
\[
 c_{\lambda,\mathfrak R_*}
  =\mathop{\rm Res}_{s=1}L_K^S(s,\lambda)
       \sum_{d\mid\rad\mathfrak R_*}\frac{\mu_K(d)\lambda(d)}{q_d}.
\]
Use a subsidiary divisor loss at most $\eps_1/4$ here, so the
products of the resulting errors below still fit within $Z^{\eps_1}$.
This proves \eqref{eq:can-exceptional-main}; its absolute
bound also follows directly from the original ideal sum, including
the empty subunit range.
The two products have identical residue term
$c_{\lambda,\mathfrak R_*}^2T^{1+it}I_1(t)I_2(t)$, which cancels
before any absolute value.  Each cross error is bounded by the
largest of the four scales, at most $TZ^{-a}$; also
$1\le TZ^{-a}$.  Division by $T^{1/2}$ proves the assertion.
\end{proof}

\subsection{The fourth-moment induction}\label{subsec:fourth-induction}

\begin{proof}[Proof of Theorem~\ref{thm:can-fourth}]
Put $A=n_1+n_2+z$. The effective row width is $M=m+q$:
the logarithmic norm of the declared moving support contributes
$q$ in addition to the row length $m$. We use the finite induction of
\cite[Sections~18.2--18.8]{Q}, organized by the following ordered state:
\[
 (\text{zero or positive slots},\ \text{width band},\
   \text{uncentered or centered stage}).
\]
All zero-slot bands precede the positive-slot bands. Within a width
band the uncentered estimate is proved first. A comparison term returns
to that earlier stage in the same band; a two-transform child has
strictly smaller width. This gives a finite dependency order.

The coefficient that must be preserved is one whole-product character
on each rectangle, with the same norm phase and coprimality mask.
Lemma~\ref{lem:paired-phase-invariant} preserves that class, and
Lemma~\ref{lem:can-exceptional} estimates its centered exceptional
members. The remaining steps below compute the lengths and the
available loss at each reduction.

\smallskip\noindent\emph{Base ranges and ordinary reflection.}
Choose a width floor $\rho>0$, a step $\sigma>0$, and a padding
$\delta_p>0$ small compared with both.  Divide the bounded width
range into bands of length $\sigma/4$.  Complete all zero-slot
bands first, then the positive-slot bands.  In each band prove
the uncentered range $A\le5M/6$ first.  For zero slots,
\eqref{eq:can-reflected-length}, applied once to each plain longer
than $M/2$, reduces all lengths to the padded core
$n_i\le M/2+\xi$, $n_1+n_2\le M+\delta_p$.
The case $M\le\rho$ follows from absolute counting with loss
$\rho+\delta_p$; positive slots add at most $\rho/6$ by
\eqref{eq:can-prime-bound} and the capacity inequality.

\smallskip\noindent\emph{Comparison rectangles.}
For $A>5M/6$ put $L=M/4$.  If one plain is shorter than $L$,
reflect the other.  Otherwise subtract the product at comparison
scales $Y_1=Z^L$, $Y_2=X_1X_2/Y_1$, keeping the same profiles
and every slot.  Reflecting the longer comparison factor gives
total length at most
\[
 A_{\rm comp}\le3M/2-A+2z+\xi.
\]
For $z=0$ this has a margin at least $M/6-\xi$ inside the
uncentered range.  For $z>0$, capacity implies $z<M/20$,
and both required boundaries have margin at least $M/15-\xi$;
these are the inequalities in \cite[(18.15)--(18.16)]{Q}.
The comparison therefore uses an already proved assertion in
the same band.  It remains to bound the whole centered difference.

After any frozen allocation $b_1,b_2$ its coefficient contains
\[
 D_b(l_1,l_2)=
  \prod_{i=1}^2W_i(q_{b_i}q_{l_i}/X_i)
    -\prod_{i=1}^2W_i(q_{b_i}q_{l_i}/Y_i),
 \qquad X_1X_2=Y_1Y_2.
\]
The same allocated powers occur in both rectangles.
An impossible rectangle is zero but is retained formally.
Including the live prime product $\Pi$, the arithmetic factor
on a separated side is one whole-product character, one norm
power, and one common mask:
\begin{equation}\label{eq:can-centered-class}
 \left(\prod_{i\in I'}a_i(p_i)W_i^{\rm slot}(q_{p_i}/P_i)\right)
 D_b(l_1,l_2)\lambda_1(\Pi l_1l_2)
 q_{\Pi l_1l_2}^{it}
 \mathbf1_{(\Pi l_1l_2,\mathfrak R)=1}
 \mathbf1_{s\mid\Pi l_1l_2}.
\end{equation}
The last divisibility condition is omitted when $s=1$.
For an uncentered input omit the second rectangle.

\smallskip\noindent\emph{First transform and its norm allowance.}
Apply the first orbital Poisson transform to the nonnegative
row norm of the whole expression.  Its zero frequency requires
the product of the two full indices to be powerful, giving
$O(Z^{m+\eps_1})$ after central normalization.
For nonzero frequencies extract the complete common support
as $Ca,Db$.  Write $c,d,p$ for the lengths of $C,D$ and their
common radical, $R_1$ for the nonzero-exponent radical length,
and $E_1$ for the divisor of the complementary mask.  Thus
$R_1+E_1\le p$.  The new displayed moving support has length
$\widetilde q=q+R_1+E_1$.
The nominal frequency length and the C-side allowance are
\begin{equation}\label{eq:can-plain-first}
 K_0=2A-c-d+R_1+E_1-m,\qquad
 N_C\ll Z^{A-c+\widetilde q+B_c-s_0+\eps_G},
 \quad B_c=\left(\frac{3c-5d-R_1}{6}\right)_+,
\end{equation}
where $s_0$ is the length of the full mutual-gcd divisor.
There is a symmetric $D$-side allowance.
The outer prefactor, common-support count, and two Cauchy
halves sum to
$M+\{B_c+B_d-(c+d-2p-R_1)\}/2$.
This is at most $M$: at a common prime of multiplicities
$i\ge j\ge1$, with $r_p=\mathbf1_{6\nmid i-j}$,
\[
 \frac{(3i-5j-r_p)_+}{6}\le i+j-2-r_p.
\]
Summing with the actual logarithmic prime norms proves the
inequality, including every ramified multiplicity.

\smallskip\noindent\emph{One prime-pool amplification.}
We explain how the C-side allowance is obtained.  Let $K$ be
the retained Gauss-row length, and let $w_o\le w$ be the
moving-support and total lengths removed by an amplifier
error.  Put
\[
 a_0=A-c-w,\qquad
 J=d-c+(K_0-K)-2w+w_o.
\]
For zero slots enlarge directly by $g=J_++\sigma$ and put
$\ell_{\rm pool}=0$ in the prefactor.
For positive slots choose one pool of good primes of length
$\ell_* =\sigma/3$, disjoint from every live slot by requiring
the mesh $\eta<\sigma/6$.
The prime ideal theorem gives $Z^{\ell_*+o(1)}$ such primes.
For each row omit its prime divisors and all frozen supports;
only $O(\log Z)$ labels are lost.
Comparing $h$ with $h\mathfrak p^6$ gives a main positive norm
with density factor $Z^{-\ell_*+o(1)}$, since each output has
boundedly many such representations.  Here take
$g=J_++2\sigma$ and $\ell_{\rm pool}=\ell_*$.
The only error valuations are $1,6,7$, by
\eqref{eq:can-local-G}.  Their squared central coefficients are
$P^{-1},(1-P^{-1})^2,P^{-1}$, respectively.
Thus $(w,w_o)=(i\ell_p,e_i\ell_p)$, with
$e_1=e_7=1$, $e_6=0$ and $\sigma/6\le\ell_p\le\sigma/3$.
Each error is treated without further amplification, using
$g=J_++\sigma$ and $\ell_{\rm pool}=0$.
The entire centered coefficient is used in each of these
identities, and the pool average is used once.

\smallskip\noindent\emph{Second transform and the width drop.}
The second transform uses the full finite correlation
\eqref{eq:can-local-G}.  After extracting its genuine
complete common supports $D_2,E_2$, write $c_2,d_2$ for their
lengths, $b_2=(c_2+d_2)/2$, $p_2$ for the common radical length,
and $g_2$ for the length of their gcd.
At equal non-six-divisible multiplicities let $t_2$ be the
radical length where the divided frequency is a unit, and
$V_2$ the corresponding nonunit radical length.
The local correlation has size at most $Z^{g_2-t_2}$ on its
retained partition and is extended by zero off that partition.
Writing the frequency as $\gcd(D_2,E_2)V_{\mathrm{id}}h'$ gives
\begin{align}\label{eq:can-plain-child}
 m'&=2a_0-K-g-g_2-V_2,&
 q'&=\widetilde q+w_o+t_2+V_2,\notag\\
 M'&=m'+q'=M+J-g-g_2+t_2\le M-\sigma,\\
 \Delta_{\rm child}&=b_2-p_2+w+B_c+\ell_{\rm pool}
                                      \ge w+\ell_{\rm pool}.\notag
\end{align}
The second line uses $g_2\ge t_2$.
The exponent available for the inner products is exactly
$M'+\Delta_{\rm child}$, after the common-support count,
Poisson factor, central normalizations, and pool density have
each been included once; see \cite[(18.32)--(18.33)]{Q}.

Insert the full residual coprimality sum before separating the
children.  At a divisor prime allocate its occurrence among
the two plains and the live slots by
\[
 \mathbf1_{p\mid\prod_i n_i}
   =\sum_{\varnothing\ne J}(-1)^{|J|+1}
                             \prod_{i\in J}\mathbf1_{p\mid n_i}.
\]
Extracting a plain power divides its two rectangle scales by the
same norm and preserves the common mask.
A selected slot is frozen.  Old masks remain common, and a
slot selected by two distinct divisor primes gives zero.
Both transforms preserve \eqref{eq:can-centered-class}.
The separated child sides have finite inducing-character
ratio in $\Theta$, although their smooth norm powers may
differ.  Hence they have the same finite exceptional status.
The live background on each side remains $\omega_0$ up to
whole conjugation, and \eqref{eq:can-pairing} remains valid.

On nonexceptional rows first take triangle between the two
rectangles, then clip a nonempty subunit plain scale to one.
With all actual boundary errors included, the positive-part
defect of a child capacity is
\[
 F_{\rm act}=\bigl(A_{\rm act}-M'_{\rm act}
                       +(6\kappa-1)z_{\rm act}\bigr)_+
 \le6(w+\ell_{\rm pool})+O(\xi).
\]
Remove whole slots by \eqref{eq:can-prime-bound} until capacity
holds or no slot remains.  One greedy prefix has at most one
overshoot, so the squared cost is
\[
 \kappa d_z\le F_{\rm act}/6+\kappa\eta
          \le\Delta_{\rm child}+O(\xi+\eta).
\]
The remaining product uses the smaller-width induction; if
no slot remains its plain lengths are unrestricted.
Divisor allocations, exact extraction coefficients, and paired Cauchy
contribute a single aggregate support error.

It remains to justify the exceptional alternative.  Every
prime in the genuine second common support is disjoint from
the old displayed moving support, because its frozen scalar
otherwise vanishes before the absolute estimate.
If the child finite inducing character belongs to $\Theta$,
each exterior row valuation is prescribed modulo six.
Thus $(h')=h_0a^6$ in finitely many fixed unit sectors.
At a nonunit common prime of equal multiplicity one the
frozen frequency factor has valuation two, forcing
$v_p(h')\equiv4\pmod6$.
If the radical length of these primes is $v_1$, put $f_1=2v_1$;
then $q_{h_0}\ge Z^{2f_1}$ and, on a nonempty range,
\begin{equation}\label{eq:can-exceptional-count}
 \#\{h':q_{h'}\ll Z^{m'_{\rm act}},\ \xi_{h'}^*\in\Theta\}
 \ll Z^{(m'_{\rm act}-f_1)/6+\eps_1}.
\end{equation}

This count includes all rows admitted by the positive enlargement.

The absolute volumes together with
\eqref{eq:can-exceptional-count} have excess
$A-5M/6-F_1-F_2$, where the local ledger of
\cite[(18.42)--(18.44)]{Q} gives
\[
 F_1\ge\frac23(c+w)-3\sigma-O(\xi),\qquad
 F_2\ge\frac23b_2.
\]
This closes the uncentered range $A\le5M/6$.
For a centered input let
$v_*=c+w+\min(c_2,d_2)$.
On the side with that smaller extraction, every formal plain
scale initially has length at least $L-v_*$.
Lemma~\ref{lem:can-exceptional} gives saving
$(L-v_* -r_1)_+$ after a final divisor extraction of total
length $r_1$ on that side; if this lower length is nonpositive
use its absolute-volume alternative.
The divisor count and the two raw extraction coefficients
retain a saving of $(L-v_*)_+$, since
\[
 t_- -r_1-r_2-(r-r_1)_+
       =(t_- -r_2)-\max(r,r_1)\le\omega_2-r,
 \qquad r=(L-v_*)_+.
\]
Here $t_-$ is the lower divisor-dyad length and $\omega_2$ is
the actual frozen-slot ratio error on the other side.
Therefore the exceptional excess is at most
\[
 A-\frac56M-\frac23v_*-(L-v_*)_++3\sigma+O(\xi)
 \le(A-M)_++3\sigma+O(\xi),
\]
because the maximum in $v_*\ge0$ occurs at $v_*=L=M/4$.
This is at most $3\sigma+O(\xi)$ in the positive-slot region,
and $\delta_p+3\sigma+O(\xi)$ in the padded zero-slot core.
Keep the centered scales formal and perform their common residue
cancellation before taking absolute values.

The second-transform diagonal has excess at most
$5\sigma/3+O(\xi)$, with an additional $\delta_p$ in the
padded core; its local diagonal is the full Euler function
of the common modulus.  Frequency tails are removed as
whole common dyads on the genuine sums, before the artificial
coprime extension and Fourier absolutization, using the fixed
rapid-decay order of the orbital kernels.
Equations~\eqref{eq:can-plain-child} and the specified clipping
give a strict width drop at least $\sigma/2$ once
$\xi$ is sufficiently small.
The number of strict calls is at most
$D=2+\lceil2M_{max}/\sigma\rceil$.
Within a band, comparison calls the earlier uncentered stage. Each
Gauss norm is amplified at most once, and each positive child removes
slots once after all capacity defects have been included.
Choose the floor and step so that
\[
 \rho+\delta_p+\rho/6+3\sigma+5\sigma/3<\eps/4,
\]
then choose $\xi$, the mesh $\eta$, and the aggregate local
power loss sufficiently small that their sum through $D$
stages is less than $\eps/2$.
The backward finite propagation below handles the analytic
orders.  This proves the natural assertion; the exact mask
deletion \eqref{eq:can-mask-delete} proves the full statement.

The field-dependent inputs in this induction have now all been
specified: the complete correlation \eqref{eq:can-local-G}, orbital
Poisson with one common kernel, the ideal prime pool, and the masked
Mellin cancellation. The local multiplicity inequalities use only
valuations. Multiplying their contributions by the actual logarithmic
prime norms gives the displayed capacities and the strict decrease of
$M$. The finite ideal fibres and sector counts enter the allocated
small-power loss; Lemma~\ref{lem:polylog-propagation} handles the
remaining profile logarithms.
\end{proof}

\subsection{Uniformity}

\begin{lemma}[Finite propagation with logarithmic profiles]\label{lem:polylog-propagation}
Consider a finite-depth canonical reduction with the arithmetic
powers, normalized outer counting measures, common Fourier measures,
and fixed nonnegative Cauchy or H\"older exponents of
\cite[Lemma 17.3 and Corollary 17.4]{Q}.
Replace a uniform bound at any fixed Euler or weighted Fourier order
by that bound times $(1+\log Z)^A$, where $A$ is finite and may
depend on the eventual fixed arithmetic data and on the order.
Then the same arithmetic scale exponents propagate through the
reduction.  At the root there are finite $A_*,J_*,H_*,b_*$ such that
its conclusion differs from the original one by at most
\begin{equation}\label{eq:polylog-root-bound}
 C(1+\log Z)^{A_*}p_{J_*}(w)^{b_*}\langle t\rangle^{H_*}.
\end{equation}
Consequently every preassigned additional loss $Z^\eta$, $\eta>0$,
absorbs all these logarithms.  The threshold and the finite orders,
but not $\eta$, may depend on those later fixed data.
\end{lemma}
\begin{proof}
Use backward induction on the finite reduction graph.
At a terminal the assertion is its finite-order estimate, with a
possibly nonzero logarithmic degree.  At an edge, the child estimates
specify only finitely many seminorm orders and height degrees.
Choose the parent Euler order large enough to bound the corresponding
weighted $L^1$ norm of the one common Fourier measure.  The modified
edge assumption adds a finite logarithmic factor.  Products,
fixed positive powers, and the finite sums in the edge add or
multiply logarithmic degrees by fixed finite numbers.  Use the Cauchy exponents attached to each edge.  Norm powers in the edge combine
exactly as in the original arithmetic estimate.  A normalized outer
measure has already had its norm count extracted, so its remaining
logarithmic mass adds only another such factor.

Iterating gives finite $A_*,J_*,H_*,b_*$ at the root and proves
\eqref{eq:polylog-root-bound}.  Finally
$(1+\log Z)^{A_*}\ll_{A_*,\eta}Z^\eta$.

\end{proof}

\begin{remark}[Application to norm sheets]\label{rem:polylog-norm-sheets}
The orbital kernels used here are rapidly decreasing at infinity and
have at most fixed-degree logarithmic growth at zero, after every
fixed Euler derivative.  On a retained annular block their nonzero
arguments lie in $[cZ^{-C_0},CZ^{C_0}]$: all moving norm lengths are
bounded, nonzero exterior ideal norms are at least one, and fixed
support constants enter only $c,C$.  Therefore each such derivative
is bounded by a finite power of $1+\log Z$ after its assigned central
powers and decay factors have been removed.  On a product of normalized
annuli, derivatives of $K(R_0\prod_i y_i^{a_i})$ are fixed linear
combinations of Euler derivatives of $K$.  Logarithmic Fourier inversion, with a sufficiently
large but finite order, gives exactly the weighted edge bounds in
Lemma~\ref{lem:polylog-propagation}.
Treat the zero orbits separately by
Lemma~\ref{lem:weighted-orbital-poisson}. The logarithmic propagation
lemma then applies to both canonical inductions.
\end{remark}

\begin{proposition}[Uniform real parameters]\label{prop:can-uniformity}
For a fixed field and fixed infinity-type orbit, all real losses,
strict-width decrements, depth bounds, and the mesh in
Theorem~\ref{thm:can-fourth} can be chosen before a later
finite-order target twist and its conductor.  The fourth-moment
mesh is also chosen before the fixed slot count, uniformly for
$\kappa\in[3/4,1]$.  Once the target, final $S$, and slot system
are fixed, both moments have finite seminorm orders and fixed
polynomial height degrees.  An external tail order chosen in
a later application does not change those internal degrees.
\end{proposition}
\begin{proof}
The arithmetic proofs use finite depth and finitely many
factorization types.  The inverse decrease and its losses are
\eqref{eq:can-child-margins}; the ordinary decrease is
\eqref{eq:can-plain-child}.  They are affine identities and
positive parts of affine functions of a fixed list of total
lengths.  In the ordinary capacity the Lipschitz bound is
\[
 |\Delta M|+|\Delta A|+5|\Delta z|,
 \qquad 0\le6\kappa-1\le5.
\]
It is independent of the number of slots.
For a fixed slot system let $H_I$ bound the sum of the absolute
logarithms of all its support endpoints and of the finitely
many later support enlargements.  Every subset then has one
total ratio error at most $H_I/\log Z$.  Keep the extracted prime norms exact.  This makes every stage
error at most $C_*(\xi+\eta+\eps_0)$ with a numerical $C_*$
independent of the slot count, after the threshold for
$H_I/\log Z$ is imposed.  Choose $\rho,\sigma,\delta_p$, then
$D$, and finally $\xi,\eta,\eps_0$ so that
$C_*D(\xi+\eta+\eps_0)<\eps/2$ and $\xi$ preserves the
strict width and comparison margins.  These are real choices
before the arithmetic data and the slot count.

For those later fixed data, choose the subsidiary divisor,
Rankin, and prime-estimate exponents whose aggregate is
$\eps_0$.  Divisor allocations for $I$ slots obey
$C_I^{\omega(\mathfrak a)}\ll_{I,a}q_{\mathfrak a}^a$
for every fixed $a>0$.  Fixed local sector counts are constants.
The logarithms introduced by norm-sheet profiles and compact
unit representatives have a fixed degree.  Every preassigned
positive loss absorbs them at a larger threshold, independently
of that degree.  Fixed conductor and support constants likewise
enter lengths only through $O(1/\log Z)$.

Here is the analytic propagation rule, as in
\cite[Lemma 17.3 and Corollary 17.4]{Q}.
After its central powers are removed, every coupled smooth
factor is one function of normalized norms on a fixed product
of annuli.  A kernel evaluated at $R_0\prod y_i^{a_i}$ obeys
\[
 (y_i\partial_{y_i})K(R_0\prod y_j^{a_j})
       =a_i(x\partial_x)K(x)\big|_{x=R_0\prod y_j^{a_j}}.
\]
  The theta bounds give the assigned decay, and the orbital estimates
give Euler bounds with a possible finite logarithmic degree, at every
fixed order.  Remark~\ref{rem:polylog-norm-sheets} and
Lemma~\ref{lem:polylog-propagation} propagate these factors with any
previously fixed tail decay.  Logarithmic Fourier
inversion of this whole profile has weighted $L^1$ norm
controlled by a finite larger seminorm.  The measure is
common to all rows and live labels in the current Hilbert
space.  Row masks remain outside it.

If a child estimate has height degree $H$ and uses a profile
seminorm of order $J$, the edge's Fourier measure need only
be integrated against a fixed weight determined by $H,J$.
Its weighted bound chooses a finite larger parent order.
Select these orders backwards through the finite depth;
Cauchy edges use the weighted sum of the two child degrees.
Normalize the outer counting measures by their extracted norm counts.
Lemma~\ref{lem:polylog-propagation} absorbs the resulting root
logarithm in one previously reserved loss.  This proves finite seminorm
and height dependence at the root.  It also justifies the permitted scale suprema by a fixed
number of Sobolev derivatives on logarithmic parameter boxes.

  After their combined polynomial height degree
$A_{\mathrm{ht}}$ is fixed, an allotted height loss
$\eps_{\mathrm{ht}}>0$ on a physical range $U\ge Z^{d_{\min}}$
is respected, for example, by
\[
 0<\tau\le
 \frac{d_{\min}\eps_{\mathrm{ht}}}{20(A_{\mathrm{ht}}+1)},
 \qquad T=Z^\tau.
\]
Indeed, every retained height factor is then at most
$Z^{d_{\min}\eps_{\mathrm{ht}}/20}$ up to a fixed constant.
One may subsequently increase the external integration-by-parts
order to make its discarded tail negligible.  That later choice
differentiates only the external tests.  The internal degree
$A_{\mathrm{ht}}$, real losses, and mesh remain fixed.
This completes both quantitative assertions.
\end{proof}

For the physical application the declared moving support has $q=0$
and $u$ is sixth-power-free.  If any exterior prime divides $u$,
its valuation is between one and five and its local Kummer
character is ramified and nonprincipal there.  Hence
$\xi_u^*\in\Theta$ is possible only for the finitely many
fixed unit-sector rows of exterior norm one.  The zero-slot
principal exception is contained in these bounded rows.
Thus every sufficiently large physical row satisfies precisely
the eligibility conditions needed for both moments.

\section{Sixth-power amplification}\label{sec:amplification}

The following amplification keeps the unmarked inverse estimate
uniform as the witness length approaches the row length.  This is the
independent-scale amplification of \cite[Lemma~17.6]{Q}, reproduced with
ideal indices and the original zero extensions.

\begin{lemma}\label{lem:amplification}
Let $U,D\geq1$ be independent scales and let
\[
 M_u(D;W)=D^{-1/2}\sum_n\mu(n)\eta(n)\rho_u(n)W(q_n/D)
\]
be an inverse presentation admitted by Theorem~\ref{thm:can-inverse}.
For every fixed $c>0$ and $\eps>0$,
\begin{equation}\label{eq:amplification}
 \sum_{\substack{u\asymp U\\u\ \mathrm{sixth\mbox{-}power\mbox{-}free}}}
       |M_u(D;W)|^2
 \ll\max\{U,U^{1/6}D^{5(1+c)/6}\}(UD)^\eps.
\end{equation}
The estimate also holds with the permitted scale and smooth-profile
suprema, with their fixed polynomial height factor.
\end{lemma}
\begin{proof}
The retained presentation zeros give
\begin{align}
 \rho_{ua^6}(n)&=\rho_u(n)\ind_{(n,a)=1},\label{eq:power-mask}\\
 M_u(D;W)&=\sum_{d\mid\rad(a)}
  \mu(d)\eta(d)\rho_u(d)q_d^{-1/2}
          M_{ua^6}(D/q_d;W).\label{eq:amp-identity}
\end{align}
The second identity follows by separating the part of a squarefree $n$
supported on $a$.  

Set $H\asymp\max\{U,D^{1+c}\}$ and $A\asymp(H/U)^{1/6}$, allowing
fixed constants in the row cutoff.  Fixed-field ideal counting supplies
$\gg A$ permitted ideal multipliers $a$ with $q_a\ll A$, in a fixed
allowed sector; for bounded $A$ include the unit multiplier and change the
constant.  Average \eqref{eq:amp-identity} over these multipliers, apply
weighted Cauchy--Schwarz and the ideal divisor bound, and use the permitted
supremum in $D/q_d$.  The map $(u,a)\mapsto ua^6$ is injective on
sixth-power-free $u$: reduction of each ideal valuation modulo six recovers
$u$, and division recovers $a$.  Finite class and unit sectors change the
bound by a fixed factor.

The no-slot case of Theorem~\ref{thm:can-inverse}, with base scale $H$,
has a fixed positive margin because $D\leq H^{1/(1+c)}$.
It bounds the resulting row energy by $H(UD)^\eps$.  Division by the
number of multipliers gives
\[
 H/A\asymp U^{1/6}H^{5/6}
       \asymp\max\{U,U^{1/6}D^{5(1+c)/6}\}.
\]
The scale and profile supremum is the finite-seminorm one from
\cite[Lemma~4.5]{Q}.
\end{proof}

For $D=U^r$ in any fixed bounded range of $r$, choose $c$ and the
preliminary loss in terms of the requested final loss.  The effective
exponent in \eqref{eq:amplification} is
\begin{equation}\label{eq:e-r}
 e(r)=\max\left\{1,\frac{1+5r}{6}\right\}.
\end{equation}
If an inverse witness has size at least $U^{\delta r/2-O(\eps)}$,
division by its squared size gives
\begin{equation}\label{eq:unmarked-count}
 \#\cB\ll U^{e(r)-\delta r+\eps}(1+T)^A.
\end{equation}
This is the independent-scale argument of \cite[Lemma~17.6, equation~(19.3), and Remark~19.3]{Q}.  It applies throughout $1/50<\delta\leq1$.

\section{The physical probe}\label{sec:physical}

Fix a target $\psi\in\cF_\omega$, with a global norm twist removed and
the orientation of Proposition~\ref{prop:fixed-type-theta}.  Use one final
set $S$ and the arithmetic normalization of
Proposition~\ref{prop:arithmetic-calibration}.  Put
$R=\mathcal O_{K,S}$, $T_0=K_S^1$, and $\Gamma=(R^\times)^6$.
All the sums in this section run over exterior integral ideals with their
allowed generators, except for the explicitly displayed additive rows in $R$.

\subsection{Definition and Poisson expansion}

Let $\Lambda_K$ and $G$ be the cubic and quadratic Gauss corrections and set
\[
 \vartheta=\bar\Lambda_K\bar G\psi.
\]
Thus $\gamma_1(\mathfrak p)\gamma_2(\mathfrak p)
\vartheta(\mathfrak p)=-\psi(\mathfrak p)$ at every good prime, in the
normalization of \cite[Lemma~3.4]{P}.  Define
\[
 P_G(t)=\frac1{2\sqrt\pi}\exp\!\left(-\frac{(\log t)^2}{4}\right),
 \qquad \widehat P_G(s)=e^{s^2},
\]
and the marked completed row
\begin{equation}\label{eq:physical-B}
 \mathcal B_{m,\psi;D}(Z)
  =\sum_{\substack{c\ \mathrm{squarefree}\\n}}
    \frac{\gamma_2(c)\vartheta(cn^3)}{q_c^{1/2}q_n}
       \chi_{cn^3}(m)\ind_{D\mid cn^3}
       P_G(q_cq_n^3/Z).
\end{equation}
The variable $n$ is unrestricted and may share primes with $c$.
The target acts on the whole completed index $cn^3$.
The Gaussian makes the series absolutely convergent; for example,
$P_G(t)\ll_\sigma t^{-\sigma}$ with $\sigma>1/2$ reduces it to
$Z^\sigma\zeta_K^S(\sigma+1/2)\zeta_K^S(3\sigma+1)$.

Choose $W_{1,\mathrm{add}}\in C_c^\infty((0,\infty))$ nonnegative and
nonzero, and a tensor Bruhat--Schwartz function
$W_{0,\mathrm{add}}$ compactly supported in $K_S^\times$ such that
\begin{equation}\label{eq:base-test}
 \widehat W_{0,\mathrm{add}}(0)=0,\qquad
 \mathcal M(1/6)\ne0,
 \quad
 \mathcal M(z)=\int_{K_S^\times}
       \widehat W_{0,\mathrm{add}}(y)|y|_S^z\,d^\times y.
\end{equation}
These conditions are compatible.  At one complex place subtract two
radial bumps so that the additive integral is zero, while their
$|y|^{5/6}d^\times y$ pairing is nonzero.  At the remaining places take
nonnegative annular factors.  The local zeta functional equation then gives
\eqref{eq:base-test}, as in \cite[equations~(7.1)--(7.2)]{P}.  Notice that only
$W_{1,\mathrm{add}}$ is required to be nonnegative.

For an exterior ideal $b$ put
\[
 g_b(k)=\sum_{d\bmod b}\chi_b(d)e(kd/b).
\]
With the norm section $\ell_x$ and the quotient measure fixed in
Section~\ref{sec:transfer}, define
\begin{equation}\label{eq:physical-I}
 I_{\psi;D}(X,Y,Z)=
 \int_{T_0/\Gamma}\frac1Y\sum_b
 \frac{W_{1,\mathrm{add}}(q_b/Y)}{\sqrt{q_bX}}
 \sum_{m\in R}q_b^{-1/2}g_b(-m)
 W_{0,\mathrm{add}}\!\left(\frac{m}{t_0\ell_{q_bX}}\right)
 \mathcal B_{m,\psi;D}(Z)\,dt_0.
\end{equation}
We write $I_\psi=I_{\psi;1}$.  Each additive row sum has finite support
on a bounded set of quotient representatives, and $m=0$ is absent by the
support of $W_{0,\mathrm{add}}$.  The expression descends to $T_0/\Gamma$.
This is the norm-form probe of \cite[equation~(7.3)]{P}, with the stated mark.
Finite local sector decompositions, when used to estimate it, are recombined
in this definition.

The finite transform underlying its Poisson expansion is worth recording.
For $A=cn^3$ and $h\in R$, let
\[
 \mathscr C(b,A,h)=
 \sum_{\substack{d\bmod b\\h\equiv Ad\ (b)}}
   \chi_b(d)g_{\chi_A}\!\left(A,\frac{h-Ad}{b}\right),
 \qquad
 g_{\chi_A}(A,k)=\sum_{v\bmod A}\chi_A(v)e(kv/A).
\]
Direct finite Fourier inversion gives
\begin{equation}\label{eq:finite-probe}
 \sum_{m\bmod bA}q_b^{-1/2}g_b(-m)\chi_A(m)e(hm/(bA))
      =q_b^{1/2}\mathscr C(b,A,h).
\end{equation}
Indeed, write $m=v+Aj$ and sum first over $j\bmod b$.
This calculation also allows common factors of $b$ and $A$.
The lattice index is $q_bq_A$, whereas the dilation has additive modulus
$q_bX$.  After both square-root factors in \eqref{eq:physical-I}, the
Poisson multiplier is $X^{1/2}/q_A$.  The norm of the Fourier argument is
$Xq_h/q_A$.

Integrate over $T_0/\Gamma$, retain the nonzero orbits $h=ur^6$, and apply
Lemma~\ref{lem:weighted-orbital-poisson} with trivial continuous row
character.  The zero frequency vanishes by
\eqref{eq:base-test}.  Three Mellin inversions yield
\begin{align}
 I_\psi(X,Y,Z)
 &=\frac1{(2\pi i)^3}\int_{(4)}\int_{(4)}\int_{(2)}
 X^{1/2-z}Z^{s+z-1}Y^{w-1}e^{(s+z-1)^2}\notag\\
 &\qquad\qquad{}\times\mathcal M(z)\widehat W_{1,\mathrm{add}}(w)
 \sum_u q_u^{-z}\mathscr F_{u,\psi}(s,w,z)\,dz\,dw\,ds,
 \label{eq:triple-mellin}\\
 \mathscr F_{u,\psi}(s,w,z)
 &=\sum_{\substack{c\ \mathrm{squarefree}\\n,b,r}}
 \gamma_2(c)\vartheta(cn^3)\mathscr C(b,cn^3,ur^6)
 q_c^{-1/2-s}q_n^{-1-3s}q_b^{-w}q_r^{-6z}.
 \label{eq:full-F}
\end{align}
The row set $u$ includes the finite unit sectors and has sixth-power-free
exterior part.  The marked form inserts $\ind_{D\mid cn^3}$ in
\eqref{eq:full-F}.  All interchanges on these initial lines are absolute.
The sheet integral has the common Mellin transform $\mathcal M$.  In particular, the exponent of $Z$ in
\eqref{eq:triple-mellin} is exactly $s+z-1$.

\subsection{Compensation}

Use the fixed geometry
\begin{equation}\label{eq:geometry}
 \begin{gathered}
 \ell=\frac16,\quad b_0=\frac18,\quad M=\frac56,\quad
 l_x=\frac{17}{48},\quad l_y=\frac{23}{48},\quad h=\frac{13}{16},\\
 C(s)=s-\frac{11}{16},\qquad X=Z^{l_x},\quad Y=Z^{l_y}.
 \end{gathered}
\end{equation}
Fix positive slot lengths $\ell_i$ with $\sum_i\ell_i=\ell$ and set
$P_i=Z^{\ell_i}$.  Choose fixed nonnegative nonzero smooth windows $W_i$
with disjoint prime supports before any row masks are inserted.  All primes
outside $S$ are allowed.  At a good prime of norm $Q$, the compensation is
\begin{equation}\label{eq:one-compensation}
 \bar\psi(\mathfrak p)I_{\psi;\mathfrak p}(X,Y,ZQ)
       -Q^{-3/2}I_\psi(X/Q,Y/Q,Z).
\end{equation}
For $J\subseteq\{1,\ldots,J_0\}$ write
$p_J=\prod_{i\in J}p_i$.  The complete finite operator is
\begin{equation}\label{eq:full-compensation}
 I_{\psi,\mathrm{mod}}(Z)=
 \sum_{(p_i)}\prod_{i=1}^{J_0}W_i(q_{p_i}/P_i)
 \sum_{J\subseteq\{1,\ldots,J_0\}}
  (-1)^{|J|}q_{p_J}^{-3/2}\bar\psi(p_{J^c})
 I_{\psi;p_{J^c}}(X/q_{p_J},Y/q_{p_J},Zq_{p_{J^c}}).
\end{equation}
Every subset and its original mask is present.
The target factors have modulus one.

Let $P_{\mathfrak p}$ denote the local sum in \eqref{eq:full-F}, and
$P_{\mathfrak p}^*$ the part with $\mathfrak p\mid cn^3$.
The monomial in \eqref{eq:triple-mellin} shows that
\eqref{eq:one-compensation} replaces the local factor, after extraction of
$Q^{z-1}$, by
\[
 \bar\psi(\mathfrak p)Q^sP_{\mathfrak p}^*-Q^{-w}P_{\mathfrak p}.
\]
Consequently its leading prime sum is
\begin{equation}\label{eq:physical-prime}
 \mathcal Q_i(u;z)=-\sum_{\mathfrak p\notin S}
       W_i(Q/P_i)Q^{z-1}\bar\chi_u(\mathfrak p).
\end{equation}
This is a finite actual prime character, even for an infinite-order target.
For example, orient a detector presentation as $\psi\theta\bar\chi_u$.
The inverse relative slot coefficient is $(\psi\theta)^{-1}$.
For the plain moment conjugate both witnesses and all selected slots;
the live character becomes $\bar\psi\bar\theta\chi_u$ and the relative
slot coefficient is $\psi\theta$.  Thus the actual prime character in
Theorem~\ref{thm:can-fourth} is precisely the one required by
\eqref{eq:physical-prime}.  

\subsection{Euler factors}

At an unramified row prime $\mathfrak p\nmid u$ put
\[
 V=Q^{-6z},\quad W=\chi_u(\mathfrak p)Q^{-w},\quad
 D=\psi(\mathfrak p)\bar\chi_u(\mathfrak p)Q^{-s},\quad
 E=\vartheta(\mathfrak p)^6Q^{4-6s}V.
\]
The finite identity \eqref{eq:finite-probe} and the local Gauss relations give
\begin{equation}\label{eq:unramified-P}
 P_{\mathfrak p}=
 \frac1{1-V}+\frac{W}{1-W}
 +\frac{E}{1-E}\left(\frac{1-Q^{-1}}{1-V}+W\right)
 -\frac{D}{1-E}\left(1+\frac{(Q-1)WV}{1-V}\right).
\end{equation}
This is \cite[equation~(7.19)]{P}.  At primes dividing $u$ the six valuation classes
use the full finite Gauss sum in \eqref{eq:finite-probe}, as in
\cite[equations~(7.21)--(7.22)]{P} and \cite[Lemma~7.1]{Q}.
These ramified factors retain the phase of the local unit part.
Multiplication of all local identities gives
\begin{equation}\label{eq:factorization}
 \mathscr F_{u,\psi}(s,w,z)=
 \zeta_K^S(6z)\frac{L_K^S(w,\chi_u)}{L_K^S(s,\psi\bar\chi_u)}
          H_{u,\psi}(s,w,z).
\end{equation}
The principal poles are at $w=1$ and $z=1/6$.

Using $P_{\mathfrak p}-P_{\mathfrak p}^*=(1-V)^{-1}+W(1-W)^{-1}$,
write the full compensated factor as
\begin{equation}\label{eq:quotient-free}
 G_{\mathfrak p}=
 \frac{(\bar\psi(\mathfrak p)Q^s-Q^{-w})(1-V)(1-W)P_{\mathfrak p}^*
                 -Q^{-w}(1-VW)}{1-D}.
\end{equation}

On the principal residue put $v=Q^{-1}$, $D=\psi(\mathfrak p)Q^{-s}$,
$E=\vartheta(\mathfrak p)^6Q^{3-6s}$, and
$K_{\mathfrak p}=1-vE-(1-v)D$.  Formula \eqref{eq:unramified-P} simplifies to
\begin{align}
 P_{\mathfrak p}&=\frac{1+v}{1-v}\frac{K_{\mathfrak p}}{1-E},&
 P_{\mathfrak p}^*&=\frac{(1+v)(E-D)}{1-E},\label{eq:principal-P}\\
 H_{\mathfrak p}&=\frac{(1-v^2)K_{\mathfrak p}}{(1-E)(1-D)}.
 \label{eq:principal-H}
\end{align}

\begin{lemma}\label{lem:principal-nonzero}
The function $H_\psi^S(s)=H_{1,\psi}(s,1,1/6)$ is holomorphic and nonzero
for $\Rea s>2/3$.  On $\sigma=\Rea s\geq3/4$, the principal compensated
quotient $B_{\mathfrak p}=G_{\mathfrak p}/H_{\mathfrak p}$ satisfies
\begin{equation}\label{eq:B-local}
 B_{\mathfrak p}+1=
 \frac{(1-v)\{E/D-vE-(1-v)D\}}{K_{\mathfrak p}}
       =O(Q^{-\sigma}),
\end{equation}
uniformly in all local phases and in $\Ima s$.
\end{lemma}
\begin{proof}
For $\sigma>1/2$, both $|D|$ and $|E|$ are less than one, so
\[
 |vE+(1-v)D|\leq v|E|+(1-v)|D|<1.
\]
Thus each factor in \eqref{eq:principal-H} is nonzero.
On a closed half-plane strictly to the right of $2/3$,
\[
 H_{\mathfrak p}-1
      \ll Q^{-2}+Q^{-1-\sigma}+Q^{3-6\sigma}.
\]
The prime-ideal sum converges normally, proving the assertion for the
product.  On the principal locus
$B_{\mathfrak p}=P_{\mathfrak p}^*/(DP_{\mathfrak p})-v$.
Substitution gives \eqref{eq:B-local}, and
$3-5\sigma\leq-\sigma$ when $\sigma\geq3/4$ proves its bound.
\end{proof}

\subsection{The principal normalization}

The only physical row whose exterior sextic character is trivial is the
principal Kummer class.  Indeed, Chebotarev detects every nontrivial Kummer
extension.  Its sixth-power frequencies are $h=r^6$, with unit ambiguity
in $\Gamma$; finer finite sectors used for positive estimates are recombined.
The double residue of \eqref{eq:triple-mellin} has a fixed scalar
\begin{equation}\label{eq:residue-c}
 c_\psi=\frac16\bigl(\Res_{s=1}\zeta_K^S(s)\bigr)^2
                  \mathcal M(1/6)\widehat W_{1,\mathrm{add}}(1)\ne0.
\end{equation}
The exact theta normalization has already removed its nonzero
archimedean scalar.  A different fixed measure normalization merely changes
$c_\psi$ by a nonzero constant.  Its sign is immaterial.

Set
\begin{align}
 S_i(Z)&=\sum_{\mathfrak p\notin S}W_i(Q/P_i)Q^{-5/6},&
 \mathcal N(Z)&=Z^{-\ell/6}\prod_{i=1}^{J_0}S_i(Z),
 \label{eq:positive-N}\\
 J_\psi(Z)&=\frac{(-1)^{J_0}I_{\psi,\mathrm{mod}}(Z)}{c_\psi\mathcal N(Z)},&
 f_\psi(Z)&=\frac1{2\pi i}\int_{(2)}
        Z^{C(s)}e^{(s-5/6)^2}
                  \frac{H_\psi^S(s)}{L_K^S(s,\psi)}\,ds.
 \label{eq:signal}
\end{align}
The prime ideal theorem gives
$S_i(Z)>0$ eventually and $\mathcal N(Z)\asymp(\log Z)^{-J_0}$.
Thus its reciprocal costs an arbitrarily small power.

The weights
$\mu_{i,Z}(\mathfrak p)=W_i(Q/P_i)Q^{-5/6}/S_i(Z)$ are probabilities.
Disjointness makes the normalized principal correction exactly
\[
 H_\psi^S(s)\prod_{i=1}^{J_0}
       \sum_{\mathfrak p}\mu_{i,Z}(\mathfrak p)(-B_{\mathfrak p}(s)).
\]
By Lemma~\ref{lem:principal-nonzero}, this is
$H_\psi^S(s)\{1+O_\psi(Z^{-(7/8)\min_i\ell_i})\}$ on $\Rea s>7/8$.
For equal lengths reserve the positive margin
\begin{equation}\label{eq:mu0}
 \mu_0=\frac7{16}\min_i\ell_i=\frac7{96J_0}.
\end{equation}
All these expressions use the same $S$ and are independent of the
analysis height cutoff.

\subsection{The low estimate}

\begin{proposition}\label{prop:physical-low}
For each fixed slot system and every $\eps>0$, the same physical function
in \eqref{eq:signal} satisfies
\begin{equation}\label{eq:all-low}
 |J_\psi(Z)|\ll_{\psi,\eps} Z^{3/16+\eps}.
\end{equation}
The exponent is uniform in the slot mesh.
\end{proposition}
\begin{proof}
Fix the rescaled subset $J$ in \eqref{eq:full-compensation}, put
$d=\sum_{i\in J}\ell_i$, and first hold its prime tuple fixed.  Write
\[
 X'=X/q_{p_J},\quad Y'=Y/q_{p_J},\quad
 M'=\frac56-2d,\quad\ell'=\frac16-d,\qquad0\leq d\leq\frac16.
\]
Sum the surviving marked slots in the completed factor before taking its
row norm.  Lemma~\ref{lem:can-marked-energy} gives
\begin{equation}\label{eq:marked-low-energy}
 \sum_{m\text{ in the actual row box}}
            |\mathcal B^J_{m,\psi}(Z)|^2\ll Z^{M'+\eps}.
\end{equation}
In the reflected energy notation, let $O$ be the powerful
row length, $H$ the actual residual row length, and
$\Delta_H=M'-O-H\geq-O(1/\log Z)$.  The local conductor ledger gives
$2A_0\leq O+O(1/\log Z)$, $N_0\leq A_0$, and $z_a\leq\ell'$.
On a Gaussian annulus the completed length is
$N_*=1+\ell'+\theta_N$, with $|\theta_N|$ as small as prescribed.
The dual length is
\[
 T_d=2H+2A_0+2z_a-1-\ell'-\theta_N-N_0-3B_0,
 \qquad T_d\leq H-3d+O(|\theta_N|+1/\log Z).
\]
The two branches of the hybrid minimum in the reflected energy are at most
\[
 M'-O/2-\Delta_H+O(\eps_0+|\theta_N|),\qquad
 M'+\frac{d-1/6-2\Delta_H+\theta_N}{4}+O(\eps_0).
\]
Both are bounded by $M'+O(\eps_0+|\theta_N|+1/\log Z)$.
This is the algebra in \cite[Lemma~15.1]{Q}, applied to the actual row
lengths.  The common Gaussian profile has summable homogeneous seminorms
over its annuli.  Extracting each small profile amplitude before the
positive enlargement and taking the preliminary losses smaller proves
\eqref{eq:marked-low-energy}.  

For the additive factor of \eqref{eq:physical-I}, logarithmic Fourier
separation of $q_b/Y'$ leaves one common profile and the polynomials
$\mathcal A_m(Y')$ of Proposition~\ref{prop:additive-gram}.  Put
$Q_0\asymp X'Y'$ and $P_a=Y'^2/Q_0\asymp Z^{1/8}$; fixed arithmetic
scale factors are harmless.  Since
\[
 l_y-d-\frac{11}{6}\cdot\frac18\geq\frac1{12},
\]
the Gram bound reduces to
\begin{equation}\label{eq:gram-specialized}
 \sum_m|\mathcal A_m(Y')|^2
       \ll (1+|\nu|)^{J_{\rm Gr}}(Q_0/Y')P_a^{1/6}Z^\eps.
\end{equation}
Cauchy--Schwarz in the same row box and the rapid decay of the separating
Fourier density bound this tuple's separation by
\[
 Q_0^{-1/2}\{(Q_0/Y')P_a^{1/6}\}^{1/2}Z^{M'/2+\eps}
       \ll (X')^{1/2}P_a^{1/12}Z^\eps.
\]
There are at most $Z^{d+\eps}$ rescaled tuples, their coefficient has size
$O(Z^{-3d/2})$, and $(X')^{1/2}\asymp X^{1/2}Z^{-d/2}$.
Their extra exponent is $d-3d/2-d/2=-d\leq0$.
Summing the finitely many subsets gives
\[
 |I_{\psi,\mathrm{mod}}(Z)|
       \ll Z^{l_x/2+b_0/12+\eps}=Z^{3/16+\eps}.
\]
Divide by \eqref{eq:positive-N}, decreasing the preliminary loss.
\end{proof}

\subsection{The high decomposition}

We record the precise physical estimates used in the row argument.
Their real parameters are chosen before the later finite-order target;
their finite height degrees may depend on that target.

\begin{proposition}\label{prop:physical-high}
Suppose $B>7/8$.  The fully compensated probe has the row decomposition,
zero witnesses, and local-error estimates of \cite[Sections~8--10 and~16]{Q}, with ideal norms and the family \eqref{eq:family}.  More explicitly:
\begin{enumerate}[label=\textup{(\roman*)}]
\item The detector uses $a\in[51/100,B]$, with
$\delta=2a-1\in[1/50,2B-1]$.  Its saturated witnesses have
$r\leq t+O(\eps)$, $r+m\geq t-O(\eps)$, and
$0\leq m\leq1/2+O(\eps)$ for any fixed $t\in[1,3/2]$.
Every live Hecke denominator belongs to $\cF_\omega$.
\item If a moderate row class of norm $U=Z^d$ has cardinality
$\ll U^{R+\eps}(1+T)^A$ and aggregate prime-amplitude parameter $q$,
its exponent relative to $C(7/8)=3/16$ is at most
\begin{equation}\label{eq:physical-row}
 E(d;R,q)= -\frac1{48}+\frac23\delta+\frac q6
   -h(1-R)+(d-h)\left(R+\frac\delta2-\frac{17}{50}\right),
\end{equation}
up to the allocated bin, contour, and correction losses.  Both selected
and unselected prime slots remain in the full correction.
\item The floor bin, small rows, and the absolute large-row contour have
the estimates used in Section~\ref{sec:rows}.  The principal $w$ and $z$
contours have respective margins $23/960$ and $13/9600$, measured against
$C(B)$.  The principal slot approximation has margin $\mu_0$.
\end{enumerate}
All assertions allow arbitrarily high external tail orders after their
internal profile and height degrees have been fixed.
\end{proposition}
\begin{proof}
The detector is \cite[Section~8]{Q}, with ordinary Hecke continuation and
the fixed-type functional equation.  The gamma estimates in
Lemma~\ref{lem:gamma-reflection} give its finite-seminorm height bounds.
Its argument allows $a\leq1$; the endpoint $B=1$ is also included in
\cite[Lemma~10.5]{Q}.  The conductor of a sextic row is bounded by its
squarefree radical with fixed local factors.  Thus the conductor powers in
the witness construction are unchanged when element norms are replaced
by ideal norms.  The live-background verification in
Section~\ref{sec:canonical} puts every companion denominator under the
same boundary $B$.

For the local correction use \eqref{eq:quotient-free} and the full
ramified Gauss table.
The cancellation producing \eqref{eq:physical-prime} is a local identity
for every good prime.  The remaining phase factors have modulus one.
Consequently the two Euler regions and joint error-slot estimates of
\cite[Section~16]{Q} apply with the same powers of $Q$.  On their dynamic contours,
for a fixed positive allowance $e$, the unmarked corrections satisfy
\[
 |H_{u,\psi,\mathfrak p}-1|\leq
 \begin{cases}
 C Q^{-10e},&\mathfrak p\mid u,\\
 C Q^{-1-10e},&\mathfrak p\nmid u.
 \end{cases}
\]
Choose the usual finite cutoff $P_0(e)$ before splitting these factors.
It concerns the nonprincipal dynamic correction and is fully compatible
with Lemma~\ref{lem:principal-nonzero}.
All masks generated by the local table are retained in the common
coefficients of Theorems~\ref{thm:can-inverse} and~\ref{thm:can-fourth}.
Finite sector multiplicities and the ideal divisor bound change constants
and small losses only.  This proves the transfer of the full error
decomposition, including slots of zero amplitude.

Inserting a row cardinality into the norm-form contour estimate gives
\eqref{eq:physical-row}; it is the scale identity in \cite[equation~(20.4)]{Q},
with $l_x,l_y,\ell,h$ from \eqref{eq:geometry}.  The term $q/6$ is measured in base $Z$.
The outer-row comparisons are \cite[Lemma~10.6, equations~(10.21)--(10.25)]{Q}.  The principal contour estimates are
\cite[Lemma~10.5, equations~(10.17)--(10.19)]{Q}, with the numerical
substitution in \cite[equation~(20.3)]{Q}.
On the principal absolute contour one can use
$G_{\mathfrak p}=-H_{\mathfrak p}+O(Q^{-7/8})$ and
$H_{\mathfrak p}=1+O(Q^{-363/200})$.  The selected factors are bounded,
the unselected product converges absolutely, and summing each prime
window gives the full bound $\ll Z^{\ell\Rea z}$.

Finally the initial common contour has $s=B+20e$ and
$w=1-a-6e$, so
$\Rea(s+w)\geq1+(B-a)+14e$.  The final dynamic region still has margin
$10e$.    Fixed-degree norm-sheet
integration uses only the finite-seminorm bounds already proved, and
arbitrary integration-by-parts orders give the external tails.  
\end{proof}

\section{Row counting and the high exponent}
\label{sec:rows}

Throughout this section we use the common boundary \(B\) in
\eqref{eq:common-boundary} and assume \(B>7/8\).  Thus
\[
 \Delta=B-\frac78\in(0,1/8],\qquad
 \kappa=2B-1=\frac34+2\Delta\in(3/4,1].
\]
The fixed geometric constants are
\begin{equation}\label{row:constants}
 h=\frac{13}{16},\qquad \ell=\frac16,\qquad
 \alpha=\frac56,\qquad \delta_0=\frac1{50},\qquad
 d_{\min}=\frac1{16}.
\end{equation}
Fix \(0<\zeta<1/48\).  A moderate dyadic row range has
\[
 U=Z^d,\qquad d_{\min}\le d\le h+\zeta.
\]
Write the physical prime scales as
\[
 P_i=Z^{\ell_i}=U^{w_i},\qquad
 w_i=\frac{\ell_i}{d},\qquad
 \ell_i>0,\qquad \sum_{i=1}^{J}\ell_i=\ell.
\]
The slot system is fixed independently of \(Z\); its underlying prime
supports are disjoint before any common mask or row zero extension is
imposed.  Its mesh will be chosen below.

\subsection{The analytic inputs and the pointwise subdivisions}
\label{row:inputs}

We state explicitly the inputs used in the counting argument.  A live
detector bin has
\[
 a>\frac{51}{100},\qquad
 \delta=2a-1\in(\delta_0,\kappa].
\]
The floor bin has \(a=51/100\) and is treated by the total row count.
Proposition~\ref{prop:physical-high}, using the buffered detector and
saturated witnesses of \cite[Lemmas 8.1--8.2 and Proposition 8.3]{Q},
has the following consequences.  For every \(1\le t\le3/2\) and every sufficiently small
prescribed loss \(\eps_{\mathrm w}>0\), each row in a live bin has
centrally normalized inverse and plain witnesses \(M_{u,r}\) and
\(S_{u,m}\), for the same row character, such that
\begin{align}
 &r\le t+O(1/\log U),\qquad
 r+m\ge t-O(1/\log U),\qquad
 0\le m\le\tfrac12+O(\eps_{\mathrm w}),\label{row:witness-lengths}\\
 &|M_{u,r}|^2\gg U^{\delta r-\eps_{\mathrm w}},\qquad
 |S_{u,m}|^2\gg U^{\delta m-\eps_{\mathrm w}}.
 \label{row:witness-spikes}
\end{align}
Their lengths are \(U^r\) and \(U^m\).  There are finitely many character
presentations and \(O((\log U)^2)\) dyadic length pairs.  On each such
subdivision, their remaining rowwise choices are smooth parameters of
common annular profiles.  In particular, the inverse base profile has
the form
\[
 W_{\mathrm{base}}(y)=W_1(y)V_{\le}(U^r y/U^t),
\]
whose fixed logarithmic seminorms are bounded uniformly in the displayed
scales.  Taking the allowed suprema in the smooth parameters costs only
a fixed power of \(1+T\), where \(T\) bounds the retained heights.  All
height assertions are subject to the single cumulative detector
allowance.  
The inequalities in \eqref{row:witness-lengths} also give
\(r\ge t-\tfrac12-O(\eps_{\mathrm w})\).

The required moment estimates, in row base \(U\), are precisely
\begin{align}
 \sum_u\left|M_{u,r}\prod_{i\in I}Q_{u,i}\right|^2
   &\ll U^{1+\eps_{\mathrm m}}(1+T)^A,
 &r+2z&<1,\quad 2r+8z<3,\label{row:moment-inverse}\\
 \sum_u |S_{u,m}|^4\prod_{i\in I}|Q_{u,i}|^2
   &\ll U^{1+\eps_{\mathrm m}}(1+T)^A,
 &2m+6\kappa z&\le1\quad(z>0),\label{row:moment-plain}\\
 \sum_u|M_{u,r}|^2
   &\ll U^{e(r)+\eps_{\mathrm m}}(1+T)^A,
 &e(r)&=\max\{1,(1+5r)/6\}.
 \label{row:moment-amplified}
\end{align}
Here \(z=\sum_{i\in I}w_i\), the two strict inequalities in
\eqref{row:moment-inverse} have fixed positive margins, and the selected
slot lengths in \eqref{row:moment-plain} are below its prescribed mesh.
Every positive plain slot is the complete smooth sum over the good
primes in its annulus, with fixed finite-character coefficients and
only the permitted common masks and original row zero extensions.

The marked inverse estimate admits its larger class of fixed bounded
coefficients and prime lists; the present application uses complete
annuli for both moments.
At \(z=0\), \eqref{row:moment-plain} holds without a restriction on the
bounded nonnegative length \(m\); its actual inducing character must
be nonprincipal after a global norm twist is separated.  For \(z>0\),
the finite inducing character must lie outside the fixed finite family
\(\Theta\).  These are
Theorems~\ref{thm:can-inverse} and~\ref{thm:can-fourth}, and
Lemma~\ref{lem:amplification}, respectively.  The exponent in
\eqref{row:moment-amplified} is uniform as \(r\) approaches one.
The inverse estimates allow their full row families; the plain estimate
is summed over the indicated admissible family.  Positivity of the
moment sums makes each estimate valid on every subset of its row family.

For clarity, the coefficient requirement is preserved when these
estimates are applied to physical slots.  On a fixed presentation write
the common witness character as
\(\lambda_u=\omega_0\xi_u\), with \(\omega_0\) fixed.  Each physical prime
factor can be written, with a fixed smooth profile \(V_i\), as
\begin{equation}\label{row:common-coefficients}
 Q_{u,i}
 =P_i^{-1/2}\sum_{\mathfrak p\notin S}
   \lambda_u(\mathfrak p)\,
   \omega_0(\mathfrak p)^{-1}b_i(\mathfrak p)
   V_i(\Norm\mathfrak p/P_i)
 =P_i^{-1/2}\sum_{\mathfrak p\notin S}
   \xi_u(\mathfrak p)b_i(\mathfrak p)
   V_i(\Norm\mathfrak p/P_i).
\end{equation}
The \(b_i\) are fixed finite combinations of the allowed finite
characters.  Thus the actual prime characters remain finite, and the
relative coefficients are independent of the moving row.  All
reciprocity phases, common masks and original zero extensions remain
in the displayed factors.  If an orientation change is needed, conjugate
the whole product, including both copies of the plain witness, all
selected prime factors and their profiles.  This changes neither its
absolute square nor the relevant exceptional-family condition.
For a retained sixth-power-free physical row, an exceptional finite
inducing character would force every prime outside the fixed set \(S\)
to have valuation divisible by six.  The exceptional physical rows
therefore have bounded norm.  Apart from the principal row, which is
treated separately, they belong to the small-row range below.

The fourth-moment mesh is uniform for
\(3/4\le\kappa\le1\), independently of the eventual fixed slot count;
see Proposition~\ref{prop:can-uniformity} and
\cite[Lemma 18.1]{Q}.  If \(\kappa<1\), its boundary premise holds
because \((1+\kappa)/2=B\) is the common family boundary.  At
\(\kappa=1\) its prime estimate uses absolute counting and requires
no zero-free premise.

We next specify the amplitude sets, including their zero bins.
Fix the external physical parameters and a subset of error slots only
after moving the original fixed row sum to its retained central
contours.  The shared detector includes the finitely many actual
prime-character presentations in \eqref{row:common-coefficients} and
their conjugates, as well as the witness presentations; they all belong
to the common family defining \(B\).  The prime estimate in a detector bin gives, for any
prescribed \(\eps_{\mathrm p}>0\),
\begin{equation}\label{row:prime-upper}
 |Q_{u,i}|\ll U^{\eps_{\mathrm p}}P_i^{\delta/2+O(e)}
\end{equation}
on every main slot; this is the annular logarithmic-derivative estimate
of \cite[Lemma 19.1]{Q}.  Choose an amplitude width
\(\vartheta>0\).  After the slot system is fixed, \(e\) and
\(\eps_{\mathrm p}\) may be reduced so that the right side of
\eqref{row:prime-upper} is at most \(P_i^{\delta/2+\vartheta}\),
for all sufficiently large \(Z\) and every slot.
For a main slot with \(Q_{u,i}\ne0\), set
\[
 g_i=\min\left\{\frac\delta2,\,
       \max\left\{0,\,
          \vartheta\left\lfloor
           \frac{\log|Q_{u,i}|}{\vartheta\log P_i}
          \right\rfloor\right\}\right\}.
\]
Set \(g_i=0\) when \(Q_{u,i}=0\), and also for every error slot.
Then, on a main slot,
\begin{equation}\label{row:amplitude-bins}
 |Q_{u,i}|\le P_i^{g_i+\vartheta},\qquad
 g_i>0\ \Longrightarrow\ |Q_{u,i}|\ge P_i^{g_i}.
\end{equation}

The finitely many vectors \((g_i)\) partition the rows pointwise.
On one such set define
\begin{equation}\label{row:mean-amplitude}
 q=\frac{\sum_i\ell_i g_i}{\ell},
 \qquad 0\le q\le\frac\delta2.
\end{equation}
Errors are included with their assigned zero amplitudes in this
average.  Each pointwise set is further subdivided by the witness
presentation and dyadic lengths.  Within the resulting row sum, the
arithmetic coefficients, the amplitude vector and any selected slot
indices are common.  These pointwise sets bound the already continued integrand.

\begin{lemma}[Selection of positive amplitudes]\label{row:selection}
For \(0\le z\le\ell/d\), a subcollection of whole main slots with
positive amplitude has total length at most \(z\) and satisfies
\begin{equation}\label{row:selection-spike}
 \prod_{i\ {\rm selected}}|Q_{u,i}|^2
 \ge U^{\,2qz-\delta\max_i w_i}.
\end{equation}
If the requested capacity is first reduced by a fixed \(\nu>0\),
the total additional loss in the exponent is at most
\(2q\nu+\delta\max_i w_i\), whenever \(z>\nu\).
\end{lemma}

\begin{proof}
Order the positive \(g_i\) decreasingly and fill a length \(z\)
fractionally in that order.  If they fill \(z\), their initial weighted
average is at least the average \(q\) of all slots.  If they do not,
their entire gain is \(q\ell/d\ge qz\).  At most one terminal slot is
fractional.  Deleting it loses at most
\((\delta/2)\max_iw_i\) in the unsquared exponent.
Equation~\eqref{row:amplitude-bins} proves
\eqref{row:selection-spike}.  Replacing \(z\) by \(z-\nu\) loses a
further \(2q\nu\).  Only positive amplitudes have been used, so the
argument also applies when some or all main factors vanish.
\end{proof}

\subsection{A row bound throughout the moderate range}
\label{row:counts}

\begin{proposition}\label{row:count-proposition}
Given \(\zeta\in(0,1/48)\) and a requested loss \(\eps>0\), a sufficiently
fine fixed physical mesh has the following property.  For every
moderate dyadic range, every retained amplitude set and each witness
subdivision of a live detector bin, one has
\begin{equation}\label{row:cardinality}
 \#\cB\ll U^{R+\eps}(1+T)^A.
\end{equation}
The floor bin admits \(R=1\).  For
\(\alpha\le\delta\le\kappa\le1\), one may take
\begin{equation}\label{row:high-delta}
 R=1-\delta.
\end{equation}
For \(\delta_0<\delta\le\alpha\), put
\begin{equation}\label{row:j-t}
 y=\frac12-\frac q\delta,\qquad
 j=185+(144y-138)\delta,\qquad
 t=1+\frac{6(7+12y)\delta}{j},
\end{equation}
and define
\begin{equation}\label{row:bar-R}
 \overline R=1-\delta+
       \frac{(5-6\delta)(7+12y)\delta}{j}.
\end{equation}
Then \(70\le j\le185\), \(1\le t\le3/2\), and
\begin{equation}\label{row:selected-R}
 R=\overline R+\frac{\Delta}{4}
\end{equation}
is admissible.

The requested loss and the physical mesh can be fixed independently of
the later target twist, uniformly over the displayed ranges of
\(\delta,q,\kappa,d\).  For each fixed arithmetic datum and slot
system, the implied constant, the finite height order \(A\), and the
lower threshold for \(Z\) are fixed.  The remaining detector and
amplitude losses are then chosen sufficiently small, and the retained
heights must satisfy the common detector allowance.
\end{proposition}

\begin{proof}
We first suppress the prescribed witness, moment and rounding losses,
and restore them using fixed margins.

For \(\delta\ge\alpha\), take \(t=3/2\).  The saturated witness lengths
give \(r\ge1-O(\eps_{\mathrm w})\).  Dividing
\eqref{row:moment-amplified} by the inverse spike uses the exact identity
\[
 e(r)-\delta r
 =1-\delta+\delta(1-r)_++(\alpha-\delta)(r-1)_+.
\]
Since \(\alpha-\delta\le0\), this is at most
\(1-\delta+O(\eps_{\mathrm w})\), proving
\eqref{row:high-delta}.  This unmarked estimate applies also when \(B=1\).
The floor bin uses the total row count \(O(U)\).

Now suppose \(\delta_0<\delta<\alpha\), and write
\(x=q/\delta\in[0,1/2]\).  If \(r\ge1\), the same amplified moment,
together with \(r\le t+O(\eps_{\mathrm w})\), gives
\begin{equation}\label{row:long-line}
 R\le L(t):=1-\delta+(\alpha-\delta)(t-1).
\end{equation}
If \(m\ge1/2\), the zero-slot plain moment gives \(R\le1-\delta\).
It remains to treat \(r<1\), \(m<1/2\), when both prospective
capacities are positive.

The inverse capacity is \(z_I=(1-r)/2\).  The actual plain capacity is
\begin{equation}\label{row:plain-capacity}
 z_P=\frac{1-2m}{6\kappa}
     =\frac{1-2m}{9/2+12\Delta}.
\end{equation}
First compare the inverse count with the plain count at the baseline
capacity \(2(1-2m)/9\).  Division by the individual spikes in
\eqref{row:witness-spikes} and the squared prime spike gives the two
affine comparison lines
\[
 A(r)=1-\delta[x+(1-x)r],\qquad
 S_t(r)=1-\delta\{4x/9+(2-8x/9)(t-r)\}.
\]
The replacement of \(m\) by \(t-r\) is legitimate up to
\(O(\eps_{\mathrm w})\): both the actual and baseline plain exponents
decrease with \(m\), since their derivatives are
\[
 -2\delta+\frac{4q}{9/2+12\Delta}<0,\qquad
 -2\delta+\frac{8q}{9}<0.
\]
Choose the inverse estimate when \(A(r)\le S_t(r)\), and the plain
estimate otherwise.  The slopes of \(A\) and \(S_t\) have opposite
signs.  Moreover,
\begin{align*}
 A(3/5)-S_t(3/5)
   &\ge\delta(1/5-14x/45)\ge2\delta/45>0,\\
 S_t(t-1/3)-A(t-1/3)
   &=\delta\{(1-x)(t-1)+5x/27\}\ge0.
\end{align*}
Consequently an inverse selection occurs only for \(r\ge3/5\), while
a plain selection occurs only for \(r\le t-1/3\), hence
\(m\ge1/3-O(\eps_{\mathrm w})\).  In particular,
\begin{equation}\label{row:capacity-bounds}
 z_I\le\frac15,\qquad z_P\le\frac2{27}+O(\eps_{\mathrm w}).
\end{equation}

Reduce a chosen positive capacity by a fixed small \(\nu>0\).
On the inverse side this gives the strict width margins
\[
 1-r-2z\ge2\nu,\qquad
 3-2r-8z
 =4(1-r-2z)+(2r-1)\ge8\nu+\frac15.
\]
The plain capacity is reduced as well.  If \(z_I\le\nu\), use the
unmarked inverse count
\[
 1-\delta r\le1-\delta+2\delta\nu.
\]
If \(z_P\le\nu\), use the unmarked plain count
\[
 1-2\delta m\le1-\delta+6\delta\nu,
\]
because \(1-2m\le6\kappa\nu\le6\nu\).

There is a fixed supply margin for every moderate \(d\):
\begin{equation}\label{row:supply}
 \frac{\ell}{d}\ge\frac{\ell}{h+\zeta}>\frac15.
\end{equation}
Lemma~\ref{row:selection} therefore supplies every requested capacity
in \eqref{row:capacity-bounds}.  Its exponent loss is at most
\[
 2q\nu+\delta\max_iw_i
 \le \nu+\frac{\max_i\ell_i}{d_{\min}}.
\]
Choose \(\nu\) within the requested loss, then require
\(\max_i\ell_i/d_{\min}=16\max_i\ell_i\) to be smaller than the
rounding allowance and the uniform fourth-moment mesh.  The moment contains the selected slots, of total length at most
\(1/5\).
The bounded effective length ranges used to choose its mesh therefore
remain fixed independently of the physical slot count.

Replacing the baseline plain capacity by
\eqref{row:plain-capacity} increases the count exponent by
\begin{align}
 2q(1-2m)
       \left(\frac29-\frac1{9/2+12\Delta}\right)
 &=\frac{24q(1-2m)\Delta}{(9/2)(9/2+12\Delta)}
 \le\frac{\Delta}{4}+O(\eps_{\mathrm w}).
 \label{row:dynamic-cost}
\end{align}
Here \(m\ge1/3-O(\eps_{\mathrm w})\), \(2q\le1\), and
each denominator is at least four.  The estimate is uniform for
\(0\le\Delta\le1/8\).

The smaller comparison line is bounded above by the weighted average
\[
 \frac{(2-8x/9)A(r)+(1-x)S_t(r)}{3-17x/9}
 =1-\delta+\delta p(x)(3/2-t),
 \qquad
 p(x)=\frac{(2-8x/9)(1-x)}{3-17x/9}.
\]
Its \(r\)-coefficient is zero.  On \(0\le x\le1/2\), use the linear
majorant
\begin{equation}\label{row:linear-majorant}
 p(x)\le g(x):=\frac{26-24x}{37},\qquad
 g(x)-p(x)=
 \frac{4(1-2x)(9-14x)}{37(27-17x)}\ge0.
\end{equation}
Put \(L_0=\alpha-\delta+\delta g(x)>0\) and choose
\[
 t=1+\frac{\delta g(x)}{2L_0}.
\]
Then \(1\le t\le3/2\) and
\[
 \delta g(x)(3/2-t)=(\alpha-\delta)(t-1).
\]
The short count and the long count \eqref{row:long-line} are therefore
both bounded by
\[
 1-\delta+
 \frac{(\alpha-\delta)\delta g(x)}
      {2(\alpha-\delta+\delta g(x))}
 +\frac{\Delta}{4},
\]
apart from the requested losses.  Substitution of \(y=1/2-x\) gives
\eqref{row:j-t}--\eqref{row:selected-R}, with \(j=222L_0\).
The bounds \(70\le j\le185\) follow directly from
\(0\le y\le1/2\), \(0\le\delta\le5/6\).  At \(\delta=\alpha\),
the displayed formulas have \(t=3/2\) and
\(\overline R=1-\delta\); the already proved unmarked case gives
the stronger bound \eqref{row:high-delta}.

To restore the losses, first fix the moment loss and \(\nu\), then
the physical mesh and a fixed slot system.  The supply gap and the
inverse widths above are fixed, and the fourth-moment mesh is uniform
for \(\kappa\in[3/4,1]\).  Next reduce the witness, detector and
amplitude losses.  The inequalities \(\delta\ge1/50\) and \(j\ge70\)
make all length perturbations and rational expressions uniform;
the near-zero capacities have already been covered by unmarked
moments.  The finitely many amplitude vectors and the
\(O((\log U)^2)\) witness choices cost at most a further arbitrarily
small power.  The permitted common-profile calculus supplies the
fixed height factor in \eqref{row:cardinality}.  This proves the
proposition with the stated order of choices.
\end{proof}

\subsection{A fixed numerical margin}
\label{row:certificate}

In the present geometry,
\[
 C(s)=s-\frac{11}{16},\qquad C(7/8)=\frac3{16},
 \qquad C(B)-C(7/8)=\Delta.
\]
For the compensated physical expression, the exponent associated with
a pointwise row count \(U^{R+\eps}\) is given by
\eqref{eq:physical-row}.  Relative to \(C(7/8)\), it is
\begin{equation}\label{row:physical-E}
 E(d;R,q)
 =-\frac1{48}+\frac{2\delta}{3}+\frac q6-h(1-R)
  +(d-h)\left(R+\frac\delta2-\frac{17}{50}\right).
\end{equation}
Thus the central contribution of the original fixed bin is bounded
using the supremum of \(E\) over its pointwise sets, with the prescribed
real losses and a fixed polynomial height factor.  In particular,
\(q/6=q\ell\) is a base-\(Z\) exponent.

\begin{lemma}\label{row:certificate-lemma}
Let \(c_*=49/301920\).  Before the adjustable real losses, every
moderate bin has an admissible row exponent from
Proposition~\ref{row:count-proposition} such that
\begin{equation}\label{row:uniform-saving}
 E(d;R,q)-\Delta\le-c_*+\zeta.
\end{equation}
This is uniform for \(B\in(7/8,1]\).
\end{lemma}

\begin{proof}
For \(0\le\delta\le5/6\) and \(0\le y\le1/2\), let
\(\overline E=E(h;\overline R,(1/2-y)\delta)\), with \(j\) and
\(\overline R\) as in \eqref{row:j-t}--\eqref{row:bar-R}.
Expansion gives the exact identity
\begin{equation}\label{row:square-certificate}
 \boxed{\;
 1632j(-\overline E)
 =[(204+170y)\delta-79]^2+49
  +68y\delta[37+(48+151y)\delta].
 \;}
\end{equation}
Every term on the right is nonnegative, and \(j\le185\).  Hence
\[
 -\overline E\ge\frac{49}{1632\cdot185}
 =c_*>\frac1{10000}.
\]
At \(d=h\), the capacity cost \(\Delta/4\) increases \(E\) by
\(h\Delta/4\).  Therefore
\begin{equation}\label{row:relative-saving}
 E(h;\overline R+\Delta/4,q)-\Delta
 \le-c_*-\frac{51}{64}\Delta.
\end{equation}

For \(\alpha\le\delta\le1\), choose \(R=1-\delta\) and use
\(q\le\delta/2\).  Then
\[
 E(h;1-\delta,q)
 \le-\frac1{48}-\frac\delta{16}\le-\frac7{96}.
\]
In the floor bin, \(R=1\), \(\delta=\delta_0\) and
\(q\le\delta_0/2\) give
\[
 E(h;1,q)\le-\frac1{48}+\frac{3\delta_0}{4}
 =-\frac7{1200}.
\]

It remains to extend these bounds in \(d\).  Since
\(0<g(x)\le1\), the selected exponent satisfies
\[
 1-\delta\le\overline R\le1-\delta/2.
\]
Using \(\Delta/4\le1/32\), all the upper exponents just used,
including the upper-\(\delta\) and floor cases, obey
\begin{equation}\label{row:positive-slope}
 \frac4{25}\le R+\frac\delta2-\frac{17}{50}<1.
\end{equation}
The estimates at \(h\) therefore hold for every smaller moderate
\(d\), and their extension to \(h+\zeta\) costs at most \(\zeta\).
Both \(7/96\) and \(7/1200\) exceed \(c_*\), proving
\eqref{row:uniform-saving}.
\end{proof}

\subsection{The two outer ranges}
\label{row:outer}

The outer ranges use the full correction before its central
main/error decomposition.  Let \(\mathcal R(U;Z)\) denote the
contribution of nonprincipal physical rows
\(U\le\Norm u<2U\), and let \(\mathcal H_{u,Z}\) denote their full
correction.  The quotient-free bounds in the physical expansion are,
with \(z_0=17/50\),
\begin{align}
 |\mathcal H_{u,Z}(B+e+it,\,1/2+iv,\,z_0+iw)|
 &\ll Z^{\ell z_0+\eps}(\Norm u)^\eps
        (1+|t|+|v|+|w|)^A,\label{row:small-correction}\\
 |\mathcal H_{u,Z}(2+it,\,2+iv,\,z_\infty+iw)|
 &\ll Z^{\ell z_\infty+\eps}(\Norm u)^\eps
        (1+|t|+|v|+|w|)^A,
 \qquad z_\infty>2.
 \label{row:large-correction}
\end{align}
These bounds retain every local factor and are valid at its zeros.
The contours can thus be shifted for each original dyadic sum.

On the first lines the reciprocal remains on the common global line
\(\Rea s=B+e\).  The nonprincipal numerator contributes at most
\(U^{3/5+\eps}\) times a fixed height polynomial.  Counting the rows
and including their factor \((\Norm u)^{-z_0}\) gives the row power
\[
 1+\frac35-\frac{17}{50}=\frac{63}{50}.
\]
The \(Z\)-power outside this row count is
\(C(B)+h(z_0-1/6)-l_y/2+e\), where \(l_y=23/48\).
This proves the outer-row estimate
\cite[Lemma 10.6]{Q} in the present normalization:
\begin{equation}\label{row:small-dyad}
 |\mathcal R(U;Z)|
 \ll Z^{C(B)-79/800+e+\eps}U^{63/50+\eps}.
\end{equation}
The fixed tests integrate the height polynomial.  Summing the dyads
\(1\le U\le Z^{d_{\min}}\), and decreasing the preliminary powers,
gives
\begin{equation}\label{row:small-saving}
 \sum_{U\le Z^{1/16}}|\mathcal R(U;Z)|
 \ll Z^{C(B)-1/50+e+\eps},
 \qquad
 -\frac{79}{800}+\frac{63}{50}\frac1{16}
 =-\frac1{50}.
\end{equation}
This estimate includes bounded exceptional rows with nonprincipal numerator.
It also applies at \(B=1\), when the reciprocal line is in the region
of absolute convergence.

On the absolute lines in \eqref{row:large-correction}, the scalar
\(L\)-factors are absolutely bounded.  With
\[
 B_0=\frac{l_x}{2}+1+l_y=\frac{53}{32},
 \qquad l_x=\frac{17}{48},
\]
the same dyadic calculation gives
\begin{equation}\label{row:large-dyad}
 |\mathcal R(U;Z)|
 \ll Z^{B_0+h z_\infty+\eps}U^{1+\eps-z_\infty}.
\end{equation}
Choose the explicit fixed line \(z_\infty=1+3/\zeta\).
For \(\eps<z_\infty-1\), geometric summation over
\(U>Z^{h+\zeta}\) yields
\begin{align}
 \sum_{U>Z^{h+\zeta}}|\mathcal R(U;Z)|
 &\ll Z^{B_0+(h+\zeta)(1+\eps)-\zeta z_\infty+\eps}\notag\\
 &=Z^{-17/32+(1+h+\zeta)\eps}.
 \label{row:large-saving}
\end{align}
As \(C(B)\ge3/16\), its saving relative to \(C(B)\) is at least
\(23/32-O(\eps)\).

Equations~\eqref{row:uniform-saving},
\eqref{row:small-saving} and~\eqref{row:large-saving} cover every
nonprincipal physical row.  Choose
\(\zeta<c_*/4\), then choose the requested row-count loss and its
uniform moment mesh within fixed fractions of \(c_*\).
After fixing the slot system, choose the remaining real detector,
amplitude and correction losses so that their total exponent cost is
less than \(c_*/4\).  This leaves a fixed positive nonprincipal saving,
independent of \(B\) and of the later target.  The finite height orders
and constants may depend on that target and the fixed slot system.
Their absorption, the principal remainder and the common
normalization are handled in the final choice of heights.

\section{Mellin continuation}\label{sec:closure}

\subsection{Real parameters and heights}\label{subsec:parameter-order}

The choices are made in the following order. Steps 1--5 determine a
common real saving; Steps 6--9 specialize the estimates to a target.
\begin{enumerate}[label=\textup{\arabic*.},leftmargin=*,itemsep=5pt]
\item \textbf{Field and infinity types.}
Fix $K$ and the three types $0,\omega_\infty,\bar\omega_\infty$.
The final excluded set $S$ is chosen at Step 6, after the target's
finite conductor is known.

\item \textbf{Geometric margins.}
Fix $c=10^{-4}$, $d_{\min}=1/16$, and
$0<\zeta<\min\{1/96,c/16\}$. The certificate and the outer-row
estimates in Section~\ref{sec:rows} leave the required positive reserves.

\item \textbf{Moment losses and mesh.}
Choose the moment exponent losses and capacity decrements within small
fixed fractions of $c$. Proposition~\ref{prop:can-uniformity} then gives
a prime-slot mesh, uniform on $3/4\le\kappa\le1$. These choices precede
the number of slots because the capacity-error coefficients depend on
total lengths.

\item \textbf{The finite slot system.}
Take $\ell_i=1/(6J_0)$ and choose $J_0$ so that
$16\max_i\ell_i$ is smaller than both the mesh and the rounding reserve.
Fix nonnegative nonzero tests supported in the disjoint intervals
\[
 \left(1+\frac{2i-1}{2J_0+1},1+\frac{2i}{2J_0+1}\right),
 \qquad 1\le i\le J_0.
\]
This fixes the analytic slot system and
$\mu_0=7/(96J_0)>0$. The low exponent for this system is the one in
Proposition~\ref{prop:physical-low}.

\item \textbf{Remaining real losses.}
Choose the amplitude bins, detector width, contour shifts, and local
correction losses. They may depend on $J_0$ and are taken within the
remaining nonprincipal reserve and the principal margins
$23/960$, $13/9600$, and $\mu_0$. The principal losses are chosen
separately from the earlier moment losses. Set
$z_\infty=1+3/\zeta$. Reserve a common number $m_*>0$ for the final
height estimate.

\item \textbf{Target and arithmetic data.}
Choose the later finite-order twist and hence a target in $\cF_\omega$.
Remove its norm twist and use the permitted conjugate, writing $\psi$
for that representative. Fix its finite sectors and final $S$, including
the cutoff $P_0(e)$. This same $S$ is used in the auxiliary function,
residue scalar, and Mellin signal.

\item \textbf{Internal analytic orders.}
Fix the internal seminorm and Sobolev orders and the resulting finite
height degrees $A_\psi,B_\psi$. Propositions~\ref{prop:can-uniformity}
and~\ref{prop:physical-high} give, in a positive target-dependent
height range,
\begin{equation}\label{eq:before-height}
 |J_\psi(Z)-f_\psi(Z)|
 \ll_{\psi,N}Z^{C(B)-2m_*}(1+T)^{A_\psi}
                       +Z^{B_\psi}T^{-N}.
\end{equation}
The degrees in this estimate are fixed before the external tail
order $N$.

\item \textbf{Height and external tail.}
Choose $\tau_\psi>0$ within that height range and with
$\tau_\psi\le m_*/[2(A_\psi+1)]$, and set $T=Z^{\tau_\psi}$.
Then choose $N$ with
$N\tau_\psi>B_\psi-3/16+m_*$.

\item \textbf{Constants and thresholds.}
Choose the implied constants and a lower threshold for $Z$ for these
fixed data and orders. They may depend on the target. The common
real number $m_*$ was fixed at Step 5.
\end{enumerate}

Since $C(B)\ge3/16$, both terms in \eqref{eq:before-height} are now
$\ll_\psi Z^{C(B)-m_*}$. Thus
\begin{equation}\label{eq:fixed-high}
 |J_\psi(Z)-f_\psi(Z)|\ll_\psi Z^{C(B)-\eta_0},
 \qquad \eta_0=m_*>0.
\end{equation}
The saving is common to the later targets and independent of
$\Delta=B-7/8$. Conjugation and vertical translation restore the
original target after its chosen representative has been treated.

\subsection{Family closure}

We use the common-signal continuation principle of
\cite[Proposition~2.1]{Q}.  The formulation below retains the maximum
of the low exponent and the fixed high-side exponent, and applies it to
the fixed-type family of this paper.

\begin{lemma}\label{lem:mellin-closure}
Suppose \eqref{eq:fixed-high} holds for every representative of the fixed
family with a common $\eta_0>0$, and the same fixed probe satisfies
\eqref{eq:all-low} for every $\eps>0$.  Then every member of the family is
zero-free in
\begin{equation}\label{eq:improved-halfplane}
 \Rea s>\max\{7/8,B-\eta_0\}.
\end{equation}
Consequently $B\leq7/8$.
\end{lemma}
\begin{proof}
Combining the two estimates gives, for every $\eps>0$,
\[
 |f_\psi(Z)|\ll_{\psi,\eps}
     Z^{C(7/8)+\eps}+Z^{C(B)-\eta_0}.
\]
As $Z\downarrow0$, moving the defining line in \eqref{eq:signal}
arbitrarily far to the right gives decay faster than any fixed power.
The Gaussian supplies absolute vertical convergence on each such line.
Therefore
\begin{equation}\label{eq:mellin-continuation}
 F_\psi(s)=\int_0^\infty f_\psi(Z)Z^{-C(s)}\,\frac{dZ}{Z}
\end{equation}
is holomorphic in \eqref{eq:improved-halfplane}.  To verify local uniform
convergence, choose $\eps$ smaller than the distance of a given compact
set from $\Rea s=7/8$.  

Mellin inversion on an initial absolute line identifies
\eqref{eq:mellin-continuation} with
\[
 e^{(s-5/6)^2}\frac{H_\psi^S(s)}{L_K^S(s,\psi)}.
\]
The Gaussian and $H_\psi^S$ are nonzero in the asserted half-plane,
by Lemma~\ref{lem:principal-nonzero}.  A zero of $L_K^S(s,\psi)$ would
create a pole, contrary to holomorphy.  At a norm-character pole the reciprocal is holomorphic and vanishes.  Deleted Euler factors
are nonzero in $\Rea s>0$.

If $B>7/8$, the number $\max\{7/8,B-\eta_0\}$ is strictly smaller
than $B$.  By the definition \eqref{eq:common-boundary}, some member of
$\cF_\omega$ has a zero strictly to its right.  This contradicts
\eqref{eq:improved-halfplane} and proves $B\leq7/8$.
\end{proof}

\begin{proof}[Proof of Theorem~\ref{thm:main}]
Given $F$ and $\psi$, take $K=F(\zeta_{12})$ and
$\omega=\psi\circ\Norm_{K/F}$.  The arithmetic, reflection, and canonical
arguments above apply over $K$.  If its common boundary exceeded $7/8$,
the row estimates and the fixed choices just made would give
\eqref{eq:all-low} and \eqref{eq:fixed-high} for the same probes.
Lemma~\ref{lem:mellin-closure} rules this out.  In particular
$L_K(s,\omega)$ is zero-free for $\Rea s>7/8$.
Lemma~\ref{lem:descent} gives the assertion over $F$, including its
possible norm-character pole and any additionally deleted Euler factors.
\end{proof}

\begin{remark}
The maximum in \eqref{eq:improved-halfplane} is necessary when
$B-7/8<\eta_0$: the low estimate still contributes
$Z^{C(7/8)+\eps}$.  
\end{remark}

\section{Weil representations and applications}\label{sec:weil}

We deduce Theorem~\ref{thm:weil-main} and its applications from
Theorem~\ref{thm:main} by unitary Brauer induction and pole counting.

\subsection{Unitary induction and Euler factors}\label{subsec:unitary-induction}

We use the classical global Weil group, with its local maps and its
finite-index subgroups $W_E$ for finite extensions $E/F$, as in
\cite[Chapters~1--2]{Langlands} and
\cite[Section~3.1]{Johansson}.    A finite-dimensional continuous representation
factors through a relative Weil group $W_{K/F}$ for some finite
Galois extension $K/F$; this is the no-small-subgroups argument recorded
in \cite[proof of Proposition~7]{Johansson}.  We retain the norm
convention from the introduction:
\begin{equation}\label{eq:weil-norm-compatibility}
 \mathcal N_F|_{W_E}=\mathcal N_E,\qquad
 \nu_{F,t}|_{W_E}=\nu_{E,t}.
\end{equation}
Indeed a Frobenius at a prime of residue degree $f$ maps to the
$f$th power of the downstairs Frobenius, and its norm is
$(\Norm\mathfrak p)^f$.  

\begin{lemma}[Unitary Brauer induction]\label{lem:unitary-brauer}
Let $r$ be a finite-dimensional continuous unitary representation of
$W_F$.  There are finite extensions $E_j/F$, continuous unitary
Hecke characters $\psi_j$ of $E_j$, and integers $n_j$ such that
\begin{equation}\label{eq:weil-brauer}
 [r]=\sum_{j=1}^m n_j
  \left[\operatorname{Ind}_{W_{E_j}}^{W_F}\psi_j\right].
\end{equation}
This is an equality in the Grothendieck group of semisimple
representations.  
\end{lemma}
\begin{proof}
Apply \cite[Lemma~2.2]{Langlands} to a relative Weil group through
which $r$ factors.    Every Hecke
quasicharacter has a unique expression
$\psi=\psi^0\mathcal N_E^a$, where $\psi^0$ is unitary and
$a\in\RR$: its absolute value is trivial on the compact norm-one
idele class group and therefore is a real power of the norm.
The tensor-induction identity
\cite[Lemma~2.3]{Langlands} and
\eqref{eq:weil-norm-compatibility} give
\[
 \operatorname{Ind}_{W_E}^{W_F}\psi
   =\left(\operatorname{Ind}_{W_E}^{W_F}\psi^0\right)
      \otimes\mathcal N_F^a.
\]
Induction from a finite-index subgroup preserves unitarity: use the
sum of the invariant inner products on a set of coset representatives.
Thus every irreducible constituent of $\operatorname{Ind}\psi^0$
is unitary.  If $a\ne0$, none of its twists by $\mathcal N_F^a$
is unitarizable.  To see this, choose $g$ with $\mathcal N_F(g)>1$;
all eigenvalues of such a twist at $g$ have absolute value
$\mathcal N_F(g)^a\ne1$.

The Grothendieck group is freely generated by irreducible classes.
Project its underlying abelian group onto the classes of the
unitarizable irreducibles.  This projection fixes $[r]$ and discards
all summands with $a\ne0$.  The remaining inducing characters have
$a=0$, which proves the required refinement of Langlands's lemma.
\end{proof}

The standard induction identity for local $L$-factors, including
ramified primes, gives
\begin{equation}\label{eq:weil-L-induction}
 L_F\!\left(s,\operatorname{Ind}_{W_E}^{W_F}\psi\right)
   =L_E(s,\psi).
\end{equation}
See \cite[Chapter~20, pp.~211--213]{Langlands}.  One can read the
finite-place calculation directly on inertia invariants: Frobenius
permutes the cosets in cycles, and a cycle of residue degree $f$
contributes the determinant
$\det(1-\psi(\operatorname{Frob}_{\mathfrak P})
 q_{\mathfrak p}^{-fs}\mid V_\psi^{I_{\mathfrak P}})$.
It is $1$ when these invariants vanish.  Since
$q_{\mathfrak P}=q_{\mathfrak p}^f$, its reciprocal is the
upstairs Euler factor.  Consequently \eqref{eq:weil-brauer} implies
\begin{equation}\label{eq:weil-L-product}
 L_F(s,r)=\prod_{j=1}^m L_{E_j}(s,\psi_j)^{n_j}.
\end{equation}
This identity first holds in the absolute half-plane and then defines
the usual meromorphic continuation.  

If $S$ is a finite set of finite primes of $F$ and $S_E$ is the set
above it, induction also gives
$L_F^S(s,\operatorname{Ind}\psi)=L_E^{S_E}(s,\psi)$.
Moreover the eigenvalues of Frobenius on the inertia invariants of
any unitary representation have absolute value one.  Thus
\begin{equation}\label{eq:weil-deleted-factors}
 L_F^S(s,r)=L_F(s,r)\prod_{\mathfrak p\in S}
 \det\!\left(1-r(\operatorname{Frob}_{\mathfrak p})
 q_{\mathfrak p}^{-s}\bigm|V_r^{I_{\mathfrak p}}\right)
\end{equation}
is multiplication by a holomorphic nowhere-zero function on
$\Rea s>0$.  Deleting these factors changes neither a zero nor a
pole in our half-plane.

\subsection{The pole divisor}\label{subsec:weil-divisor}

\begin{proposition}[Transfer of a common Hecke half-plane]
\label{prop:weil-halfplane-transfer}
Let $0<\beta<1$.  Suppose every continuous unitary Hecke character
over every number field has a zero-free $L$-function on
$\Rea s>\beta$.  Then every finite-dimensional continuous unitary
Weil representation $r$ has no zeros there.  Away from the points
$1+it$ it is holomorphic, and at those points
\[
 \operatorname{ord}_{1+it}L_F(s,r)
    =-\dim\operatorname{Hom}_{W_F}(\nu_{F,t},r).
\]
The same assertion holds for $L_F^S(s,r)$.
\end{proposition}
\begin{proof}
Use \eqref{eq:weil-L-product} with the unitary inducing characters of
Lemma~\ref{lem:unitary-brauer}.  Each factor is zero-free in the
specified half-plane.  By the Hecke pole classification
\cite{Tate}, its only possible pole is a simple pole at $1+it$,
and this occurs precisely when $\psi_j=\nu_{E_j,t}$.
At all other points every factor and its reciprocal are holomorphic
and nonzero.

At $1+it$ take meromorphic orders.  The multiplicity of a fixed
irreducible character is additive on the Grothendieck group of
semisimple representations, and Frobenius reciprocity gives
\begin{align}
 \operatorname{ord}_{1+it}L_F(s,r)
 &=-\sum_j n_j\ind_{\{\psi_j=\nu_{E_j,t}\}} \notag\\
 &=-\sum_j n_j\dim\operatorname{Hom}_{W_{E_j}}
        (\nu_{F,t}|_{W_{E_j}},\psi_j) \notag\\
 &=-\dim\operatorname{Hom}_{W_F}(\nu_{F,t},r).
 \label{eq:weil-Frobenius-order}
\end{align}
The last dimension is nonnegative, so the displayed order is nonpositive.  The deleted-factor assertion follows from
\eqref{eq:weil-deleted-factors}.
\end{proof}

\begin{proof}[Proof of Theorem~\ref{thm:weil-main}]
Apply Proposition~\ref{prop:weil-halfplane-transfer} with
$\beta=7/8$ and Theorem~\ref{thm:main}.
\end{proof}

Only finitely many $t$ have $m_t(r)>0$, and their multiplicities
sum to at most $\dim r$.  An equivalent form of the result is that
\begin{equation}\label{eq:weil-regularized}
 L_F^S(s,r)\prod_{t:m_t(r)>0}(s-1-it)^{m_t(r)}
\end{equation}
is holomorphic and nowhere zero on $\Rea s>7/8$.

\begin{remark}[Completed $L$-functions]\label{rem:weil-actual-and-completed}

The completed Weil $L$-function has the same zeros and poles as its
finite part on $\Rea s>0$.  Indeed the unitary archimedean factors
are products of $\Gamma_{\RR}(s+\epsilon+it)$, with
$\epsilon\in\{0,1\}$, and $\Gamma_{\CC}(s+a+it)$, with
$a\geq0$; they and any conductor exponential are holomorphic and
nonzero there.  This follows from the local formulas of
\cite[Chapter~20]{Langlands}.  
\end{remark}

\subsection{Artin twists and weight-one forms}\label{subsec:Artin-applications}

\begin{corollary}\label{cor:Artin}
Let $\rho:G_F\longrightarrow\mathrm{GL}_d(\CC)$ be a continuous
finite-image representation.  Then $L_F(s,\rho)$ is zero-free on
$\Rea s>7/8$, and its only possible pole there is at $s=1$, of
order $\dim\rho^{G_F}$.  In particular every nontrivial irreducible
Artin $L$-function is holomorphic and nowhere zero in this half-plane.
For any continuous unitary Hecke character $\psi$ of $F$,
Theorem~\ref{thm:weil-main} also applies to $\rho\otimes\psi$,
with pole multiplicities
$\dim\operatorname{Hom}_{W_F}(\nu_{F,t},\rho\otimes\psi)$.
\end{corollary}
\begin{proof}
Average an inner product over the finite image of $\rho$ and pull
the representation back to $W_F$.  No $\nu_{F,t}$ with $t\ne0$
has finite image, so only the trivial character can contribute a
pole to the untwisted representation.  Tensoring by a unitary
character preserves unitarity, and the general assertion follows.
\end{proof}

The finite-image assertion needs only the finite-order subcase of
Theorem~\ref{thm:main} over all number fields: ordinary Brauer
induction uses finite-order characters.  The infinite-order twists
and the general Weil statement use its full Hecke scope.  More
explicitly, if
$[\rho]=\sum_j n_j[\operatorname{Ind}_{G_{E_j}}^{G_F}\chi_j]$
is a Brauer decomposition, then the tensor-induction identity gives
\begin{equation}\label{eq:Artin-Hecke-twist}
 L_F(s,\rho\otimes\psi)=
  \prod_j L_{E_j}\!\left(s,
       \chi_j(\psi\circ\Norm_{E_j/F})\right)^{n_j}.
\end{equation}
The source of the induction theorem is \cite{Brauer}; the Weil and
$L$-function compatibilities are those of
\cite[Lemma~2.3 and Chapter~20]{Langlands}.
  If $\rho$ is irreducible of dimension at
least two, any such twist has no one-dimensional constituent and
is holomorphic as well as nonzero on the asserted half-plane.

\begin{corollary}\label{cor:weight-one}
For every normalized cuspidal weight-one newform $f$ over $\QQ$,
its standard $L$-function is holomorphic and nowhere zero for
$\Rea s>7/8$.
\end{corollary}
\begin{proof}
Deligne--Serre attach a finite-image irreducible representation
$\rho_f:G_{\QQ}\to\mathrm{GL}_2(\CC)$ and prove
$L(s,f)=L_{\QQ}(s,\rho_f)$, including the bad finite Euler
factors \cite[Theorems~4.1 and~4.6]{DS}.  Apply
Corollary~\ref{cor:Artin}.
\end{proof}

\subsection{Effective Chebotarev}\label{subsec:chebotarev}

Let $L/F$ be a finite Galois extension with group $G$, and let $C$ be
one conjugacy class. Write $c_C=|C|/|G|$,
$n_L=[L:\QQ]$, $D_L=|\operatorname{disc}(L)|$, and
$\mathcal D_L=\log(eD_L)+n_L$. Let $\pi_C(x)$ count the unramified
prime ideals $\mathfrak p$ of $F$ with $q_{\mathfrak p}\le x$ and
$\operatorname{Frob}_{\mathfrak p}\in C$.

\begin{corollary}[Prime ideals with a prescribed Frobenius class]
\label{cor:chebotarev}
For every $0<\eps<1/8$ and $x\ge2$,
\begin{equation}\label{eq:chebotarev-count}
 \pi_C(x)=c_C\operatorname{Li}(x)
       +O_\eps\!\left(c_C\mathcal D_Lx^{7/8+\eps}\right).
\end{equation}
The implied constant is effectively computable and independent of
$L,F,C$. There is an unramified prime ideal in the class $C$ with
\begin{equation}\label{eq:chebotarev-least}
 q_{\mathfrak p}\ll\mathcal D_L^8,
\end{equation}
where the implied constant is absolute and effectively computable.
\end{corollary}

\begin{proof}
We use the classical cyclic reduction and the Hecke explicit formula,
in the form recorded in \cite[Section~3, Lemmas~4--5 and Proposition~6]{DKN}.
Here are the estimates with the zero boundary supplied by
Theorem~\ref{thm:main}.

Choose $g\in C$, put $H=\langle g\rangle$ and $E=L^H$, and regard
$\chi\in\widehat H$ as a finite-order Hecke character of $E$.
Finite-group induction gives
$\ind_C=c_C\sum_{\chi\in\widehat H}\bar\chi(g)
\operatorname{Ind}_H^G\chi$. Thus the associated logarithmic derivative is
\begin{equation}\label{eq:chebotarev-logder}
 D_C(s)=-c_C\sum_{\chi\in\widehat H}
          \bar\chi(g)\frac{L_E'}{L_E}(s,\chi).
\end{equation}
Its coefficient at an unramified prime power is
$\ind_C(\operatorname{Frob}_{\mathfrak p}^{\,m})\log q_{\mathfrak p}$.
At a ramified prime it is the corresponding inertia average, of
absolute value at most $\log q_{\mathfrak p}$.
For $A(\chi)=D_E\Norm\mathfrak f_\chi$, conductor--discriminant and
$[L:\QQ]=|H|[E:\QQ]$ give
$\sum_\chi\log A(\chi)=\log D_L$ and
$\sum_\chi[E:\QQ]=n_L$.

All nontrivial zeros of these Hecke functions have real part at most
$7/8$. The local zero count and logarithmic-derivative expansion in
\cite[Lemma~5, equations~(3.18) and~(3.21)]{DKN} therefore give,
uniformly for $7/8+\eps/2\le\sigma\le2$,
\begin{equation}\label{eq:chebotarev-line}
 D_C(s)-\frac{c_C}{s-1}
 \ll_\eps c_C\bigl(\log D_L+n_L\log(|t|+3)\bigr),
 \qquad s=\sigma+it.
\end{equation}
Apply truncated Perron summation at a half-integer $x$ and height
$T=x^2$, and move the line to $7/8+\eps/2$.
The Dirichlet coefficient at the integer $n$ is bounded in absolute
value by $[F:\QQ]\Lambda(n)$, so the truncation error is
$O([F:\QQ]\log^2(3x))$. The only crossed pole contributes $c_Cx$.
The new vertical integral is
$O_\eps(c_C\mathcal D_Lx^{7/8+\eps/2}\log^2(3x))$;
the horizontal integrals are smaller. Replacing $x$ by the adjacent
half-integer preserves the prime-power sum and changes the main term
by $O(1)$. Absorbing the logarithms gives the same estimate with
exponent $7/8+\eps$.

The contribution of powers $m\ge2$ is
$O([F:\QQ]x^{1/2}\log^2(3x))$.
At a ramified prime, the discriminant exponent is at least
$|G|/2$, and hence
$\sum_{\mathfrak p\,\mathrm{ram}}\log q_{\mathfrak p}
\le2\log D_L/|G|$.
Their prime-power contribution is
$O(|G|^{-1}\log D_L\log(3x))$.
Both fit the stated error since
$[F:\QQ]=n_L/|G|\le c_Cn_L$.
After removing these terms, partial summation proves
\eqref{eq:chebotarev-count}.

For the least-prime bound, fix a nonnegative nonzero
$W\in C_c^\infty((1,2))$ once and for all and write
$a_W=\widehat W(1)>0$.
Mellin inversion of \eqref{eq:chebotarev-logder}, followed by a shift
to $\Rea s=-1/4$, gives the annular explicit formula
\begin{align}\label{eq:chebotarev-smooth}
 \sum_{\mathfrak p,m\ge1}
   \vartheta_C(\mathfrak p^m)\log q_{\mathfrak p}
          W(q_{\mathfrak p}^m/x)
 ={}&c_Ca_Wx
  -c_C\sum_{\chi\in\widehat H}\bar\chi(g)
       \sum_{\rho_\chi}x^{\rho_\chi}\widehat W(\rho_\chi)
  +O_W(c_C\mathcal D_L).
\end{align}
Here $\vartheta_C$ is the inertia-average coefficient just defined,
and $\rho_\chi$ runs over nontrivial zeros with multiplicity.
This is the contour argument of \cite[Proposition~6]{DKN} with the
entire Mellin transform of an annular test. The residue at zero has
size $O_W([E:\QQ])$ for each character, and the functional equation
bounds the remaining line by
$O_W(x^{-1/4}(\log A(\chi)+[E:\QQ]))$.

Uniformly for $0\le\Rea\rho\le1$,
$\widehat W(\rho)\ll_W(1+|\Im\rho|)^{-2}$.
The local zero count therefore implies
$\sum_{\rho_\chi}|\widehat W(\rho_\chi)|
\ll_W\log A(\chi)+[E:\QQ]$.
The zero sum in \eqref{eq:chebotarev-smooth} is consequently
$O_W(c_C\mathcal D_Lx^{7/8})$.
Removing ramified primes and higher powers as above, using
$x^{1/2}\log^2(3x)\ll x^{7/8}$, gives
\[
 \sum_{\substack{\mathfrak p\,\mathrm{unramified}\\
                  \operatorname{Frob}_{\mathfrak p}\in C}}
 \log q_{\mathfrak p}\,W(q_{\mathfrak p}/x)
  =c_Ca_Wx+O_W(c_C\mathcal D_Lx^{7/8}).
\]
For $x$ a sufficiently large absolute multiple of $\mathcal D_L^8$
this sum is positive. Some required prime then has norm between $x$
and $2x$, proving \eqref{eq:chebotarev-least}. All constants in these
contour and zero-count estimates are effective; the weight is fixed.
\end{proof}

\subsection{Geometric and dihedral consequences}\label{subsec:potentially-abelian}

\begin{remark}[Pure potentially abelian realizations]\label{cor:potentially-abelian}\label{subsec:CM}\label{cor:CM}
Let $V$ be a semisimple geometric $\ell$-adic representation over $F$
whose restriction to some finite extension is a sum of characters.
Here geometric means unramified outside finitely many primes and de
Rham above $\ell$. Fix compatible coefficient embeddings, and suppose
its good Euler roots are pure of weight $w\in\ZZ$, with
$|\alpha_{\mathfrak p,j}|=q_{\mathfrak p}^{w/2}$ for every complex
embedding. For $S$ containing the ramified primes and those above
$\ell$, the algebraic Weil correspondence gives a parameter $R$ with
these good factors \cite[Theorem~9 and Corollary~10]{Johansson}.
Its normalization $r_0=R\otimes\mathcal N_F^{-w/2}$ is unitary: after
a finite extension, purity fixes the absolute value of every character
to be $\mathcal N^{w/2}$, and an invariant inner product descends by
averaging over the finite cosets \cite[Lemma~14(1)]{Johansson}.
Consequently
\begin{equation}\label{eq:geometric-L-shift}
 L_F^S(s,V)=L_F^S(s-w/2,r_0),\qquad
 L_F^S(s,V)\ne0\quad(\Rea s>w/2+7/8).
\end{equation}
The only possible pole there is at $1+w/2$, with order
$\dim\operatorname{Hom}_{W_F}(\mathcal N_F^{w/2},R)$.
Indeed an algebraic norm constituent $\mathcal N_F^{w/2+it}$ has
trivial angle at a complex local Weil subgroup, so its two integral
exponents agree; this forces $t=0$ and $w$ even. In particular the
function is holomorphic in this half-plane for odd $w$.

This includes pure CM motives in the category generated by CM abelian
varieties and Artin motives using Hodge cycles. Their full local Weil
realization gives the same assertion for the complete $L$-function
\cite[Corollaries~7.6, 7.9 and Proposition~7.12]{Fargues}.
For a motive over $F$, apply that realization over $\QQ$ to its additive
restriction of scalars: its realizations are induced, and
$L_{\QQ}(s,\operatorname{Res}_{F/\QQ}M)=L_F(s,M)$ by local induction.
Exterior and symmetric powers, tensor products, duals, Tate twists,
Artin twists, and direct summands preserve potential abelianity.
For CM abelian varieties this yields boundaries $j/2+7/8$ for
$H^j(A)$, $m/2+7/8$ for $\operatorname{Sym}^mH^1(A)$, and $15/8$
for $H^1(A)\otimes H^1(B)$.
\end{remark}

\begin{remark}[Dihedral decompositions]\label{subsec:dihedral}\label{prop:dihedral}
Let $K/F$ be quadratic with automorphism $c$ and character
$\epsilon=\epsilon_{K/F}$. For a unitary Hecke character $\theta$ put
$r=\operatorname{Ind}_{W_K}^{W_F}\theta$ and $\delta=\det r$.
Index-two induction gives
\[
 \operatorname{Ad}r\simeq\epsilon\oplus
        \operatorname{Ind}(\theta/\theta^c),\qquad
 \operatorname{Ind}\theta_1\otimes\operatorname{Ind}\theta_2
 \simeq\operatorname{Ind}(\theta_1\theta_2)
          \oplus\operatorname{Ind}(\theta_1\theta_2^c),
\]
and
\begin{equation}\label{eq:dihedral-symmetric}
 \operatorname{Sym}^m r\simeq
 \bigoplus_{0\le j<m/2}
       \operatorname{Ind}_{W_K}^{W_F}
             \bigl(\theta^{m-j}(\theta^c)^j\bigr)
 \oplus
 \begin{cases}
   0,&m\text{ odd},\\
   (\delta\epsilon)^{m/2},&m\text{ even}.
 \end{cases}
\end{equation}
The tensor identity is \cite[Lemma~2.3]{Langlands}. To check the middle
line in the symmetric power, choose $\tau\notin W_K$ and a basis with
$r(\tau)=\left(\begin{smallmatrix}0&a\\1&0\end{smallmatrix}\right)$,
$a=\theta(\tau^2)$. The monomials of indices $j$ and $m-j$ form the
induced pairs; for $m=2k$ the middle monomial has $\tau$-value $a^k$,
equal to $(\delta\epsilon)^k(\tau)$. The adjoint decomposition follows
from the diagonal trace-zero line and the two off-diagonal lines.

Local induction converts all these decompositions into Hecke
factorizations, so Theorem~\ref{thm:main} applies to every factor.
For an induced CM elliptic parameter with $\delta=1$ and
$\theta^c=\theta^{-1}$, for example,
$L_F(s,\operatorname{Sym}^4r)=
L_K(s,\theta^4)L_K(s,\theta^2)\zeta_F(s)$.
The two Hecke characters retain nonzero angular weights, and the
remaining factor gives a simple pole at $s=1$; in the arithmetic
variable it lies at $s=3$.
\end{remark}

\section*{AI provenance}
The manuscript text and mathematical arguments were developed predominantly
with OpenAI language models. The analytic core in
Sections~\ref{sec:transfer}--\ref{sec:closure} builds on the AI-generated
preprints \cite{Q,P}.

\clearpage
\appendix
\section{Transfer of the source arguments}\label{app:transfer-map}

The table separates the arithmetic or analytic input that changes with
the field from the norm calculation used in the source induction.
All ideals in the replacement column are exterior ideals. The set $S$
contains the fixed local data and generators of the class group;
$R=\mathcal O_{K,S}$ and $\Gamma=(R^\times)^6$.

\begingroup
\small
\setlength{\tabcolsep}{4pt}
\renewcommand{\arraystretch}{1.15}
\begin{longtable}{@{}>{\raggedright\arraybackslash}p{.17\textwidth}>{\raggedright\arraybackslash}p{.21\textwidth}>{\raggedright\arraybackslash}p{.40\textwidth}>{\raggedright\arraybackslash}p{.16\textwidth}@{}}
\caption{Source inputs, field replacements, and preserved calculations.}
\label{tab:transfer-map}\\
\toprule
Source step & Field-sensitive input & Replacement and reason the exponent is preserved & Location here\\
\midrule
\endfirsthead
\multicolumn{4}{l}{\textit{Table \thetable\ continued}}\\
\toprule
Source step & Field-sensitive input & Replacement and reason the exponent is preserved & Location here\\
\midrule
\endhead
\midrule
\multicolumn{4}{r}{\textit{Continued on the next page}}\\
\endfoot
\bottomrule
\endlastfoot
Generators and arithmetic indexing; $\Qsrc$; $\Psrc$, \S3
& Primary generators in the Eisenstein ring and its finite unit group ($\Qsrc$); cyclotomic fields containing $\mu_{12}$ ($\Psrc$).
& Localize to make $R$ principal. Allowed generators are unique modulo
$\Gamma$, and residual unit classes form a finite set. Each ideal is
counted with its displayed multiplicity; the finite sector count changes
only the constant.
& Proposition~\ref{prop:arithmetic-calibration}; start of
Section~\ref{sec:transfer}.\\[3pt]
Gauss calibration and reciprocity; $\Psrc$, Lemma~3.4
& Cubic reciprocity, Gauss phases, and the primary-generator convention.
& Hilbert reciprocity on an isotropic complement supplies allowed
reciprocity. Norm pullback of the cubic Jacobi character and the Weil
index supply $\Lambda_K,G$. The CRT correlations use the same local
powers of $q_{\mathfrak p}$.
& Proposition~\ref{prop:arithmetic-calibration};
\eqref{eq:can-local-G}.\\[3pt]
Orbital Poisson and complete correlations; $\Qsrc$, \S\S13, 17--18
& Additive lattice shape, finite units, and common support of the moduli.
& Integrate on the compact quotient $T_0/\Gamma$ and balance generators.
The multiplier is the row scale divided by the lattice index, and the
dual length is the conductor norm divided by that scale. Higher
dimension changes derivative orders; norm-sheet logarithms enter the
finite propagation bound.
& Lemmas~\ref{lem:weighted-orbital-poisson}--\ref{lem:balanced-lattice-decay};
Lemma~\ref{lem:polylog-propagation}.\\[3pt]
Theta and marked reflection; $\Qsrc$, \S14; $\Psrc$, \S\S5--6
& The finite-background theta expansion and its archimedean test.
& A character-weighted norm-sheet period gives the fixed-type Gamma
multiplier. Local mark digits separate the row, mark, and dual
coefficients. Their remaining interactions are the quadratic and cubic
symbols in \eqref{eq:marked-only-couplings}; these preserve the
source conductor powers and collision masks.
& Proposition~\ref{prop:fixed-type-theta};
Theorem~\ref{thm:marked-expansion}.\\[3pt]
Reflected norm and short inverse terminal; $\Qsrc$, \S\S14, 17
& Quadratic and cubic sieve families, powerful rows, and local dual data.
& Fixed ray-class splitting identifies the sieve families. Ideal
factorization gives the same powerful-row count. Finite-$S$ frequencies
are included in the common profile. Combining its squared decay with
the explicit local weights gives \eqref{eq:can-Eref} and the
three terminal bounds \eqref{eq:can-terminal-ledger}.
& Lemmas~\ref{lem:exterior-cubic-sieve}--\ref{lem:can-marked-energy};
Proposition~\ref{prop:can-canonical}.\\[3pt]
Two inverse transforms and initialization; $\Qsrc$, \S17
& Complete Gauss correlations, primary-generator fibres, and finite
backgrounds in the two transforms.
& The local tables and CRT are unchanged after replacing prime sizes by
norms. Allowed-generator fibres have divisor-bounded multiplicity and
finitely many unit sectors. The fixed background stays on the columns.
The child margins and row decrease are the affine norm identities
\eqref{eq:can-child-identities}--\eqref{eq:can-child-margins}; the
prefactor identity \eqref{eq:can-energy-identity} counts each outer
measure once.
& Proposition~\ref{prop:can-canonical}; proof of
Theorem~\ref{thm:can-inverse}.\\[3pt]
Ordinary reflection and the prime premise; $\Qsrc$, \S18.2
& Rank-one functional equation, fixed gamma factors, and finite
characters in the prime sums.
& Hecke reflection has dual norm length $Q/X$ for the actual conductor
$Q$. The paired coefficient cancels the infinite-order background at
primes. Its small annular terms are summable, including the fixed
logarithmic factors. Thus the capacity remains
$n_1+n_2+6\kappa z\le m+q$.
& Lemma~\ref{lem:gamma-reflection};
\eqref{eq:can-pairing};
Section~\ref{subsec:ordinary-reflection}.\\[3pt]
Fourth-moment recursion and exceptional rows; $\Qsrc$, \S\S18.2--18.8
& Finite two-column correlations, the prime pool, exceptional rows, and
common-product cancellation.
& Ideal valuations give the local inequality for each pair of
multiplicities and the same count of sixth-power classes. The prime
ideal theorem supplies the pool with its single density gain. A common
masked Mellin residue cancels on equal-product rectangles. The resulting
child width is \eqref{eq:can-plain-child}; all remaining losses are
positive parts of the displayed log-norm expressions.
& Lemma~\ref{lem:can-exceptional};
Section~\ref{subsec:fourth-induction};
Proposition~\ref{prop:can-uniformity}.\\[3pt]
Physical estimates and Euler compensation; $\Qsrc$, \S\S5--16;
$\Psrc$, \S7
& Additive Gram estimate, theta terminal, and local Euler algebra.
& The Gram proof uses the common two-column orbital kernel. The complete
marked reflection supplies its row energy. Local compensation is a
rational identity in unit-modulus character values, so it gives the same
physical norm monomial and row exponents.
& Proposition~\ref{prop:additive-gram};
Section~\ref{sec:physical}.\\[3pt]
Exponent comparison and closure; $\Qsrc$, \S\S19--20 and Proposition~2.1
& Numerical lengths and the common zero boundary.
& The comparison is field-independent algebra. The linear envelope and
\eqref{row:square-certificate} give the fixed positive reserve; the
Mellin transform uses the same normalized function on both sides.
& Section~\ref{sec:rows};
Lemma~\ref{lem:mellin-closure}.\\
\end{longtable}
\endgroup

The rows on the two inverse transforms and on the fourth-moment
recursion invoke the recursive arguments of \cite[\S\S17--18]{Q} by
locator rather than reproducing them. Checking the transfer amounts to
confirming, against those sections, that the inputs in the second
column are the only field- or type-dependent ones.

The two recursion rows use three concrete inputs: the complete local
correlation, one common smooth kernel, and the explicitly counted ideal
fibres. Once those have been supplied, the child lengths, capacities,
and normalization factors are identities among ideal norms. The
finite-depth estimate in Lemma~\ref{lem:polylog-propagation} accounts for
the additional logarithms of the unit-quotient kernels.

\begingroup
\small
\begin{thebibliography}{99}

\bibitem{BGL}
V.~Blomer, L.~Goldmakher, and B.~Louvel,
\emph{$L$-functions with $n$-th-order twists},
International Mathematics Research Notices \textbf{2014} (2014), no.~7,
1925--1955.
\href{https://doi.org/10.1093/imrn/rns257}{doi:10.1093/imrn/rns257}.
Preprint consulted for theorem numbering:
\href{https://arxiv.org/abs/1112.1650v1}{arXiv:1112.1650v1}, 7 December 2011.

\bibitem{Brauer}
R.~Brauer,
\emph{On Artin's $L$-series with general group characters},
Annals of Mathematics (2) \textbf{48} (1947), no.~2, 502--514.
\href{https://doi.org/10.2307/1969183}{doi:10.2307/1969183}.

\bibitem{CT}
D.~Cureton and J.~Thorner,
\emph{A quasi-Riemann hypothesis for Dedekind zeta functions},
Hexagon preprint hexagon:2610.00200v1, 12 October 2026.
\url{https://hexagonmath.org/2610.00200v1}.
\bibitem{DS}
P.~Deligne and J.-P.~Serre,
\emph{Formes modulaires de poids $1$},
Annales scientifiques de l'\'Ecole Normale Sup\'erieure (4)
\textbf{7} (1974), no.~4, 507--530.
\href{https://doi.org/10.24033/asens.1277}{doi:10.24033/asens.1277}.

\bibitem{DR}
A.~Dunn and M.~Radziwi{\l}{\l},
\emph{Bias in cubic Gauss sums: Patterson's conjecture},
Annals of Mathematics (2) \textbf{200} (2024), no.~3, 967--1057.
\href{https://doi.org/10.4007/annals.2024.200.3.3}{doi:10.4007/annals.2024.200.3.3}.
Preprint consulted for section and theorem numbering:
\href{https://arxiv.org/abs/2109.07463v3}{arXiv:2109.07463v3}, 14 May 2024.

\bibitem{Fargues}
L.~Fargues,
\emph{Motives and automorphic forms: the (potentially) abelian case},
expository notes, undated, 42-page author-hosted version.
\url{https://webusers.imj-prg.fr/~laurent.fargues/Motifs_abeliens.pdf}.
The locators used here are in Sections~2 and~7; consulted 12 October 2026.

\bibitem{GL}
L.~Goldmakher and B.~Louvel,
\emph{A quadratic large sieve inequality over number fields},
Mathematical Proceedings of the Cambridge Philosophical Society
\textbf{154} (2013), no.~2, 193--212.
\href{https://doi.org/10.1017/S0305004112000370}{doi:10.1017/S0305004112000370}.
Preprint consulted for definition and theorem numbering:
\href{https://arxiv.org/abs/1112.1642v2}{arXiv:1112.1642v2}, 31 May 2012.

\bibitem{HBQ}
D.~R.~Heath-Brown,
\emph{A mean value estimate for real character sums},
Acta Arithmetica \textbf{72} (1995), no.~3, 235--275.
\href{https://doi.org/10.4064/aa-72-3-235-275}{doi:10.4064/aa-72-3-235-275}.

\bibitem{HBC}
D.~R.~Heath-Brown,
\emph{Kummer's conjecture for cubic Gauss sums},
Israel Journal of Mathematics \textbf{120} (2000), 97--124.
\href{https://doi.org/10.1007/s11856-000-1273-y}{doi:10.1007/s11856-000-1273-y}.

\bibitem{HK}
A.~Hoshi and K.~Kanai,
\emph{Davenport and Hasse's theorems and lifts of multiplication
matrices of Gaussian periods},
Finite Fields and Their Applications \textbf{84} (2022), article 102101.
\href{https://doi.org/10.1016/j.ffa.2022.102101}{doi:10.1016/j.ffa.2022.102101}.
Extended preprint consulted for Theorem~1.1:
\href{https://arxiv.org/abs/2105.14872v2}{arXiv:2105.14872v2}, 4 August 2022.

\bibitem{Johansson}
C.~Johansson, with an appendix by F.~Fit\'e,
\emph{On the Sato--Tate conjecture for non-generic abelian surfaces},
Transactions of the American Mathematical Society
\textbf{369} (2017), no.~9, 6303--6325.
\href{https://doi.org/10.1090/tran/6847}{doi:10.1090/tran/6847}.
Numbering follows \href{https://arxiv.org/abs/1307.6478v5}{arXiv:1307.6478v5},
8 October 2015, Section~3.

\bibitem{KP}
D.~A.~Kazhdan and S.~J.~Patterson,
\emph{Metaplectic forms},
Publications Math\'ematiques de l'IH\'ES \textbf{59} (1984), 35--142.
\href{https://doi.org/10.1007/BF02698770}{doi:10.1007/BF02698770}.
Read with the corrections \cite{KPcorr}.

\bibitem{KPcorr}
D.~A.~Kazhdan and S.~J.~Patterson,
\emph{Corrections to: ``Metaplectic forms''},
Publications Math\'ematiques de l'IH\'ES \textbf{62} (1985), 203.
\href{https://doi.org/10.1007/BF02698809}{doi:10.1007/BF02698809}.

\bibitem{Langlands}
R.~P.~Langlands,
\emph{On the functional equation of the Artin $L$-functions},
unpublished preprint, 1970.
IAS electronic edition:
\url{https://publications.ias.edu/sites/default/files/on-the-functional-equation-of-the-artin-l-functions_rpl_5.pdf},
240 PDF pages.  Locators refer to its printed numbering:
Lemma~2.2 on pp.~8--9 and Chapter~20 on pp.~211--213.

\bibitem{Milne}
J.~S.~Milne,
\emph{Class Field Theory}, version 4.03, 6 August 2020.
\url{https://www.jmilne.org/math/CourseNotes/CFT.pdf}.

\bibitem{Q}
OpenAI,
\emph{The Quasi-Riemann Hypothesis: A Zero-Free Half-Plane
$\operatorname{Re}(s)>7/8$},
preprint, 30 September 2026, 199-page edition.
Release identifier:
\nolinkurl{OAI:The-Quasi-Riemann-Hypothesis-September-30-2026}.
Dated release record:
\url{https://github.com/openai/math/blob/main/preprints/The-Quasi-Riemann-Hypothesis-September-30-2026/README.md}.
Paper:
\url{https://github.com/openai/math/blob/main/preprints/The-Quasi-Riemann-Hypothesis-September-30-2026/paper.pdf}.
All locators for $\Qsrc$ refer to this dated edition.

\bibitem{P}
OpenAI,
\emph{Primitive roots for every admissible integer base},
preprint, 4 October 2026, 92-page edition.
Release identifier:
\nolinkurl{OAI:Primitive-roots-for-every-admissible-integer-base-October-4-2026}.
Dated release record:
\url{https://github.com/openai/math/blob/main/preprints/Primitive-roots-for-every-admissible-integer-base-October-4-2026/README.md}.
Paper:
\url{https://github.com/openai/math/blob/main/preprints/Primitive-roots-for-every-admissible-integer-base-October-4-2026/primitive-roots-all-integer-bases.pdf}.
All locators for $\Psrc$ refer to this dated edition.

\bibitem{Patterson}
S.~J.~Patterson,
\emph{A cubic analogue of the theta series},
Journal f\"ur die reine und angewandte Mathematik
\textbf{296} (1977), 125--161.
\href{https://doi.org/10.1515/crll.1977.296.125}{doi:10.1515/crll.1977.296.125}.

\bibitem{Tate}
J.~T.~Tate,
\emph{Fourier analysis in number fields and Hecke's zeta-functions},
in J.~W.~S.~Cassels and A.~Fr\"ohlich (eds.),
\emph{Algebraic Number Theory}, Academic Press, London, 1967, 305--347.

\bibitem{Weil}
A.~Weil,
\emph{Sur certains groupes d'op\'erateurs unitaires},
Acta Mathematica \textbf{111} (1964), 143--211.
\href{https://doi.org/10.1007/BF02391012}{doi:10.1007/BF02391012}.

\bibitem{DKN}
S.~Das, H.~Kadiri, and N.~Ng,
\emph{An effective version of Chebotarev's density theorem},
preprint, arXiv:2508.09480v1, 13 August 2025.
\url{https://arxiv.org/abs/2508.09480v1}.
The cyclic reduction and explicit-formula locators refer to Section~3
of this version.

\end{thebibliography}
\endgroup

\end{document}
